% WZ-LOCKED-CLASS-BEGIN
\begin{filecontents*}[overwrite]{elegantbook-original-adapter.cls}
\NeedsTeXFormat{LaTeX2e}
\ProvidesClass{elegantbook-original-adapter}[2026/09/23 v3.0 fixed mathematical book environments]
\LoadClass[10pt,letterpaper,oneside]{book}
\RequirePackage[UTF8,scheme=plain,fontset=fandol]{ctex}
\RequirePackage{fontspec}
\setmainfont{cmunrm.otf}[BoldFont=cmunbx.otf,ItalicFont=cmunti.otf,BoldItalicFont=cmunbi.otf]
\setCJKfamilyfont{sourceheader}[ItalicFont=FandolKai-Regular.otf]{FandolKai-Regular.otf}
\RequirePackage{amsmath,amssymb,mathrsfs}
\RequirePackage{amsthm,etoolbox}
\RequirePackage{xcolor,geometry,setspace,fancyhdr,enumitem,needspace,xurl}
\definecolor{structurecolor}{HTML}{880E4F}
\definecolor{main}{HTML}{9C27B0}
\geometry{paperwidth=612bp,paperheight=792bp,left=20mm,right=20mm,top=95.04bp,bottom=20mm,headheight=12pt,headsep=25pt,footskip=30pt}
\setstretch{1.6}
\setlength{\parindent}{2em}
\setlength{\parskip}{0pt}
\setlist{nosep,leftmargin=2em}
\AtBeginDocument{\setlength{\abovedisplayskip}{4pt plus 1pt minus 1pt}\setlength{\belowdisplayskip}{4pt plus 1pt minus 1pt}\setlength{\abovedisplayshortskip}{2pt}\setlength{\belowdisplayshortskip}{3pt}}
\fancyhf{}
\fancyhead[L]{{\itshape\CJKfamily{sourceheader}\rightmark}}
\fancyhead[R]{{\itshape\CJKfamily{sourceheader}\leftmark}}
\fancyfoot[C]{\thepage}
\renewcommand{\headrulewidth}{0.4pt}
\renewcommand{\headrule}{{\color{structurecolor}\hrule height\headrulewidth width\headwidth\vskip-\headrulewidth}}
\pagestyle{fancy}
\fancypagestyle{cover}{\fancyhf{}\fancyfoot[C]{\thepage}\renewcommand{\headrulewidth}{0pt}\renewcommand{\headrule}{}}
\RequirePackage{hyperref}
\hypersetup{unicode=true,colorlinks=true,linkcolor=structurecolor,urlcolor=structurecolor,citecolor=structurecolor,linktoc=section,bookmarksnumbered=true,bookmarksopen=true,pdfborder={0 0 0},pdfstartview=FitH}
\setcounter{tocdepth}{1}
\renewcommand*\l@section{\@dottedtocline{1}{1.5em}{2.4em}}
\renewcommand{\thesection}{\thechapter.\arabic{section}}
\newcommand{\originalchapter}[1]{\par\addvspace{14pt}\Needspace{5\baselineskip}\refstepcounter{chapter}\setcounter{section}{0}\addcontentsline{toc}{chapter}{\protect\numberline{\thechapter}#1}\markboth{\thechapter\quad #1}{}\begingroup\centering\color{structurecolor}\bfseries\fontsize{14.4}{20}\selectfont\thechapter\quad #1\par\endgroup\nobreak\vspace{8pt}}
\newcommand{\originalsection}[1]{\par\addvspace{9pt}\Needspace{4\baselineskip}\refstepcounter{section}\addcontentsline{toc}{section}{\protect\hspace*{1.5em}\protect\numberline{\protect\textcolor{structurecolor}{\thesection}}#1}\markright{\thesection\quad #1}\noindent{\fontsize{12}{17}\selectfont\color{structurecolor}\thesection\quad}{\bfseries\fontsize{12}{17}\selectfont #1}\par\nobreak\vspace{4pt}}

% Semantic environments are registered once here, not re-designed in manuscripts.
\newtheoremstyle{wzstatement}{6pt}{5pt}
  {\normalfont\kaishu}{0pt}{\normalfont\bfseries\color{structurecolor}}
  {}{.7em}{\thmname{#1}\thmnumber{ #2}\thmnote{\normalfont\ (#3)}}
\newtheoremstyle{wzdefinition}{6pt}{5pt}
  {\normalfont\songti}{0pt}{\normalfont\bfseries\color{structurecolor}}
  {}{.7em}{\thmname{#1}\thmnumber{ #2}\thmnote{\normalfont\ (#3)}}
\newtheoremstyle{wzproblem}{7pt}{5pt}
  {\normalfont\songti}{2em}{\normalfont\bfseries\color{main}}
  {}{.7em}{\thmname{#1}\thmnumber{ #2}}
\newtheoremstyle{wzremark}{4pt}{4pt}
  {\normalfont\songti}{0pt}{\normalfont\bfseries}
  {}{.7em}{\thmname{#1}}
\theoremstyle{wzstatement}
\newtheorem{theorem}{定理}[chapter]
\newtheorem{lemma}[theorem]{引理}
\newtheorem{proposition}[theorem]{命题}
\newtheorem{corollary}[theorem]{推论}
\theoremstyle{wzdefinition}
\newtheorem{definition}[theorem]{定义}
\theoremstyle{wzproblem}
\newtheorem{example}{例题}[chapter]
\newtheorem{exercise}{习题}[chapter]
\theoremstyle{wzremark}
\newtheorem*{remark}{注}
\renewcommand{\proofname}{证明}
% Retain the native amsthm QED stack; only adjust the fixed typography.
\expandafter\patchcmd\csname\string\proof\endcsname{\normalfont \topsep6\p@\@plus6\p@\relax}
  {\normalfont\songti \topsep5\p@\@plus1\p@\relax}
  {}{\ClassError{elegantbook-original-adapter}{Cannot patch proof spacing}{Check the amsthm version.}}
\expandafter\patchcmd\csname\string\proof\endcsname{\itshape}{\normalfont\bfseries}
  {}{\ClassError{elegantbook-original-adapter}{Cannot patch proof heading}{Check the amsthm version.}}
\expandafter\patchcmd\csname\string\proof\endcsname{\ignorespaces}
  {\setlength{\parindent}{2em}\ignorespaces}
  {}{\ClassError{elegantbook-original-adapter}{Cannot patch proof paragraphs}{Check the amsthm version.}}
\AtBeginEnvironment{proof}{\par\Needspace{3\baselineskip}}
\renewcommand{\qedsymbol}{\ensuremath{\square}}
\newenvironment{solution}{\begin{proof}[解答]}{\end{proof}}
\newcommand{\wzcheckchapter}{\ifnum\value{chapter}<1
  \ClassError{elegantbook-original-adapter}{Numbered mathematics before chapter 1}{Move it after the first originalchapter.}\fi}

\newcommand{\wzregisterguard}[1]{\AtBeginEnvironment{#1}{\wzcheckchapter\par\Needspace{3\baselineskip}}}
\forcsvlist{\wzregisterguard}{theorem,lemma,proposition,corollary,definition,example,exercise}
% Front matter and bibliography are not numbered mathematical chapters.
\newcommand{\originalpreface}{\par\addvspace{10pt}\Needspace{5\baselineskip}
  \noindent{\bfseries\fontsize{12}{17}\selectfont 前言}\par\nobreak\vspace{4pt}}
\newcommand{\originalcontents}{\par\addvspace{14pt}
  \begin{center}{\color{structurecolor}\bfseries\fontsize{14.4}{20}\selectfont 目录}\end{center}
  \vspace{-10pt}\@starttoc{toc}}
\renewcommand{\bibname}{参考文献}
\renewenvironment{thebibliography}[1]{
  \par\addvspace{12pt}\Needspace{3\baselineskip}\phantomsection
  \addcontentsline{toc}{chapter}{\bibname}
  \markboth{\bibname}{}
  \noindent{\color{structurecolor}\bfseries\fontsize{14.4}{20}\selectfont\bibname}\par\nobreak\vspace{6pt}
  \list{\@biblabel{\@arabic\c@enumiv}}
    {\settowidth\labelwidth{\@biblabel{#1}}\leftmargin\labelwidth
     \advance\leftmargin\labelsep\usecounter{enumiv}
     \let\p@enumiv\@empty\renewcommand\theenumiv{\@arabic\c@enumiv}
     \setlength{\itemsep}{3pt}\setlength{\parsep}{0pt}}
  \sloppy\clubpenalty4000\widowpenalty4000\sfcode`\.=1000\relax
}{\def\@noitemerr{\@latex@warning{Empty bibliography}}\endlist}

\numberwithin{equation}{chapter}
\renewcommand{\theequation}{\thechapter.\arabic{equation}}
\allowdisplaybreaks[2]
\emergencystretch=2em
\clubpenalty=6000
\widowpenalty=6000
\raggedbottom
\endinput
\end{filecontents*}
% WZ-LOCKED-CLASS-END
\documentclass{elegantbook-original-adapter}
% ===== 元数据内容槽 =====
\newcommand{\BookTitle}{凸分析与最优化}
\newcommand{\BookAuthor}{据 Dimitri P. Bertsekas 课程讲义整理}
\newcommand{\BookDate}{2026年10月}
\newcommand{\PDFTitle}{凸分析与最优化: MIT 6.253 中文讲义}
\newcommand{\PDFAuthor}{Dimitri P. Bertsekas; 中文整理}
\hypersetup{pdftitle={\PDFTitle},pdfauthor={\PDFAuthor}}
% 数学记号可以增加；禁止重定义版式、环境或证毕符号.
\newcommand{\dd}{\,\mathrm d}
\newcommand{\R}{\mathbb R}
\newcommand{\norm}[1]{\lVert #1\rVert}
\newcommand{\inner}[2]{\langle #1,#2\rangle}
\usepackage{tikz,longtable,booktabs}
\usetikzlibrary{arrows.meta,calc}
\DeclareMathOperator{\dom}{dom}
\DeclareMathOperator{\epi}{epi}
\DeclareMathOperator{\conv}{conv}
\DeclareMathOperator{\aff}{aff}
\DeclareMathOperator{\cone}{cone}
\DeclareMathOperator{\ri}{ri}
\DeclareMathOperator{\cl}{cl}
\DeclareMathOperator{\rec}{rec}
\DeclareMathOperator{\argmin}{argmin}
\DeclareMathOperator{\argmax}{argmax}
\DeclareMathOperator{\dist}{dist}
\DeclareMathOperator{\tr}{tr}
\newcommand{\T}{^{\mathsf T}}
\newcommand{\ind}[1]{\delta_{#1}}
\begin{document}
\frontmatter
\thispagestyle{cover}
\begin{center}
{\color{structurecolor}\bfseries\fontsize{20}{28}\selectfont 凸分析与最优化\par}
\vspace{8pt}
{\large MIT 6.253 中文重构讲义\par}
\vspace{10pt}
Dimitri P. Bertsekas 原课程讲义, Spring 2012\par
2026年10月中文整理\par
\end{center}
\originalpreface
本书以 MIT OpenCourseWare 的 6.253 \emph{Convex Analysis and Optimization} 2012 年春季课程完整讲义为主源, 按凸集与函数、最优解与分离、对偶、结构化模型和数值算法的依赖关系重新组织. 原资料是教学幻灯片; 本稿把其数学陈述转写为可连续阅读的简体中文正文, 补充定义之间的联系、证明中省略的论证及例题计算. 补充论证由整理者负责, 不能理解为原作者逐字表述.

预备知识是有限维线性代数、多元微积分及实分析中的序列、闭集与紧集. 所有向量默认是列向量, $x\T y$ 表示欧氏内积, $\norm{x}=\sqrt{x\T x}$. 最优值允许为无穷; “最优解存在”与“最优值有限”分别讨论. 为避免遗漏条件, 后文使用闭、真、凸函数时均说明其含义.

原作者署名为 Dimitri P. Bertsekas, 课程提供者为 Massachusetts Institute of Technology / MIT OpenCourseWare. 原始公开课程材料的适用许可见 \href{https://ocw.mit.edu/pages/privacy-and-terms-of-use/}{MIT OCW 使用条款}及 \href{https://creativecommons.org/licenses/by-nc-sa/4.0/}{CC BY-NC-SA 4.0}. 本稿为翻译、重排与补充证明的改编, 按同一非商业、署名、相同方式共享许可提供; MIT 与原作者未为本稿背书. 原PDF第1页说明图形由 Athena Scientific 授权使用. 本稿不复制这些图形, 不转载外部教材全文; 所绘示意图是按数学定义重新绘制. 外部书籍、论文及课程链接不因此自动取得 OCW 许可.

各章首给出对应原讲次与原PDF页范围, 便于对照; 同一原讲的复习内容可在数章出现. 数学正文不沿用幻灯片式省略, 也不把算法列表当作收敛证明.
\originalcontents
\clearpage
% 按当前任务填写正文. 模板不规定篇幅、证明深度或章节内容.
% 需要前言/目录时使用 \originalpreface / \originalcontents.
\mainmatter
\originalchapter{凸集与凸函数}\label{chap:1}
本章对应原讲义第1讲的优化概念及第2--3讲, 原PDF第2--38页. 先理解凸性提供了什么, 再把这一几何性质转写为分析和计算工具.

\originalsection{优化问题与凸性}
最优化问题的基本形式是求 $\inf_{x\in C}f(x)$. 集合 $C$ 是可行集, $f$ 是目标函数. 记最优值为 $f^*$; 若某个 $x^*\in C$ 满足 $f(x^*)=f^*$, 则称其为最优解. 写 $\inf$ 而非 $\min$ 是为了不预设最优解存在. 例如 $f(x)=e^x$ 在 $\R$ 上的下确界为0, 却没有取到0的点.

\begin{definition}
集合 $C\subseteq\R^n$ 称为凸集, 若任意 $x,y\in C$ 及 $t\in[0,1]$ 都满足 $(1-t)x+ty\in C$. 在凸集 $C$ 上, 函数 $f:C\to\R$ 称为凸函数, 若
\[
f((1-t)x+ty)\le(1-t)f(x)+tf(y).
\]
若对 $x\ne y$ 及 $0<t<1$ 不等号严格, 则称 $f$ 严格凸. 凸规划指 $C$ 与 $f$ 都凸的最优化问题.
\end{definition}
凸集要求连接两个可行点的整条线段仍然可行; 凸函数要求线段上的函数图像不高于两端图像的弦. 这不是“图像看起来像碗”的模糊描述, 因为可行集可以位于低维平面上, 目标也可以不可微.

\begin{proposition}\label{prop:localglobal}
凸函数在凸集上的每个局部极小点都是全局极小点. 最优解集合是凸集; 严格凸函数至多有一个最优解.
\end{proposition}
\begin{proof}
设 $x^*$ 是局部极小点, 若存在 $y\in C$ 使 $f(y)<f(x^*)$, 则对任意 $0<t<1$,
$f((1-t)x^*+ty)\le(1-t)f(x^*)+tf(y)<f(x^*)$. 取足够小的 $t$ 就与局部极小性矛盾. 若 $x,y$ 都最优, 凸性给出线段上函数值不大于 $f^*$, 而最优值的定义给出反向不等式. 若严格凸且 $x\ne y$, 线段内部的值严格小于 $f^*$, 产生矛盾.
\end{proof}
上述结论不保证最优解存在, 也不保证任何数值算法都能快速找到它. 原课程强调对偶的另一作用: 即使原问题含有整数限制或其他非凸结构, 由拉格朗日乘子得到的对偶目标仍是凹函数, 因而对偶最大化具有凸优化结构. 我们将在对偶章节证明这一点及其下界意义.

\begin{example}
比较 $f(x)=x^2$ 在 $C_1=\R$ 与 $C_2=\{-2,1\}$ 上的局部、全局最优点.
\end{example}
\begin{solution}
在 $C_1$ 上, $f$ 严格凸, 唯一最优点为0. 在离散可行集 $C_2$ 中, 两个点都在相对拓扑意义下局部最优, 因为可取一个小邻域不包含另一个可行点. 但 $f(-2)=4>1=f(1)$, 所以 $-2$ 不是全局最优. 关键障碍在于两点间线段不在 $C_2$ 中, 命题\ref{prop:localglobal}的证明无法使用.
\end{solution}

\originalsection{保持凸性的运算}
半空间 $\{x:a\T x\le b\}$ 和超平面 $\{x:a\T x=b\}$ 都是凸集. 有限个半空间的交称为多面体, 它闭且凸, 但不一定有界. 锥是对正数乘法封闭的集合; 锥本身不一定凸或闭. 本书涉及生成锥时将明确包括原点.

\begin{proposition}
任意族凸集的交为凸集. 凸集的线性像、仿射像、逆像、数乘、向量和、闭包与内部均为凸集.
\end{proposition}
\begin{proof}
交集的两点属于每一个原集合, 因而线段逐一属于它们. 对线性映射 $A$, 有 $A((1-t)x+ty)=(1-t)Ax+tAy$; 这一恒等式同时证明线性像与逆像的结论, 平移不改变线段结构. 若 $x_i\in C_i$, 则 $(1-t)(x_1+x_2)+t(y_1+y_2)$ 是两个可行线段点之和. 闭包的结论通过用 $C$ 中序列逼近端点, 然后取极限得到. 若端点属于内部, 它们各有一小球包含于 $C$; 对 $0<t<1$, 将一端小球与另一端固定点作凸组合, 得到以相应线段点为中心的小球. 端点本身已经是内点.
\end{proof}
内部可以为空, 例如 $\R^2$ 中的一条直线; 后面用相对内部处理这一情形. 线性像保持凸性却未必保持闭性, 两种性质不能混为一谈.

\begin{proposition}\label{prop:convexoperations}
设各 $f_i:\R^n\to(-\infty,+\infty]$ 为凸函数. 下列运算保持凸性: 正系数有限和 $\sum_i\alpha_i f_i$; 仿射复合 $f(Ax+b)$; 逐点上确界 $\sup_i f_i(x)$. 若参与运算的函数都闭, 上述得到的函数也闭.
\end{proposition}
\begin{proof}
凸性不等式可以乘以正数后相加, 也可以在仿射映射下代入. 对上确界, 每个 $i$ 都有
$f_i((1-t)x+ty)\le(1-t)\sup_jf_j(x)+t\sup_jf_j(y)$, 再对 $i$ 取上确界. 闭性将由下一节的下半连续性等价刻画得到. 仿射映射连续, 因而保持下半连续性; 上确界的上图是各上图的交. 对有限和, 在任一极限点附近, 每个下半连续函数都具有一个有限下界, 所以可以对下极限使用有限和不等式, 不会出现 $+\infty-\infty$.
\end{proof}
常用的凸函数包括仿射函数、范数及半正定二次型. 范数的凸性由三角不等式与齐次性给出. 二次函数 $f(x)=\tfrac12 x\T Qx+b\T x+c$ 的凸性取决于对称矩阵 $Q$ 是否半正定; 若原矩阵不对称, 先以 $(Q+Q\T)/2$ 代替它, 因为反对称部分对二次型没有贡献.

\begin{example}
证明 $f(x)=\max\{a_1\T x+b_1,\ldots,a_m\T x+b_m\}$ 凸, 并说明为何凸函数的逐点最小值未必凸.
\end{example}
\begin{solution}
每个仿射函数凸, 所以命题\ref{prop:convexoperations}适用. 反例为 $g(x)=\min\{x,-x\}=-|x|$. 取 $x=-1,y=1,t=1/2$, 有 $g(0)=0>-1=(g(-1)+g(1))/2$. 上确界与下确界在凸性运算中并不对称.
\end{solution}

\originalsection{上图、闭性与扩展实值}
约束可以放入目标函数中. 对集合 $C$, 定义指示函数 $\ind{C}(x)=0$ 当 $x\in C$, 在其他点取 $+\infty$. 则在 $C$ 上最小化 $f$ 等价于在全空间最小化 $f+\ind{C}$. 这里的 $+\infty$ 是禁止不可行点的记号, 不参与普通实数算术.

\begin{definition}
对于扩展实值函数 $f:\R^n\to[-\infty,+\infty]$, 定义有效域 $\dom f=\{x:f(x)<+\infty\}$ 及上图
$\epi f=\{(x,w)\in\R^{n+1}:f(x)\le w\}$. 若上图凸, 称 $f$ 凸; 若上图闭, 称 $f$ 闭. 若 $\dom f\ne\varnothing$ 且 $f$ 从不取 $-\infty$, 称 $f$ 为真函数. 对序列 $x_k\to x$, 若总有 $f(x)\le\liminf_kf(x_k)$, 则称 $f$ 在 $x$ 下半连续.
\end{definition}
对有限值函数, 上图凸性与此前的函数凸性定义等价. 一方面, 取上图中两点并使用函数凸性; 另一方面, 对图像上两点 $(x,f(x))$ 与 $(y,f(y))$ 使用上图凸性. 有效域也因此凸. 函数的下水平集 $\{x:f(x)\le\gamma\}$ 和严格下水平集 $\{x:f(x)<\gamma\}$ 都凸, 但反过来, 仅有水平集凸只说明拟凸, 不足以说明函数凸.

\begin{center}
\begin{tikzpicture}[scale=.85]
\draw[->] (-.2,0)--(4.3,0) node[right]{$x$};
\draw[->] (0,-.2)--(0,3.8) node[above]{$w$};
\draw[thick,domain=.2:3.6,samples=70] plot (\x,{.35*(\x-1.7)^2+.6});
\draw[dashed] (.6,1.024)--(3.2,1.3875);
\node at (2.1,2.8){$\epi f$};
\node at (3.4,.6){$f$};
\fill (.6,1.024) circle(1.5pt); \fill (3.2,1.3875) circle(1.5pt);
\end{tikzpicture}
\end{center}
图中虚线为两端图像点之间的线段, 凸性要求它处在上图中. 图只辅助定义, 不代替其量词条件.

\begin{theorem}\label{thm:lsc}
对定义在整个 $\R^n$ 上的扩展实值函数, 下列三条等价: 每个实数 $\gamma$ 的下水平集都闭; 函数处处下半连续; 上图闭.
\end{theorem}
\begin{proof}
若函数下半连续, 取 $(x_k,w_k)\in\epi f$ 且 $(x_k,w_k)\to(x,w)$, 则 $f(x)\le\liminf f(x_k)\le w$, 所以上图闭. 若上图闭, 由 $x_k$ 属于 $\gamma$ 水平集可知 $(x_k,\gamma)\in\epi f$, 极限 $(x,\gamma)$ 仍在上图, 因而水平集闭.

最后假设所有水平集闭. 若在某个 $x$ 有 $f(x)>\liminf_kf(x_k)$, 可在二者之间取实数 $\gamma$. 存在子序列满足 $f(x_{k_j})\le\gamma$, 又有 $x_{k_j}\to x$. 水平集闭便给出 $f(x)\le\gamma$, 矛盾. 这也适用于无穷值, 因为两者严格不等时仍能选取有限的中间实数.
\end{proof}
“定义在整个空间”不能省略. 在 $(0,1)$ 上恒为0的函数在其定义域内下半连续, 但上图 $(0,1)\times[0,\infty)$ 并不闭. 若把端点函数值定义为 $+\infty$, 上图不变, 而端点处的下半连续性失败. 因此约束进入目标后, 应在全空间上检查闭性.

\begin{example}
给出一个闭真凸函数, 其有效域不是闭集.
\end{example}
\begin{solution}
取 $f(x)=1/x$ 当 $x>0$, 在 $x\le0$ 取 $+\infty$. 在正半轴上 $f''(x)=2/x^3>0$, 上图凸. 对有限 $\gamma>0$, 下水平集为 $[1/\gamma,\infty)$; 对 $\gamma\le0$ 为空集. 它们都闭, 所以函数闭, 但 $\dom f=(0,\infty)$ 不闭. 此例说明闭性由上图决定, 不能简单替换为有效域闭.
\end{solution}
若有效域闭且函数在有效域内下半连续, 则把域外取 $+\infty$ 后的函数确实闭. 证明时, 上图中的收敛序列的横坐标极限仍在有效域, 然后使用相对下半连续性即可.

原讲义也提到非真闭凸函数的特殊性. 若它不恒为 $+\infty$ 而在某点取 $-\infty$, 则它在每一个有效域点都取 $-\infty$. 设 $f(x_0)=-\infty$, $y\in\dom f$. 凸性使 $(1-t)y+tx_0$ 在任意 $t>0$ 时都具有任意低的上图高度; 令 $t\downarrow0$, 上图闭性使 $(y,r)$ 对任意实数 $r$ 都属于上图, 所以 $f(y)=-\infty$. 因而后续优化理论主要使用真函数, 避免这类退化情形.

\originalsection{可微凸函数与最优性}
\begin{theorem}\label{thm:firstorder}
设 $C$ 凸, $f$ 在包含 $C$ 的开集上可微. 函数 $f$ 在 $C$ 上凸, 当且仅当对任意 $x,z\in C$,
\begin{equation}\label{eq:tangent}
f(z)\ge f(x)+\nabla f(x)\T(z-x).
\end{equation}
若对不同的 $x,z$ 不等号总严格, 则 $f$ 严格凸.
\end{theorem}
\begin{proof}
设 $f$ 凸. 对 $0<t\le1$, 凸性给出
$[f(x+t(z-x))-f(x)]/t\le f(z)-f(x)$. 令 $t\downarrow0$, 左侧趋于 $\nabla f(x)\T(z-x)$, 得到\eqref{eq:tangent}.

反过来, 设\eqref{eq:tangent}成立. 对 $w=(1-t)x+tz$, 分别在基点 $w$ 使用它, 得
$f(x)\ge f(w)+\nabla f(w)\T(x-w)$ 与 $f(z)\ge f(w)+\nabla f(w)\T(z-w)$. 乘以 $1-t$ 与 $t$ 后相加, 梯度项抵消, 得到凸性. 当 $0<t<1$ 且 $x\ne z$ 时, $x,w,z$ 两两不同; 严格版本给出严格凸性.
\end{proof}
\eqref{eq:tangent}把局部一阶信息变成了全局下界. 这就是凸优化中驻点容易判定的原因. 但有约束时, 梯度不必为零: 在边界处, 所有可行移动方向都可能让目标增加.

\begin{theorem}\label{thm:vi}
在定理\ref{thm:firstorder}的条件下再设 $f$ 凸. 点 $x^*\in C$ 最优, 当且仅当
\[
\nabla f(x^*)\T(x-x^*)\ge0\quad(x\in C).
\]
\end{theorem}
\begin{proof}
充分性由切平面下界直接得到. 对必要性, 固定 $x\in C$, 线段 $x^*+t(x-x^*)$ 在 $0\le t\le1$ 时可行. 最优性使差商 $[f(x^*+t(x-x^*))-f(x^*)]/t\ge0$. 令 $t\downarrow0$ 即得结论.
\end{proof}
无约束 $C=\R^n$ 时, 可以取梯度方向与反方向, 条件等价于 $\nabla f(x^*)=0$. 若 $C$ 是线性等式约束, 条件表示梯度垂直于所有可行方向. 一般凸集的情况将在法锥与次梯度章节统一处理.

\begin{theorem}\label{thm:hessian}
设 $f$ 二阶连续可微. 若在凸集 $C$ 上 $\nabla^2f(x)\succeq0$, 则 $f$ 在 $C$ 上凸; 若处处正定, 则严格凸. 若 $C$ 还是开集, 则凸性反过来蕴含 Hessian 半正定.
\end{theorem}
\begin{proof}
对 $x,z\in C$ 令 $h(t)=f(x+t(z-x))$. 它满足 $h''(t)=(z-x)\T\nabla^2f(x+t(z-x))(z-x)\ge0$. 因而 $h'$ 单调, 积分得 $h(1)\ge h(0)+h'(0)$; 由定理\ref{thm:firstorder}即得凸性. 当 $x\ne z$ 且 Hessian 正定时, $h''>0$, 积分不等式严格.

反向结论在内点 $x$ 处对任意方向 $d$ 考察 $f(x+td)+f(x-td)-2f(x)\ge0$. 除以 $t^2$ 并令 $t\to0$, 得 $d\T\nabla^2f(x)d\ge0$. 开集条件保证足够小的正负方向都可行.
\end{proof}
正定是严格凸的充分条件, 不是必要条件. 例如 $x^4$ 严格凸, 但在0处二阶导数为0. 低维可行集上的凸性也不要求在所有环境方向上 Hessian 半正定.

\originalsection{投影与有限表示}
\begin{theorem}[投影定理]\label{thm:projection}
非空闭凸集 $C\subseteq\R^n$ 上, 每个 $z\in\R^n$ 都有唯一最近点 $P_C(z)$. 点 $p\in C$ 等于 $P_C(z)$ 当且仅当
\begin{equation}\label{eq:projectionvi}
(z-p)\T(x-p)\le0\quad(x\in C).
\end{equation}
\end{theorem}
\begin{proof}
固定 $x_0\in C$. 只须在 $C\cap\{x:\norm{x-z}\le\norm{x_0-z}\}$ 上寻找最小值; 此集非空、闭且有界, 所以紧, 连续目标取到最小值. 函数 $\norm{x-z}^2$ 严格凸, 因而最小点唯一. 最优性条件\ref{thm:vi}与梯度 $2(p-z)$ 给出\eqref{eq:projectionvi}.
\end{proof}
\begin{corollary}\label{cor:nonexpansive}
投影非扩张: $\norm{P_C(u)-P_C(v)}\le\norm{u-v}$. 若 $x\in C$, 则 $\norm{P_C(z)-x}\le\norm{z-x}$.
\end{corollary}
\begin{proof}
记 $p=P_C(u),q=P_C(v)$. 在两次投影条件中分别取 $x=q,p$, 相加得到 $\norm{p-q}^2\le(p-q)\T(u-v)\le\norm{p-q}\norm{u-v}$. 对第二条, 展开 $\norm{z-x}^2=\norm{z-p}^2+\norm{p-x}^2+2(z-p)\T(p-x)$, 最后一项非负.
\end{proof}
\begin{example}
求 $z$ 在半空间 $C=\{x:a\T x\le b\}$ 上的投影, 其中 $a\ne0$.
\end{example}
\begin{solution}
若 $a\T z\le b$, 投影就是 $z$. 否则令 $p=z-[(a\T z-b)/\norm{a}^2]a$. 有 $a\T p=b$, 且 $z-p$ 是 $a$ 的正倍数. 因而对任意 $x\in C$, $(z-p)\T(x-p)\le0$, 投影定理给出答案. 这也是投影梯度算法的一种可直接计算的约束步骤.
\end{solution}

对任意集合 $X$, 其凸包 $\conv X$ 是包含它的最小凸集, 等于全部有限凸组合 $\sum_i\alpha_i x_i$, 其中 $\alpha_i\ge0$, $\sum_i\alpha_i=1$. 仿射包 $\aff X$ 允许系数为任意实数但仍要求和为1. 生成锥 $\cone X$ 允许非负系数而不要求和为1. 凸组合集合本身凸且包含 $X$, 又必须包含于每个包含 $X$ 的凸集, 这证明了两种凸包定义一致.

\begin{theorem}[Carathéodory 定理]\label{thm:caratheodory}
若 $X\subseteq\R^n$ 非空, 则 $\cone X$ 中每个非零向量都能写成至多 $n$ 个来自 $X$ 的线性无关向量的正组合. 每个 $x\in\conv X$ 都能写成至多 $n+1$ 个来自 $X$ 的向量的凸组合.
\end{theorem}
\begin{proof}
选择一个项数最少的正组合 $x=\sum_{i=1}^m\alpha_i x_i$. 若 $x_i$ 线性相关, 有非零系数 $\beta_i$ 使 $\sum_i\beta_i x_i=0$. 必要时整体换号, 使某个 $\beta_i>0$. 取 $t=\min_{\beta_i>0}\alpha_i/\beta_i>0$. 新系数 $\alpha_i-t\beta_i$ 全部非负且至少一个为0, 但仍表示 $x$, 与最少项数矛盾. 因而 $m\le n$.

对凸组合, 提升到 $\R^{n+1}$: $ (x,1)=\sum_i\alpha_i(x_i,1)$. 由锥版本可用至多 $n+1$ 个提升向量正组合表示. 最后一个坐标强制系数和为1, 所以降回原空间得到凸组合.
\end{proof}
\begin{corollary}
紧集的凸包紧. 有限集合生成的锥闭.
\end{corollary}
\begin{proof}
若 $X$ 紧, 每个凸包点可表示为 $n+1$ 项凸组合, 不足的项补零. 因而 $\conv X$ 是紧集 $\Delta_{n+1}\times X^{n+1}$ 在连续映射 $(\alpha,x_1,\ldots,x_{n+1})\mapsto\sum_i\alpha_ix_i$ 下的像, 其中 $\Delta_{n+1}$ 是系数单纯形, 所以紧.

设 $X$ 有限, $y_k\in\cone X$ 收敛. 锥版本使每个 $y_k$ 由 $X$ 的一个线性无关子集正组合表示. 可能的子集有限, 可取子序列让它固定为 $x_1,\ldots,x_m$. 线性无关保证从 $y_k$ 恢复系数的线性逆映射连续, 所以系数收敛到非负值, 极限仍在生成锥中. 对零向量也有结论, 因为原点属于生成锥.
\end{proof}
闭集的凸包却不一定闭. 例如 $X=\{(t,1/t):t>0\}\cup\{(t,-1/t):t>0\}$ 闭; 任意趋向 $t=0$ 的点都向无穷逃逸, 不产生有限漏点. 凸包包含 $(t,0)$ 对所有 $t>0$, 因而原点在其闭包中. 但每个有限凸组合的第一坐标都正, 所以原点不在凸包中. 后面讨论对偶闭性时, 这种“极限存在而表示参数逃向无穷”的现象会再次出现.

\originalchapter{相对内部与无穷远方向}\label{chap:2}
本章对应原第4--5讲, 第39--63页. 凸性不仅控制两点之间, 还控制边界与无穷远处的行为. 相对内部处理低维可行集, 衰退方向则把不有界性变成可计算的代数条件.

\originalsection{相对内部与线段原理}
\begin{definition}
设 $C$ 非空, $L=\aff C-x_0$ 为其仿射包的方向子空间. 点 $x\in C$ 属于相对内部 $\ri C$, 若某个 $r>0$ 满足 $x+\{u\in L:\norm{u}<r\}\subseteq C$.
\end{definition}
例如线段的相对内部是去掉两个端点后的开线段, 单点集的相对内部就是它自己. 相对内部总在自身的仿射包中计算, 不能把它当成环境空间内部.

\begin{theorem}[线段原理]\label{thm:segment}
若 $C$ 凸, $x\in\ri C$, $y\in\cl C$, 则 $(1-t)x+ty\in\ri C$ 对所有 $0\le t<1$ 成立.
\end{theorem}
\begin{proof}
取定义中的相对球半径 $r$. 当 $t=0$ 结论已知. 固定 $0<t<1$, 取 $y_k\in C$ 趋于 $y$, 并令 $z=(1-t)x+ty$. 对任意 $u\in L$ 满足 $\norm{u}<(1-t)r/2$, 选足够大的 $k$ 使 $t\norm{y-y_k}/(1-t)<r/2$. 则
\[
z+u=(1-t)\left(x+\frac{t(y-y_k)+u}{1-t}\right)+ty_k\in C.
\]
括号内点位于以 $x$ 为中心半径 $r$ 的相对球内. 因而 $z$ 也具有包含于 $C$ 的相对球.
\end{proof}
\begin{theorem}\label{thm:ri-basic}
非空凸集的相对内部非空、凸且具有相同的仿射包. 点 $x\in C$ 属于 $\ri C$, 当且仅当每条从 $y\in C$ 到 $x$ 的线段都能越过 $x$ 延长而仍留在 $C$ 中.
\end{theorem}
\begin{proof}
设仿射维数为 $m$. 取 $m+1$ 个仿射无关的点 $z_0,\ldots,z_m\in C$. 它们的单纯形包含于 $C$, 且重心附近有一个相对于 $\aff C$ 的小球, 所以相对内部非空. 该小球也说明相对内部的仿射包等于 $\aff C$. 两个相对内点之间的线段用各自小球作凸组合即可证明仍是相对内点.

相对内点允许沿任何方向小幅移动, 因而线段可以延长. 反过来, 取 $z\in\ri C$. 若 $x=z$, 已得结论; 否则假设的延长性给出 $w=x+\varepsilon(x-z)\in C$. 点 $x$ 是 $z,w$ 间的真内分点, 由线段原理属于 $\ri C$.
\end{proof}
\begin{proposition}
非空凸集满足 $\cl(\ri C)=\cl C$ 及 $\ri(\cl C)=\ri C$. 两个非空凸集有相同闭包当且仅当有相同相对内部. 若 $\ri C\subseteq D\subseteq\cl C$ 且 $D$ 凸, 则 $D$ 也有相同闭包和相对内部.
\end{proposition}
\begin{proof}
固定 $z\in\ri C$. 对 $y\in\cl C$, 线段原理给出 $(1-t)z+ty\in\ri C$ 对 $t<1$ 成立, 令 $t\uparrow1$ 得到闭包等式. 若 $x\in\ri(\cl C)$, 则从 $z$ 到 $x$ 的线段可以在 $\cl C$ 内延长到 $y$, 因而 $x$ 是 $z,y$ 的真内分点, 再由线段原理得 $x\in\ri C$. 其余结论由这两个等式及集合夹逼立即得到.
\end{proof}
\begin{proposition}
若凹函数在非空凸集 $C$ 的相对内点 $x^*$ 取到最小值, 则它在 $C$ 上恒定.
\end{proposition}
\begin{proof}
若某个 $y\in C$ 满足 $f(y)>f(x^*)$, 延长线段至 $w\in C$ 使 $x^*=(1-t)y+tw$, $0<t<1$. 由最小性 $f(w)\ge f(x^*)$, 而凹性给出 $f(x^*)\ge(1-t)f(y)+tf(w)>f(x^*)$, 矛盾.
\end{proof}
因此在可行集上非恒定的线性目标不可能在相对内部取到最小值. 注意说的是“在可行集上非恒定”, 否则环境空间的非零线性函数可能恰好在一个低维可行集上恒定.

\originalsection{相对内部的运算}
\begin{theorem}\label{thm:ri-image}
对非空凸集 $C$ 与线性映射 $A$, 有 $A(\ri C)=\ri(AC)$. 同时 $A(\cl C)\subseteq\cl(AC)$; 若 $C$ 有界, 则后者等号成立.
\end{theorem}
\begin{proof}
线性映射从方向空间 $L$ 到 $AL$ 是满射. 在有限维中可取一个连续线性右逆, 所以相对球的像包含 $AL$ 中的某个相对球. 这证明 $A(\ri C)\subseteq\ri(AC)$.

取 $z\in\ri C$ 及 $y\in\ri(AC)$. 若 $y=Az$, 无须进一步处理. 否则在 $AC$ 中把 $Az$ 到 $y$ 的线段越过 $y$ 延长到 $v=Aw$, 其中 $w\in C$. 于是 $y=(1-t)Az+tAw$, $t<1$, 并且 $(1-t)z+tw\in\ri C$, 给出反向包含.

闭包的正向包含由连续性得到. 若 $C$ 有界, 对 $y_k=Ax_k\to y$, 可选 $x_k$ 的收敛子序列, 极限 $x\in\cl C$ 满足 $Ax=y$, 得反向包含.
\end{proof}
\begin{theorem}\label{thm:ri-intersection}
若 $C_1,C_2$ 非空凸, 则相对内部与笛卡尔积、向量和可交换:
\[
\ri(C_1\times C_2)=\ri C_1\times\ri C_2,\qquad
\ri(C_1+C_2)=\ri C_1+\ri C_2.
\]
若 $\ri C_1\cap\ri C_2\ne\varnothing$, 则
\[
\ri(C_1\cap C_2)=\ri C_1\cap\ri C_2,\qquad
\cl(C_1\cap C_2)=\cl C_1\cap\cl C_2.
\]
\end{theorem}
\begin{proof}
积的结论由相对球定义得到; 和是积在线性映射 $(x_1,x_2)\mapsto x_1+x_2$ 下的像, 所以使用定理\ref{thm:ri-image}. 对交, 固定共同相对内点 $z$. 两个相对球在共同仿射方向空间的交仍含一个球, 故 $z\in\ri(C_1\cap C_2)$. 更一般地, $\ri C_1\cap\ri C_2$ 中任一点也如此.

反向取 $x\in\ri(C_1\cap C_2)$. 延长 $z$ 到 $x$ 的线段到共同可行点 $w$, 于是 $x$ 是 $z,w$ 的真内分点, 在每个 $C_i$ 中都属相对内部. 对闭包, 取 $y\in\cl C_1\cap\cl C_2$, 用 $z$ 与 $y$ 的线段内分点同时逼近, 即得到 $y\in\cl(C_1\cap C_2)$.
\end{proof}
在没有资格条件时, 通常只能得到 $\ri C_1\cap\ri C_2\subseteq\ri(C_1\cap C_2)$ 与 $\cl(C_1\cap C_2)\subseteq\cl C_1\cap\cl C_2$. 取实轴上的 $C_1=(-\infty,0]$, $C_2=[0,\infty)$, 交集的相对内部为 $\{0\}$, 而相对内部的交为空. 若改成两个开半轴, 它们交为空, 闭包的交却仍是 $\{0\}$.

闭包对向量和只有 $\cl C_1+\cl C_2\subseteq\cl(C_1+C_2)$. 若一集有界, 可对该集的表示序列取收敛子序列, 另一表示序列随之收敛, 因而等号成立. 对线性逆像, 若 $A^{-1}(\ri C)\ne\varnothing$, 把逆像看成集合 $\{(x,y):y=Ax\}$ 与 $\R^n\times C$ 的交, 可得
$\ri(A^{-1}C)=A^{-1}(\ri C)$ 及 $\cl(A^{-1}C)=A^{-1}(\cl C)$. 这两个公式需要上述非空资格条件; 不可把原讲义的运算摘要读成无条件规则.

\begin{proposition}[纤维公式]
设 $C\subseteq\R^{n+m}$ 凸, $C_x=\{y:(x,y)\in C\}$, $D=\{x:C_x\ne\varnothing\}$. 则
$\ri C=\{(x,y):x\in\ri D,\ y\in\ri C_x\}$.
\end{proposition}
\begin{proof}
投影定理\ref{thm:ri-image}说明 $x\in\ri D$ 时存在 $y_0$ 使 $(x,y_0)\in\ri C$. 纤维仿射平面 $M_x=\{x\}\times\R^m$ 因而与 $\ri C$ 相交. 用交集公式得 $\ri(M_x\cap C)=M_x\cap\ri C$, 而左边正是 $\{x\}\times\ri C_x$. 结合投影的反向包含即得结论.
\end{proof}

\originalsection{凸函数的连续性与闭包}
\begin{theorem}\label{thm:convex-continuity}
真凸函数在 $\ri(\dom f)$ 上相对连续. 特别地, 有限值凸函数在整个 $\R^n$ 上连续.
\end{theorem}
\begin{proof}
将坐标平移并限制到有效域的仿射包, 只须在内点0证明. 取一个小立方体包含于有效域, 其所有顶点函数值有限; 由凸性, 立方体内函数值不超过这些顶点值的最大值 $M$. 对趋于0的非零点 $x$, 选固定半径 $r$ 使正反两个径向点 $y=rx/\norm{x}$ 与 $z=-rx/\norm{x}$ 都在立方体中. 记 $s=\norm{x}/r$. 凸性给出
\[
f(x)\le(1-s)f(0)+sM,\qquad
f(0)\le\frac{s}{1+s}M+\frac1{1+s}f(x).
\]
后式等价于 $f(x)\ge(1+s)f(0)-sM$. 令 $x\to0$ 即得连续性.
\end{proof}
在边界可能不连续或取无穷. 闭包恰好修补从内部逼近边界时漏掉的上图极限. 对真凸函数定义 $\cl f$ 使 $\epi(\cl f)=\cl(\epi f)$. 它是被 $f$ 逐点控制的最大闭凸函数. 对一般函数, 先将上图取凸包再闭包得到闭凸包函数; 它是被原函数控制的最大闭凸函数, 但可能不真.

\begin{proposition}
真凸函数的闭包仍为真凸函数, 在 $\ri(\dom f)$ 上与原函数相同. 若 $x$ 属于该相对内部, $y\in\dom(\cl f)$, 则
\[
(\cl f)(y)=\lim_{t\downarrow0}f(y+t(x-y)).
\]
函数、其闭包与其闭凸包有相同下确界; 原函数取到下确界的点, 在这些包函数下仍为最优点.
\end{proposition}
\begin{proof}
连续性已说明内部上图不产生新的有限高度. 在内部取一小球, 凸性与局部有界性给出有限的仿射下界; 因而上图闭包不可能包含任何整条竖直线, 所以闭包真. 这一仿射下界也可以由下一章的上图分离构造, 不依赖边界连续性.

令 $g=\cl f$. 线段原理与内部相等性说明 $f(y+t(x-y))=g(y+t(x-y))$. 凸性给出该值不超过 $(1-t)g(y)+tg(x)$, 下半连续性给出极限下界不小于 $g(y)$, 所以极限恰为 $g(y)$.

若 $m=\inf f$ 有限, 整个原上图都位于闭半空间 $w\ge m$ 中, 其凸包与闭包也如此, 所以包函数下确界不小于 $m$; 又因包函数不大于原函数, 有反向不等式. 当 $m=-\infty$ 或 $+\infty$ 时由同一序关系直接得到. 原最优点的函数值被这两个界夹住.
\end{proof}
这里仿射下界的存在将在分离章节给出完整独立证明; 读者可暂把这一句话作为明确的前向引用, 不能用它替代后面的证明.

\originalsection{衰退锥与线性空间}
\begin{definition}
非空凸集 $C$ 的衰退锥为 $\rec C=\{d:x+td\in C\ \text{对一切 }x\in C,t\ge0\}$. 其线性空间为 $L_C=\rec C\cap(-\rec C)$.
\end{definition}
$L_C$ 中的方向允许正反两向无限移动. 因而它是子空间, 且 $C=L_C+(C\cap L_C^\perp)$. 确实, 将 $x\in C$ 正交分解为 $x=l+z$, 沿 $-l\in L_C$ 移动得到 $z\in C\cap L_C^\perp$; 反向包含由衰退方向定义得到. 任何 $S\subseteq L_C$ 的子空间也可替代 $L_C$ 作同样分解.

\begin{theorem}\label{thm:recession}
对非空闭凸集 $C$, 衰退锥闭且凸. 一个方向属于衰退锥, 当且仅当从某个点出发的整条对应射线包含于 $C$. 集合 $C$ 无界当且仅当衰退锥含非零方向. 对非空交集, $\rec(\bigcap_iC_i)=\bigcap_i\rec C_i$. 还满足 $\rec(\ri C)=\rec C$.
\end{theorem}
\begin{proof}
对每个固定的 $x,t$, 条件 $x+td\in C$ 对 $d$ 定义闭凸集, 取交即得第一条. 若 $x_0+sd\in C$ 对 $s\ge0$ 成立, 则对任意 $x\in C,t\ge0$ 及 $s>t$,
$(1-t/s)x+(t/s)(x_0+sd)\in C$. 令 $s\to\infty$ 得 $x+td\in C$, 用到了闭性.

无界时固定 $x_0\in C$, 取 $x_k\in C$ 满足 $r_k=\norm{x_k-x_0}\to\infty$. 单位方向的子序列趋于某个单位向量 $d$. 对固定 $t\ge0$, $x_0+(t/r_k)(x_k-x_0)\in C$ 最终成立, 极限为 $x_0+td$, 从而 $d\in\rec C$. 反向显然, 因为一条非零射线无界. 对交, 若一个方向给交集中的某点提供整条射线, 则分别用射线判据得到它属于每个衰退锥.

若 $d\in\rec C$ 且 $x\in\ri C$, 将 $x$ 附近相对球沿 $td$ 平移, 仍在 $C$, 所以 $x+td\in\ri C$. 反向若 $d\in\rec(\ri C)$, 同一射线位于 $C$, 用闭集射线判据即可.
\end{proof}
闭性条件重要. 非闭凸集可以无界而没有非零衰退方向, 所以不能对任意凸集套用无界性判据.

\originalsection{衰退函数与最优解存在性}
闭真凸函数的衰退函数 $f^\infty$ 定义为上图衰退锥的边界函数, 即 $\epi(f^\infty)=\rec(\epi f)$. 对任意 $x\in\dom f$,
\begin{equation}\label{eq:recession-function}
f^\infty(d)=\sup_{t>0}\frac{f(x+td)-f(x)}t
=\lim_{t\to\infty}\frac{f(x+td)-f(x)}t.
\end{equation}
差商随 $t$ 单调不减, 因为一维凸性使短弦斜率不大于长弦斜率. 若其上确界为 $a$, 就有 $f(x+td)\le f(x)+ta$ 对所有 $t\ge0$ 成立, 等价于 $(d,a)$ 从上图某点出发给出射线. 由闭凸射线判据, 该方向对整个上图有效, 所以上确界与基点无关, 并给出公式. 有限差商的极限不可能为 $-\infty$.

若 $f$ 在全空间可微, 则还可写 $f^\infty(d)=\lim_{t\to\infty}\nabla f(x+td)\T d$. 设 $h(t)=f(x+td)$, 凸性使 $h'(t)$ 单调; 由 $[h(t)-h(0)]/t=t^{-1}\int_0^t h'(s)\,\mathrm ds$, 其平均值趋于同一有限或无穷极限, 即得公式.

定义函数的衰退方向为 $R_f=\{d:f^\infty(d)\le0\}$, 线性空间为 $L_f=R_f\cap(-R_f)$. 前者使函数沿射线不增加, 后者使函数沿整条直线恒定. 原资料第55页的英文摘要出现“nondecreasing”; 由上图的水平衰退方向定义及第57页示意可知应为“不增加”, 本稿依数学定义表述.

\begin{proposition}\label{prop:levelrecession}
闭真凸函数的所有非空下水平集具有相同衰退锥 $R_f$. 若某个非空下水平集有界, 则所有下水平集紧.
\end{proposition}
\begin{proof}
若 $d\in R_f$, 函数沿射线不增加, 所以它保持每个水平集. 反过来, 若某个水平集从 $x$ 沿 $d$ 包含整条射线, 则函数值在这条射线上有一个有限上界. 将该上界减去 $f(x)$ 再除以 $t$, 令 $t\to\infty$, 得 $f^\infty(d)\le0$. 因而衰退锥相同. 下水平集闭凸, 有界当且仅当衰退锥没有非零方向, 故一集有界推出全部有界, 闭有界则紧.
\end{proof}
\begin{example}
计算半正定二次函数 $f(x)=x\T Qx+a\T x+b$ 的衰退方向与恒定方向, 其中 $Q$ 对称半正定.
\end{example}
\begin{solution}
差商为 $2x\T Qd+t\,d\T Qd+a\T d$. 若 $Qd\ne0$, 半正定性给出 $d\T Qd>0$, 极限为 $+\infty$. 若 $Qd=0$, 衰退函数为 $a\T d$. 因而 $R_f=\{d:Qd=0,a\T d\le0\}$, $L_f=\{d:Qd=0,a\T d=0\}$. 特别地, 正定二次函数在所有非零方向上衰退函数为 $+\infty$.
\end{solution}
若多个闭真凸函数在某点同时有限, 则有限和的衰退函数等于其衰退函数之和; 有限最大值则对应衰退函数的最大值. 证明在共同有限的基点使用\eqref{eq:recession-function}; 有限个单调差商的极限可交换有限和或最大值. 对任意上确界族, 还需上确界函数为真并在基点有限, 此时 $\rec(\epi\sup_if_i)=\bigcap_i\rec(\epi f_i)$ 给出 $ (\sup_if_i)^\infty=\sup_i f_i^\infty$. 不能忽略真性而进行无穷算术.

\begin{theorem}[最优解存在性]\label{thm:weierstrass}
设 $X$ 闭非空, $f$ 下半连续且不取 $-\infty$, $X\cap\dom f\ne\varnothing$. 若某个非空可行下水平集有界, 则最优解集合非空且紧. 这包括 $X$ 有界, 或 $f$ 在 $X$ 上强制, 即 $\norm{x_k}\to\infty$ 时 $f(x_k)\to+\infty$.
\end{theorem}
\begin{proof}
选该紧水平集中的点 $x_0$. 所有比 $f(x_0)$ 更好的候选点都在同一个紧集内. 下半连续函数在紧集上取到下确界: 取极小化序列并选收敛子序列, 下半连续性保证极限值不大于下确界, 因而等于它. 这也说明下确界有限, 否则会使极限点取 $-\infty$. 最优解集是闭水平集且包含于原紧集, 所以紧.
\end{proof}
\begin{theorem}
若再设 $X$ 闭凸且 $f$ 闭真凸, 则最优解集合非空且紧, 当且仅当 $\rec X\cap R_f=\{0\}$.
\end{theorem}
\begin{proof}
任一非空可行水平集 $X\cap\{f\le\gamma\}$ 的衰退锥等于 $\rec X\cap R_f$. 若其为 $\{0\}$, 水平集有界, 上一定理适用. 反过来, 若最优解集非空紧, 它的衰退锥为 $\{0\}$. 将它写成可行集与最优值水平集之交, 同一公式给出必要性.
\end{proof}
因此若 $f=\sum_i f_i$ 为真, 而其中某一项在每个非零方向上的衰退函数都为 $+\infty$, 总目标具有紧非空最优解集. 若该项是正定二次函数, 总目标还严格凸, 所以最优解唯一. 同样的存在性论证适用于有限最大值, 但最大值未必严格凸, 不可自动添加唯一性结论.

\originalchapter{闭性、分离与共轭}\label{chap:3}
本章对应原第6--8讲的闭集交、分离与共轭部分, 第64--103页. 三个看似不同的问题, 即最优解是否存在、线性像是否闭、消去变量后的值函数是否闭, 都可转化为嵌套闭集的交是否非空.

\originalsection{嵌套闭集与回缩性}
若 $C_{k+1}\subseteq C_k$ 都非空闭, 且某个 $C_k$ 紧, 则交集非空: 在每个后续集合选点, 紧性给收敛子序列, 嵌套性使极限属于每个固定集合. 不紧时这一结论失败, 如 $C_k=[k,\infty)$.

\begin{definition}
对嵌套非空闭凸集 $C_k$, 若 $x_k\in C_k$, $\norm{x_k}\to\infty$ 且 $x_k/\norm{x_k}\to d/\norm d$, $d\ne0$, 称其为渐近序列. 若最终 $x_k-d\in C_k$, 称这一渐近序列可回缩. 若每个渐近序列都可回缩, 称集合序列可回缩. 单个闭凸集的可回缩性由常值集合序列定义.
\end{definition}
任何无界的选择序列都存在渐近子序列. 其方向属于每个 $\rec C_j$: 固定 $j$, 后续点都在 $C_j$, 对规范化线段取极限即可使用衰退锥的证明. 回缩性进一步要求在已很远的可行点上向反方向移动固定距离仍可行, 它比“正方向可以无限移动”更强.

\begin{lemma}
闭半空间、多面体是可回缩集. 有限个可回缩集合序列的逐项交及笛卡尔积仍可回缩. 凸紧集与可回缩凸集的向量和可回缩.
\end{lemma}
\begin{proof}
半空间记为 $a\T x\le b$. 对渐近方向 $d$, 必有 $a\T d\le0$. 若等于0, 平移 $-d$ 不改变不等式; 若小于0, 则 $a\T x_k\to-\infty$, 最终也允许平移. 多面体只有有限个条件, 逐一应用即可. 交集的渐近序列是每个组成序列的渐近序列, 故可以同时回缩.

对积, 各分量中方向非零的项按相应增长比例作渐近序列, 使用该分量的回缩性; 方向为零的项无须移动. 对和, 写 $x_k=u_k+v_k$, 紧集分量 $u_k$ 有界. 则 $v_k$ 具有与 $x_k$ 相同的渐近方向, 最终 $v_k-d$ 可行, 所以 $x_k-d=u_k+(v_k-d)$ 可行.
\end{proof}
\begin{theorem}\label{thm:retractive-intersection}
可回缩的嵌套非空闭凸集合序列具有非空交集. 更一般地, 设 $X$ 为可回缩闭凸集, 每个 $X\cap C_k$ 非空, 并令 $R=\bigcap_k\rec C_k$, $L=\bigcap_kL_{C_k}$. 若 $\rec X\cap R\subseteq L$, 则 $\bigcap_k(X\cap C_k)$ 非空.
\end{theorem}
\begin{proof}
在每个 $C_k$ 中取离原点最近的点 $p_k$. 若此序列有界, 收敛子序列的极限属于交集. 若无界, 取渐近子序列, 方向为 $d$. 回缩性给出最终 $p_k-d\in C_k$, 但
$\norm{p_k-d}^2=\norm{p_k}^2-2p_k\T d+\norm d^2<\norm{p_k}^2$, 因为 $p_k\T d\to+\infty$. 这与最近点定义矛盾.

对一般情况, 渐近方向 $d$ 属于 $\rec X\cap R$, 因而属于每个 $L_{C_k}$. 所以 $x_k-d\in C_k$ 始终成立; 而 $X$ 的回缩性保证它最终属于 $X$. 于是 $X\cap C_k$ 序列可回缩, 用第一部分得结论.
\end{proof}
$X$ 的回缩性不能删掉. 取闭凸集 $X=\{(x,y):x>0,y\ge1/x\}$ 与 $C_k=\{(x,y):y\le1/k\}$, 各个交非空, 全部交却为空. $\rec X\cap\bigcap_k\rec C_k$ 是非负水平射线, 包含于每个 $C_k$ 的线性空间, 但 $X$ 不可回缩: 点 $(k,1/k)$ 向左平移1后不再属于 $X$.

\begin{corollary}\label{cor:qp-attainment}
在非空多面体上, 半正定二次函数只要下确界有限, 就取到最小值. 线性规划是其特例.
\end{corollary}
\begin{proof}
写 $f(x)=x\T Qx+c\T x$, $Q\succeq0$, 取 $\gamma_k\downarrow f^*$, $\gamma_k>f^*$. 下水平集的共同衰退锥为 $R=\{d:Qd=0,c\T d\le0\}$, 共同线性空间为 $L=\{d:Qd=0,c\T d=0\}$. 若 $d\in\rec X\cap R$ 且 $c\T d<0$, 任一可行点沿 $d$ 使目标趋于 $-\infty$, 与有限下确界矛盾. 因而 $\rec X\cap R\subseteq L$. 多面体可回缩, 上一定理使嵌套可行水平集的交非空, 交中点恰为最优点.
\end{proof}

\originalsection{线性像与部分极小化}
\begin{theorem}\label{thm:closed-image}
设 $C$ 非空闭凸, $A$ 线性. 若 $\rec C\cap\ker A\subseteq L_C$, 则 $AC$ 闭. 若 $X$ 是可回缩闭凸集, 则条件 $\rec X\cap\rec C\cap\ker A\subseteq L_C$ 保证 $A(X\cap C)$ 闭. 特别地, 多面体的线性像闭.
\end{theorem}
\begin{proof}
取 $y_k=Ax_k\to y$. 令 $\varepsilon_k\downarrow0$ 且选择子序列使 $\norm{y_k-y}\le\varepsilon_k$, 再令 $D_k=\{x\in C:\norm{Ax-y}\le\varepsilon_k\}$. 各 $D_k$ 非空闭凸, 共同衰退锥为 $\rec C\cap\ker A$, 线性空间包含 $L_C\cap\ker A$. 条件及定理\ref{thm:retractive-intersection}, 取 $X=\R^n$, 保证交非空; 交中点满足 $Ax=y$.

第二种情况在同一 $D_k$ 外再交 $X$, 并用一般回缩交定理. 多面体的线性像可令 $C=\R^n$, 此时 $L_C=\R^n$, 条件自动成立.
\end{proof}
原第1讲第13页的例子也可精确写出: $C_1=\{(u,v):u>0,v>0,uv\ge1\}$ 与 $C_2=\{(0,v):v\in\R\}$ 都闭凸, 但 $C_1+C_2=\{(u,v):u>0,v\in\R\}$ 不闭. 前者闭是因为趋于 $u=0$ 时 $v$ 必趋无穷, 不产生有限漏点; 与竖直线相加后这个阻止有限极限的条件消失. 同样, 将 $C_1$ 投影到第一坐标只得到开半轴.

最简单的充分条件是 $\rec C\cap\ker A=\{0\}$, 这时预像截集本身紧. 向量和是积在线性求和映射下的像; 因而若 $d_i\in\rec C_i$ 且 $\sum_i d_i=0$ 只能在各 $d_i=0$ 时成立, 则 $\sum_iC_i$ 闭. 特别地, 一个闭凸集与一个凸紧集的和闭, 以及 $\rec C_1\cap\rec C_2=\{0\}$ 时 $C_1-C_2$ 闭.

\begin{theorem}\label{thm:partial-min}
设 $F:\R^{n+m}\to(-\infty,+\infty]$ 闭真凸, $p(x)=\inf_zF(x,z)$. 则 $p$ 凸. 若某个纤维水平集 $\{z:F(x_0,z)\le\gamma\}$ 非空紧, 则 $p$ 闭且真, 且对每个 $x\in\dom p$, 纤维极小值都取到, 极小点集合非空紧.
\end{theorem}
\begin{proof}
对两个近似极小点作凸组合, 目标误差可任意小, 得 $p$ 凸. 记 $P(x,z,w)=(x,w)$. 总有 $P(\epi F)\subseteq\epi p\subseteq\cl(P(\epi F))$: 第二个包含由在固定 $x$ 处取任意接近下确界的 $z$ 得到.

某个纤维水平集紧, 说明上图没有非零水平纤维衰退方向 $(0,d,0)$. 否则该方向会使这一个紧集包含非零射线. 因而 $\rec(\epi F)\cap\ker P=\{0\}$, 线性像定理使 $P(\epi F)$ 闭, 两个包含变为等号. 对任一非空纤维函数 $F(x,\cdot)$, 它的所有非空水平集衰退锥也都为零; 存在性定理使其最优值有限并取到, 极小点集合紧. 因而 $p$ 不取 $-\infty$, 又在 $x_0$ 有限, 所以真, 上图闭则说明它闭.
\end{proof}
\begin{example}
原第69页反例取 $F(x,z)=e^{-\sqrt{xz}}$ 当 $x,z\ge0$, 其余处取 $+\infty$. 求 $p(x)=\inf_zF(x,z)$ 并解释其闭性.
\end{example}
\begin{solution}
几何平均 $\sqrt{xz}$ 在非负象限上凹, 所以 $-\sqrt{xz}$ 凸; 指数函数凸且单调递增, 复合仍凸. 非负象限闭, 函数在其中连续, 扩展后闭. 当 $x>0$, 令 $z\to\infty$ 得下确界0但不取到; 当 $x=0$, 值恒为1; 当 $x<0$, 无可行 $z$. 因而 $p(x)=0$ 对 $x>0$, $p(0)=1$, $p(x)=+\infty$ 对 $x<0$. 正数列趋于0使下半连续性失败. 此处每个非空纤维水平集都缺少定理要求的紧性.
\end{solution}

\originalsection{分离超平面}
\begin{definition}
超平面 $H=\{x:a\T x=b\}$ 要求 $a\ne0$. 若 $a\T x_1\le b\le a\T x_2$ 对所有 $x_i\in C_i$ 成立, 称 $H$ 分离两集. 若两边都为严格不等式, 称严格分离. 若分离且两集不全包含在 $H$ 中, 称适当分离. 若 $x\in\cl C$ 且经过 $x$ 的超平面把 $C$ 放在一侧, 称其在 $x$ 支持 $C$.
\end{definition}
\begin{theorem}\label{thm:separation}
非空凸集 $C$ 的非内点 $x$ 可与 $C$ 由经过 $x$ 的超平面分离. 两个不交非空凸集可分离. 若其中一个闭而另一个紧, 则可严格分离. 两个非空凸集可适当分离, 当且仅当它们的相对内部不交.
\end{theorem}
\begin{proof}
先对闭凸集及外点 $x$ 使用投影 $p=P_C(x)$. 向量 $a=p-x\ne0$ 满足 $a\T(y-p)\ge0$ 对 $y\in C$, 并且 $a\T p>a\T x$, 所以给出严格分离. 对边界点, 取外点 $x_k\to x$, 将各分离法向量规范化为单位向量并取收敛子序列, 极限仍为非零单位向量, 使 $a\T(y-x)\ge0$. 非闭集以闭包代替. 若仿射包是低维, 可取垂直于仿射包的法向量; 若 $x$ 不在闭包, 使用严格分离后把平面移至 $x$ 即可.

对不交集合, 在凸差集 $C_2-C_1$ 与原点间用第一条, 得 $a\T(x_2-x_1)\ge0$. 因而存在有限 $b$ 夹在 $\sup_{C_1}a\T x$ 与 $\inf_{C_2}a\T x$ 之间. 闭集与紧集的差闭且不含0; 投影原点到差集, 得到严格正的间隔, 可将平面放在间隔中点.

若相对内部不交, 则 $0\notin\ri(C_2-C_1)$, 因为相对内部与和交换. 在差集的仿射包内应用支持构造, 或当原点不在该仿射包时严格分离, 可选一个不在整个差集上恒为0的分离泛函, 因而适当分离成立. 反向若两集有共同相对内点 $z$, 任何分离泛函在该点取得两侧等号. 因为它在每个集合中都在相对内部取线性极值, 必须在两集各自恒定, 所以只能把两集都包含于平面中, 无法适当分离.
\end{proof}
适当分离在差集的仿射包内构造的细节如下. 若0位于相对边界, 外点可在仿射包内选择, 单位法向量也在其方向空间内; 极限法向量非零, 而整个差集仿射张成该空间, 所以它不能在差集上恒为0. 这避免了只取环境空间中无信息的垂直法向量.

\begin{corollary}
$\cl(\conv C)$ 等于所有包含 $C$ 的闭半空间的交.
\end{corollary}
\begin{proof}
每个包含 $C$ 的闭半空间都包含其闭凸包. 若 $x$ 在闭凸包外, 点与闭凸包的严格分离提供一个包含该包而排除 $x$ 的闭半空间, 因而反向包含成立.
\end{proof}

\originalsection{多面体分离的加强}
后续线性约束的对偶性不要求所有约束严格可行, 其基础是多面体分离可以加强到“不把非多面体一侧整个放在平面上”. 为说明有限性在哪里起作用, 先建立多面体的有限生成表示.

\begin{lemma}\label{lem:poly-finite}
每个多面体锥都由有限个向量生成. 每个非空多面体都可写成有限点的凸包与有限向量生成锥之和.
\end{lemma}
\begin{proof}
全空间由 $\pm e_i$ 有限生成. 若 $K=\cone\{v_i\}$, 则与半空间 $a\T x\le0$ 相交后, 可用满足 $a\T v_i\le0$ 的原生成向量及每一正负配对向量
\[
(-a\T v_j)v_i+(a\T v_i)v_j\quad(a\T v_i>0, a\T v_j<0)
\]
生成. 为验证, 将任意非负组合中正的 $a$ 分量与负的 $a$ 分量逐量配对抵消, 每次抵消正对应上述配对向量, 剩余只含非正分量. 有限次加入定义锥的半空间, 即得第一条; 等式用两个半空间表示.

对多面体 $P=\{x:Bx\le b,Ex=e\}$, 提升为锥 $K=\{(x,t):Bx\le bt,Ex=et,t\ge0\}$. 将有限生成向量中 $t>0$ 的项按其最后坐标归一到1, 其横坐标形成有限点集; $t=0$ 的项形成有限方向集. 切取 $t=1$ 时, 正高度项的系数和为1, 零高度项仍为非负组合, 得所述表示.
\end{proof}
\begin{theorem}\label{thm:polyseparation}
若 $C$ 非空凸, $P$ 非空多面体, 且 $\ri C\cap P=\varnothing$, 则存在分离超平面, 并且 $C$ 不全包含在该平面中.
\end{theorem}
\begin{proof}
先适当分离 $C,P$. 若平面不包含整个 $C$, 已得结论. 否则 $C$ 在平面上, 而 $P$ 在一侧且有点严格位于该侧. 将 $P$ 限制到与平面的交, 这是 $P$ 的真面, 维数下降; 若交为空, 多面体的有限生成点在该泛函上有统一正余量, 可将平面朝多面体方向小幅平移, 使它严格分离并不包含 $C$. 对剩余面与 $C$ 再作适当分离, 必要时重复. 每次若仍包含整个 $C$, 就再取一个真面. 维数有限, 不可能无限下降, 因而最终得到一个不包含整个 $C$ 的分离泛函.

还须把最后泛函扩展为对原 $P$ 的分离. 在引理\ref{lem:poly-finite}的有限点与射线表示上, 前一个泛函对不属于所取面的生成点或方向有严格正余量; 对面的生成项, 新泛函已经非负. 因而把新泛函乘以足够小的正数再加到前一个泛函上, 对全部有限生成项仍非负. 在 $C$ 上前一个泛函恒定, 所以加入新泛函后保留其非恒定性. 逐级反向合并, 就得到对原 $P$ 的分离且不包含整个 $C$ 的平面.
\end{proof}
有限生成项是此证明不可缺少的条件. 对非多面体, 无限生成方向的严格余量可能趋于0, 无法选取统一的小系数.

\originalsection{仿射下界与共轭函数}
\begin{lemma}\label{lem:affine-minorant}
每个真凸函数都具有有限的仿射下界. 若 $f$ 闭真凸且 $r<f(x)$, 则存在仿射函数 $\ell$ 满足 $\ell\le f$ 且 $\ell(x)>r$.
\end{lemma}
\begin{proof}
对真凸函数, 取 $x_0\in\ri\dom f$. 它在相对小球上有限连续. 在上图相对于 $\aff(\dom f)\times\R$ 的空间内, 分离上图和 $(x_0,f(x_0)-1)$. 分离泛函写成 $a\T x+\beta w\ge c$, 由于上图可向上无限移动, $\beta\ge0$. 若 $\beta=0$, 相对内点 $x_0$ 处的严格分离不可能成立, 所以可取得 $\beta>0$. 除以 $\beta$ 得仿射下界, 并把横向泛函从仿射包延拓到全空间. 这也完成上一章闭包真性的前向引用.

对于闭真凸函数, 点 $(x,r)$ 在闭上图外, 严格分离得到 $(a,\beta)$, $\beta\ge0$. 若 $\beta>0$ 直接得到所需仿射函数. 若 $\beta=0$, 将该严格分离不等式加上已有仿射下界对应不等式的小正倍数. 严格间隔在小扰动下仍为正, 竖直系数变为正, 除以它便得到 $\ell(x)>r$.
\end{proof}
\begin{definition}
函数 $f$ 的凸共轭为 $f^*(y)=\sup_x\{y\T x-f(x)\}$. 它记录所有斜率为 $y$ 的仿射下界中最大的截距: $y\T x-f^*(y)\le f(x)$. 二次共轭记为 $f^{**}$.
\end{definition}
共轭是仿射函数族的上确界, 所以始终闭凸, 即使 $f$ 本身不凸. 对真凸函数, 仿射下界说明其共轭至少在一个点有限, 因而共轭真; 又因原函数在某点有限, 共轭从不取 $-\infty$.

\begin{theorem}[共轭定理]\label{thm:biconjugate}
总有 $f^{**}\le f$. 若 $f$ 闭真凸, 则 $f^{**}=f$. 更一般地, 若 $f$ 的闭凸包函数真, 则 $f^{**}$ 等于这一闭凸包函数.
\end{theorem}
\begin{proof}
由 $f^*(y)\ge y\T x-f(x)$ 得 $y\T x-f^*(y)\le f(x)$, 对 $y$ 取上确界得第一条. 若 $f$ 闭真凸, 对任意实数 $r<f(x)$, 引理\ref{lem:affine-minorant}给出 $a\T z+b\le f(z)$ 且 $a\T x+b>r$. 由共轭定义, $f^*(a)\le-b$, 所以 $f^{**}(x)\ge a\T x-f^*(a)\ge a\T x+b>r$. 对所有 $r<f(x)$ 成立, 包括 $f(x)=+\infty$ 的情形, 因而 $f^{**}\ge f$.

共轭的定义只取线性泛函在上图上的下界, 对上图取凸包或闭包不改变这些下界. 因而 $f$ 与其闭凸包有相同共轭; 对真闭凸包使用已证明结论即可.
\end{proof}
非真性不能忽略. 令 $f(x)=-\infty$ 对 $x\le0$, $f(x)=+\infty$ 对 $x>0$, 它闭凸, 但 $f^*\equiv+\infty$, $f^{**}\equiv-\infty$, 不等于 $f$.

\begin{example}
\label{ex:conjugate-models}
计算仿射函数、绝对值、幂函数及正定二次函数的共轭.
\end{example}
\begin{solution}
对 $f(x)=a\T x-b$, 有 $f^*(y)=b$ 若 $y=a$, 否则 $+\infty$. 因为 $y\ne a$ 时沿 $y-a$ 方向的线性表达式无界. 对 $f(x)=|x|$, 当 $|y|\le1$ 时 $yx-|x|\le0$ 且0处取到; 当 $|y|>1$ 时沿对应符号方向无界, 所以共轭是区间 $[-1,1]$ 的指示函数.

取 $1<p<\infty$, $q=p/(p-1)$. 对 $f(x)=\sum_i|x_i|^p/p$, 各坐标独立, 最大化条件为 $y_i=\operatorname{sgn}(x_i)|x_i|^{p-1}$, 代回得到 $f^*(y)=\sum_i|y_i|^q/q$.

对 $f(x)=\tfrac12x\T Qx+a\T x+b$, $Q\succ0$, 最大化的一阶条件给出 $x=Q^{-1}(y-a)$, 完成平方得 $f^*(y)=\tfrac12(y-a)\T Q^{-1}(y-a)-b$. 一维 $f(x)=cx^2/2$ 的共轭因此为 $y^2/(2c)$.
\end{solution}
若 $f(x)=p(A(x-c))+a\T x+b$ 且 $A$ 可逆, 变量代换 $z=A(x-c)$ 给出
$f^*(y)=p^*(A^{-\mathsf T}(y-a))+c\T(y-a)-b$. 可逆条件是这里直接代换公式的一部分, 不可用于一般降维映射.

\originalsection{支撑函数、极锥与 Farkas 引理}
集合 $X$ 的支撑函数为 $\sigma_X(y)=\sup_{x\in X}y\T x$, 它是指示函数的共轭. $X,\cl X,\conv X,\cl\conv X$ 有相同支撑函数; 它的上图是闭凸锥. 支撑值是投影方向上的上确界, 可能为无穷或未取到, 因而不能总画成一个实际“最远点”.

本书使用极锥约定 $C^\circ=\{y:y\T x\le0\ \forall x\in C\}$. 在锥规划中将另用非负约定的对偶锥 $C^+=-C^\circ$, 两者的符号不可混用. 若 $C$ 是包含0的锥, 则 $\sigma_C=\ind{C^\circ}$, 因为任何正内积都可按比例放大至无穷. 共轭定理给出双极公式 $(C^\circ)^\circ=\cl\conv C$.

\begin{theorem}[Farkas 引理]\label{thm:farkas}
设 $A$ 的列为 $a_1,\ldots,a_m$. 方程 $A\lambda=b$, $\lambda\ge0$ 可解, 当且仅当对每个 $y$, $A\T y\le0$ 都蕴含 $b\T y\le0$. 等价地, 下列两者恰有一个成立: $b$ 是各列的非负组合; 存在 $y$ 使 $A\T y\le0$ 而 $b\T y>0$.
\end{theorem}
\begin{proof}
必要性由 $b\T y=\lambda\T A\T y\le0$. 对充分性, 列向量生成的锥闭, 已在 Carathéodory 推论中证明. 若 $b$ 不在该锥中, 对点与闭凸锥严格分离; 锥的齐次性强制分离泛函在整个锥上非正, 而在 $b$ 上严格正, 正是第二种条件. 两种条件不可能同时成立.
\end{proof}
这把“找一组非负系数”和“找一个证实不可能的线性泛函”连接起来. 在线性规划中, 前者给出对偶乘子, 后者给出降低目标的可行方向; 这一几何关系将在下一章展开.

\originalchapter{扰动、对偶与最优性}\label{chap:4}
本章对应原第1讲对偶导论以及第8--11讲的对偶理论, 原PDF第6--16页、第104--157页. 对偶的基本问题是: 能否用一个容易计算的下界逼近原最优值, 下界是否取到, 以及下界与原值相等时如何恢复解.

\originalsection{最小共点与最大截距}
对非空 $M\subseteq\R^{r+1}$, 把点记为 $(u,w)$. 定义最小共点值 $w^*=\inf\{w:(0,w)\in M\}$, 以及
\[
q(\mu)=\inf_{(u,w)\in M}(w+\mu\T u),\qquad q^*=\sup_{\mu\in\R^r}q(\mu).
\]
超平面 $w+\mu\T u=\xi$ 将 $M$ 放在上方的条件是 $\xi\le q(\mu)$, 所以 $q^*$ 是这类非竖直超平面在 $w$ 轴上的最大截距. 原课程称为 min common / max crossing, 简称 MC/MC. 它不是额外的计算技巧, 而是统一各种对偶问题的几何形式.

\begin{proposition}[弱对偶]\label{prop:weakdual}
$q$ 凹且上半连续, 总有 $q(\mu)\le w^*$ 及 $q^*\le w^*$. 将 $M$ 向上补齐、取凸包或闭包都不改变 $q$.
\end{proposition}
\begin{proof}
$q$ 是关于 $\mu$ 的仿射函数族的下确界, 所以凹且上半连续. 在下确界中只取 $u=0$ 的点即得 $q(\mu)\le w^*$. 线性泛函的下确界在凸包与闭包上不变; 向上增加 $w$ 也不降低下确界.
\end{proof}
若 $M=\epi p$, 则 $w^*=p(0)$, $q(\mu)=-p^*(-\mu)$, $q^*=p^{**}(0)$. 对偶把扰动函数中原点处的高度替换为其闭凸包高度. 这直接说明闭性、凸性为何进入强对偶条件.

\begin{theorem}[MC/MC 强对偶条件]\label{thm:mc1}
设 $w^*<+\infty$, 向上补齐集合 $M^+=\{(u,v):\exists(u,w)\in M, v\ge w\}$ 凸. 则 $q^*=w^*$ 当且仅当对每个 $(u_k,w_k)\in M$, $u_k\to0$, 都有 $w^*\le\liminf_kw_k$.
\end{theorem}
\begin{proof}
必要性由 $q(\mu)\le w_k+\mu\T u_k$ 取下极限, 再对 $\mu$ 取上确界得到. 若 $w^*=-\infty$, 弱对偶已经给强对偶, 因而只须考虑有限 $w^*$.

设序列条件成立. 对任意 $\varepsilon>0$, 点 $(0,w^*-\varepsilon)$ 不在 $\cl M^+$, 否则近似序列将违反条件. 闭凸集 $\cl M^+$ 也不含竖直直线: 一条这样的直线使向下方向成为其衰退方向, 而 $(0,w^*)\in\cl M^+$, 从而使所有 $(0,w^*-\varepsilon)$ 都在闭包中. 不含竖直直线的凸集有一个非竖直半空间下界; 这是上一章仿射下界论证对一般集合的同一版本. 若严格分离当前外点的平面恰为竖直, 加上该非竖直下界的小正倍数, 仍严格分离且竖直系数为正. 规范化为1得到 $w^*-\varepsilon<q(\mu)\le w^*$. 令 $\varepsilon\downarrow0$ 得结论.
\end{proof}
补充上段使用的集合版本: 闭凸集是包含它的全部闭半空间之交. 若全部界定半空间都竖直, 该交必含每个已有点上方与下方的整条竖直线. 因而至少一个界定半空间非竖直. 向上补齐使其竖直系数非负, 非零时即为正. 对非闭集, 其相对内部与闭包具有同一衰退锥, 所以闭包若含竖直线, 原集合的相对内部也含竖直线; 无竖直线性质因而保留.

\begin{theorem}[乘子存在与紧性]\label{thm:mc2}
在上述凸性条件下设 $w^*$ 有限, $D=\{u:\exists w,(u,w)\in M\}$. 若 $0\in\ri D$, 则强对偶成立且有对偶最优解. 对偶最优解集非空紧当且仅当 $0\in\operatorname{int}D$.
\end{theorem}
\begin{proof}
$(0,w^*)$ 不在 $\ri M^+$, 否则向下小幅移动仍可行. 适当分离给出
$\beta w^*\le\mu\T u+\beta w$ 对所有 $(u,w)\in M^+$, 且泛函不在整个集合上恒定. 向上封闭使 $\beta\ge0$. 若 $\beta=0$, 则 $\mu\T u$ 在 $D$ 的相对内点0取得最小值, 必在整个 $D$ 恒为0, 与适当性矛盾. 故 $\beta>0$, 除以它得到 $q(\mu/\beta)\ge w^*$, 弱对偶使等号成立.

若0为内点, 选 $\pm re_i\in D$ 及相应有限高度 $w_i^\pm$. 对任意最优乘子, $w^*\le w_i^\pm\pm r\mu_i$, 因而每个分量有上下界. 最优集由上半连续性闭, 所以紧. 反过来, 若有最优乘子而0不是内点, 支持 $D$ 于0得到非零 $a$ 使 $a\T u\ge0$ 对所有 $u\in D$. 则 $q(\mu+ta)\ge q(\mu)=w^*$ 对 $t\ge0$, 弱对偶又强制等号, 最优集含非零射线, 不紧.
\end{proof}
这里强对偶与乘子存在是不同结论. 扰动函数可能在原点闭却没有非竖直支持平面, 此时最优下界只能由乘子趋向无穷取到极限.

\begin{theorem}[多面体加强]\label{thm:mc3}
设 $\widetilde M$ 凸, $P$ 为非空多面体, $M=\widetilde M-(P\times\{0\})$, $w^*$ 有限. 若 $\ri\widetilde D\cap P\ne\varnothing$, 其中 $\widetilde D$ 是 $\widetilde M$ 的横坐标投影, 则强对偶成立且有最优乘子. 所有最优乘子满足 $\mu\T d\le0$ 对 $d\in\rec P$.
\end{theorem}
\begin{proof}
令 $C_1=\{(u,v):\exists(u,w)\in\widetilde M,v>w\}$, $C_2=P\times\{w^*\}$. 两者不交, 否则 $M$ 在竖轴上有低于 $w^*$ 的点. 多面体加强分离使 $C_2$ 位于下侧, 并且平面不包含整个 $C_1$. 写分离不等式为 $\mu\T z+\beta w^*\le\mu\T u+\beta v$. 有 $\beta\ge0$. 若 $\beta=0$, 在 $\ri\widetilde D\cap P$ 的共同点, $\mu\T u$ 在 $\widetilde D$ 上取得相对内部极小, 因而恒定, 违背分离加强性质. 故可规范化 $\beta=1$.

对 $z\in P$ 与 $(u,v)\in C_1$ 取下确界, 得 $w^*\le\inf(v+\mu\T(u-z))=q(\mu)$; 弱对偶给等号. 若某个 $d\in\rec P$ 满足 $\mu\T d>0$, 令 $z$ 沿 $d$ 趋于无穷就使上述下确界为 $-\infty$, 不可能最优.
\end{proof}
与相对内部版本相比, 这里把 $\ri P$ 放宽为 $P$ 本身. 线性不等式允许等号处可行而无需严格可行, 正是这一放宽的后果.

\originalsection{拉格朗日对偶}
考虑 $\inf\{f(x):x\in X,g(x)\le0\}$, 其中 $g=(g_1,\ldots,g_r)$. 定义扰动函数
$p(u)=\inf\{f(x):x\in X,g(x)\le u\}$. 引入拉格朗日函数 $L(x,\mu)=f(x)+\mu\T g(x)$. 对 $\mu\ge0$,
\[
q(\mu)=\inf_u(p(u)+\mu\T u)=\inf_{x\in X}L(x,\mu).
\]
若任一分量 $\mu_j<0$, 可把相应 $u_j$ 增大至无穷, 所以对偶值为 $-\infty$. 对偶问题因而在非负正交域上最大化 $q$. 原问题不必凸, 弱对偶及 $q$ 的凹性仍成立; 凸性主要用于使下界达到原最优值.

\begin{theorem}[非线性 Farkas 引理]\label{thm:nonlinear-farkas}
设 $X$ 凸, $f,g_j$ 在 $X$ 上有限凸, 且 $g(x)\le0$ 时 $f(x)\ge0$. 若存在 $\bar x\in X$ 使 $g(\bar x)<0$, 则存在 $\mu\ge0$ 使 $f(x)+\mu\T g(x)\ge0$ 对所有 $x\in X$ 成立. 此类乘子集合非空紧; 反过来, 其非空紧性蕴含严格可行点存在. 若全部 $g_j$ 仿射, 则仅要求 $\ri X$ 中有一个可行点就足以保证乘子存在.
\end{theorem}
\begin{proof}
取 $M=\{(u,w):\exists x\in X,g(x)\le u,f(x)\le w\}$. 它凸, 严格可行点使0为其横向投影内点, 而 $w^*\ge0$ 且有限. 定理\ref{thm:mc2}给出 $q(\mu)=w^*\ge0$, 并强制 $\mu\ge0$. 任一证书乘子满足 $0\le f(\bar x)+\mu\T g(\bar x)$, 因而 $\sum_j\mu_j\le f(\bar x)/\min_j(-g_j(\bar x))$. 证书集合闭, 故紧.

若没有严格可行点, 凸上闭集合 $g(X)+\R_+^r$ 与严格负正交域不交. 分离两集给非零 $d\ge0$ 使 $d\T g(x)\ge0$ 对所有 $x\in X$. 任一证书 $\mu$ 因而可加上 $td$, $t\ge0$, 仍是证书, 所以非空证书集不紧. 仿射约束的版本把 $M$ 写成仿射像上图加正交域, 用定理\ref{thm:mc3}及 $\ri X$ 中可行点得到证书.
\end{proof}
\begin{theorem}[凸规划强对偶]\label{thm:slater}
设 $X,f,g_j$ 如上, 原最优值有限. 若存在严格可行点, 或全部 $g_j$ 仿射且 $\ri X$ 中有可行点, 则 $q^*=f^*$ 且有最优乘子. 在严格可行条件下最优乘子集还紧.
\end{theorem}
\begin{proof}
将上一引理用于 $f-f^*$. 得 $f^*\le f(x)+\mu^{*\mathsf T} g(x)$ 对全部 $x\in X$. 取下确界得 $f^*\le q(\mu^*)$, 弱对偶给反向不等式. 对最优乘子, 严格可行点同样给出 $\sum_j\mu_j\le[f(\bar x)-f^*]/\min_j(-g_j(\bar x))$, 所以最优集紧.
\end{proof}
这一结论保证对偶解, 并不单独保证原解存在; 原解的存在仍由上一章的闭性、紧性或多面体结构判断.

\begin{theorem}[原始--对偶最优性]\label{thm:lagrangian-optimal}
可行点 $x^*$ 与 $\mu^*\ge0$ 同时最优且具有零对偶间隙, 当且仅当
\[
x^*\in\argmin_{x\in X}L(x,\mu^*),\qquad
\mu_j^*g_j(x^*)=0\quad(j=1,\ldots,r).
\]
\end{theorem}
\begin{proof}
若二者最优且零间隙, 有
$f^*=q(\mu^*)\le L(x^*,\mu^*)=f(x^*)+\sum_j\mu_j^*g_j(x^*)\le f(x^*)=f^*$.
全部等号迫使 $x^*$ 达到拉格朗日下确界, 且每个非正项为0. 反过来, 两条件使 $q(\mu^*)=f(x^*)$, 弱对偶夹逼给出二者最优及零间隙.
\end{proof}
乘子为正的约束必须取等号; 有松弛的约束必须乘子为零. 反向均不必成立, 因为一个活跃约束也可能乘子为零.

等式约束 $Ax=b$ 的乘子 $\lambda$ 无符号限制, 因为等式可写成两组反向不等式, 其非负乘子差可为任意实数. 拉格朗日函数为 $f(x)+\mu\T g(x)+\lambda\T(Ax-b)$. 纯等式的强对偶条件是有限原值与 $\ri X$ 中有等式可行点. 混合约束可用以下条件: 在等式可行集中有严格不等式可行点, 且 $\ri X$ 与等式平面相交. 取严格可行点与后一相对内点的合适凸组合, 可在 $\ri X$ 内保留严格不等式并满足等式, 从而同时应用相对内部与严格可行论证. 最优性条件增加等式可行性, 但等式没有额外的互补松弛条件.

\originalsection{线性与二次规划对偶}
\begin{theorem}\label{thm:lpdual}
线性规划 $\inf\{c\T x:Ax\ge b\}$ 的对偶为 $\sup\{b\T\mu:A\T\mu=c,\mu\ge0\}$. 若任一最优值有限, 两问题都有最优解且最优值相同. 若原问题无界向下, 对偶不可行; 若对偶无界向上, 原问题不可行.
\end{theorem}
\begin{proof}
拉格朗日函数为 $(c-A\T\mu)\T x+b\T\mu$. 对自由 $x$ 取下确界, 当 $A\T\mu=c$ 时为 $b\T\mu$, 否则为 $-\infty$. 若原值有限, 多面体上的线性函数取到最小值, 记解为 $x^*$. 令 $J$ 为活跃行集合. 可行方向锥为 $D=\{d:a_j\T d\ge0,j\in J\}$, 任一这样的方向都允许足够小的可行正步长. 最优性使 $c\T d\ge0$ 对 $d\in D$. Farkas 引理使 $c=\sum_{j\in J}\mu_j a_j$, $\mu_j\ge0$. 非活跃项置零, 得对偶可行乘子且 $c\T x^*=\sum_j\mu_jb_j=b\T\mu$. 弱对偶使二者最优. 对偶值有限的情况可对对偶线性规划反向应用同一论证. 两个无界结论直接由弱对偶得到.
\end{proof}
对线性问题, 原始、对偶可行性加互补松弛就是完整最优性证书. 两边都不可行也是可能的, 因而不能把无界与不可行一一对应而忽略此情形.

\begin{example}
求 $\tfrac12x\T Qx+c\T x$ 在 $Ax\le b$ 下的对偶, 其中 $Q\succ0$.
\end{example}
\begin{solution}
对固定 $\mu\ge0$, 无约束二次函数的极小点为 $x(\mu)=-Q^{-1}(c+A\T\mu)$. 代回得
\[
q(\mu)=-\tfrac12\mu\T AQ^{-1}A\T\mu-\mu\T(b+AQ^{-1}c)-\tfrac12c\T Q^{-1}c.
\]
所以对偶等价于非负正交域上的半正定二次最小化. 若原问题可行, 正定二次项保证原解存在唯一, 仿射约束强对偶给最优乘子. 最优性证书是 $Ax^*\le b$, $\mu^*\ge0$, $x^*=-Q^{-1}(c+A\T\mu^*)$ 及 $\mu_j^*((Ax^*)_j-b_j)=0$.
\end{solution}

\originalsection{Fenchel 对偶及反例}
对闭真凸函数 $f_1,f_2$, 原问题为 $\inf_x[f_1(x)+f_2(x)]$. 将同一个变量拆成 $x_1=x_2$, 拉格朗日函数取 $f_1(x_1)+f_2(x_2)+y\T(x_2-x_1)$, 分别极小化得
\[
q(y)=-f_1^*(y)-f_2^*(-y).
\]
若原值有限且 $\ri\dom f_1\cap\ri\dom f_2\ne\varnothing$, 等式约束强对偶给出零间隙与对偶解. 原点 $x^*$ 与乘子 $y^*$ 的完整证书为
$x^*\in\argmin_x[f_1(x)-y^{*\mathsf T}x]$ 及 $x^*\in\argmin_x[f_2(x)+y^{*\mathsf T}x]$. 若两者可微且全空间有限, 这等价于 $y^*=\nabla f_1(x^*)=-\nabla f_2(x^*)$.

Fenchel 两侧结构对称. 对偶值有限且 $\ri\dom f_1^*\cap\ri(-\dom f_2^*)\ne\varnothing$ 时, 反向使用强对偶可保证原解. 这个结论与原侧资格条件所保证的“对偶解存在”要分开.

\begin{example}
说明强对偶、对偶解存在与原解存在不能互相替代.
\end{example}
\begin{solution}
第一例沿用部分极小化反例: 在非负象限最小化 $e^{-\sqrt{x_1x_2}}$, 加等式 $x_1=0$. 原值为1. 对偶 $q(\lambda)=\inf_{x\ge0}[e^{-\sqrt{x_1x_2}}+\lambda x_1]$ 在 $\lambda\ge0$ 时为0, 在 $\lambda<0$ 时为 $-\infty$. 当 $\lambda\ge0$, 选 $x_1\downarrow0$ 同时 $x_1x_2\to\infty$ 就逼近0; 当 $\lambda<0$, 固定 $x_2=0$ 令 $x_1\to\infty$ 即无界. 因而间隙为1. 非负象限的相对内部不满足 $x_1=0$, 资格条件失效.

第二例最小化 $x$ 且 $x^2\le0$. 唯一可行原解为0, 原值0. 对 $\mu>0$, 完成平方给 $q(\mu)=-1/(4\mu)$; $\mu=0$ 时为 $-\infty$. 对偶上确界0而无最优乘子. 扰动函数为 $p(u)=-\sqrt u$ 对 $u\ge0$, 其余为 $+\infty$, 在0下半连续却没有有限斜率支持线.

最后, 无约束最小化 $e^x$ 的原值0未取到. 把它写成 $f_1=e^x,f_2=0$, Fenchel 对偶可在唯一可行共轭乘子0处取到0. 因而即使强对偶与对偶解都有, 原解仍可不存在.
\end{solution}

\originalsection{极小极大与零和博弈}
对任意有限值 $\phi:X\times Z\to\R$, 总有 $\sup_z\inf_x\phi(x,z)\le\inf_x\sup_z\phi(x,z)$. 因为固定 $x,z$ 时, 左侧内层不超过 $\phi(x,z)$, 对另一变量依次取界即可. 鞍点 $(x^*,z^*)$ 定义为
$\phi(x^*,z)\le\phi(x^*,z^*)\le\phi(x,z^*)$ 对全部 $x,z$ 成立. 鞍点等价于极小极大等式成立且两侧最优值分别在 $x^*,z^*$ 取到: 一方面由鞍点夹逼; 另一方面两侧共同最优值夹住 $\phi(x^*,z^*)$, 于是两条鞍点不等式成立.

约束优化是 $\inf_x\sup_{\mu\ge0}L(x,\mu)$, 因为不可行点让内层为 $+\infty$, 可行点让它为 $f(x)$. 对偶是把两层交换. 因而拉格朗日鞍点就是零间隙且存在原始--对偶最优解的情况.

更一般地取 $p(u)=\inf_x\sup_z[\phi(x,z)-u\T z]$. 对固定 $x$, 内层是 $-\phi(x,\cdot)$ 的共轭在负参数处的值. 对 $u$ 再极小化并交换两个下确界, 得 $q(\mu)=\inf_x\widehat\phi(x,\mu)$, 其中 $\widehat\phi$ 是关于第二变量的闭凹包. 若每个第二变量函数已经闭凹, 对偶正是 $\sup_z\inf_x\phi(x,z)$. 当 $p$ 真凸且在0闭, MC/MC 给极小极大等式. 原资料这一构造的作用是说明何种闭凹化进入对偶, 而不是在任意 $X,Z$ 上无条件交换极值.

\begin{example}
有限零和矩阵博弈中, 玩家一以概率 $x\in\Delta_n$ 选行, 玩家二以概率 $z\in\Delta_m$ 选列, 第一人向第二人支付 $x\T Az$. 证明双方最坏情形最优值相同.
\end{example}
\begin{solution}
对固定 $x$, 最大期望支付为 $\max_j(A\T x)_j$. 玩家一问题因此是线性规划 $\min t$, $A\T x\le t\mathbf1$, $\mathbf1\T x=1$, $x\ge0$. 其对偶可写成 $\max s$, $Az\ge s\mathbf1$, $\mathbf1\T z=1$, $z\ge0$, 正是玩家二的 $\max_z\min_i(Az)_i$. 两个单纯形紧, 目标连续, 原值有限; 线性规划强对偶保证等式与双方策略存在. 例如 $A=\left(\begin{smallmatrix}1&-1\\-1&1\end{smallmatrix}\right)$ 时, 双方各以 $(1/2,1/2)$ 混合, 期望值恒为0, 直接构成鞍点.
\end{solution}

\originalchapter{次梯度与灵敏度}\label{chap:5}
本章对应原第12讲, 第158--172页. 可微凸函数由一个切平面提供下界; 不可微点可能允许多个支撑斜率. 次梯度将这种全局下界性质保留下来, 既提供最优性证书, 也为次梯度与切平面算法提供信息.

\originalsection{定义、存在性与共轭等号}
\begin{definition}
对真凸函数 $f$, 若 $x\in\dom f$ 且 $g\in\R^n$ 满足 $f(z)\ge f(x)+g\T(z-x)$ 对全部 $z\in\R^n$ 成立, 称 $g$ 为 $f$ 在 $x$ 的次梯度. 全部次梯度组成次微分 $\partial f(x)$; 域外约定为空集.
\end{definition}
次微分是闭凸集, 因为每个 $z$ 对 $g$ 规定一个闭半空间. 当 $f$ 可微时, $\partial f(x)=\{\nabla f(x)\}$: 切平面不等式给包含, 而任一次梯度对 $z=x+td$ 的不等式在 $t\downarrow0$ 和 $t\uparrow0$ 两向给出 $g\T d=\nabla f(x)\T d$.

\begin{example}
计算 $f(x)=|x|$ 及 $h(x)=\max\{0,(x^2-1)/2\}$ 的次微分.
\end{example}
\begin{solution}
绝对值在正负半轴的导数分别为1和 $-1$, 在0处次梯度条件是 $|z|\ge gz$ 对全部 $z$, 等价于 $-1\le g\le1$. 对 $h$, 当 $|x|<1$ 时局部恒为0且可微, 次微分为 $\{0\}$; 当 $|x|>1$ 时为 $\{x\}$. 在1处, 支撑斜率可取0与1之间任意值, 由左、右差商也只能取这些值, 所以为 $[0,1]$; 在 $-1$ 处为 $[-1,0]$.
\end{solution}
\begin{theorem}\label{thm:subgrad-exists}
真凸函数在 $\ri\dom f$ 上次微分非空. 若 $S$ 是有效域仿射包的方向空间, 则在这些点 $\partial f(x)=S^\perp+G$, 其中 $G\subseteq S$ 非空紧. 对任意域内点, 次微分非空紧当且仅当该点属于有效域的环境空间内部.
\end{theorem}
\begin{proof}
在仿射包内, 对平移函数 $p(u)=f(x+u)-f(x)$ 使用 MC/MC 的相对内部支持论证, 得到 $p(u)\ge g\T u$ 的非竖直支持平面, 从而次梯度存在. 也可直接在闭上图的边界 $(x,f(x))$ 支持; 相对内部连续性保证这个边界高度未改变, 域内相对球强制竖直系数为正.

取相对球半径 $r$ 并在略大球上用凸性给有限上界 $M$. 对 $g\in\partial f(x)\cap S$, 在 $z=x+r g/\norm g$ 代入支撑不等式, 得 $r\norm g\le M-f(x)$, 所以这一集合有界, 又闭非空, 故紧. 沿 $S^\perp$ 加任何向量不改变域内支撑不等式, 给出分解. 若点是环境内点, $S=\R^n$, 所以次微分非空紧.

反过来, 若域内点不是环境内点, 分离有效域得到非零法向量 $a$, $a\T(z-x)\le0$ 对全部域内 $z$. 任一次梯度 $g$ 都可加上 $ta$, $t\ge0$, 仍是次梯度, 因而若非空就无界. 这证明必要性, 包括低维域情形.
\end{proof}
边界点可以有次梯度, 也可以没有. 指示函数在边界通常具有无界法锥; $-\sqrt x$ 在非负半轴边界0则没有有限次梯度. 后续算法在每一步调用次梯度时必须保证所选点确实有可用次梯度.

\begin{theorem}[Fenchel 不等式与等号]\label{thm:fenchel-equality}
对真凸函数, $x\T y\le f(x)+f^*(y)$. 等号成立当且仅当 $y\in\partial f(x)$. 若 $f$ 闭, 还等价于 $x\in\partial f^*(y)$.
\end{theorem}
\begin{proof}
不等式正是共轭定义. 等号使 $x$ 达到 $\sup_z[y\T z-f(z)]$, 即 $f(z)\ge f(x)+y\T(z-x)$; 反向也由同式得到. 若 $f$ 闭真凸, 双共轭等于 $f$, 对共轭函数反向使用相同等号条件即可.
\end{proof}
\begin{corollary}
闭真凸函数的全空间极小点集等于 $\partial f^*(0)$. 若 $0\in\ri\dom f^*$, 极小点存在; 极小点集非空紧当且仅当 $0\in\operatorname{int}\dom f^*$.
\end{corollary}
\begin{proof}
$x$ 最优当且仅当 $0\in\partial f(x)$, 用 Fenchel 等号等价于 $x\in\partial f^*(0)$. 其余结论由次梯度存在性与紧性定理得到.
\end{proof}

\originalsection{法锥与支撑函数}
集合 $C$ 在 $x\in C$ 的法锥为 $N_C(x)=\{g:g\T(z-x)\le0\ \forall z\in C\}$. 由定义立即有 $\partial\ind C(x)=N_C(x)$. 内点的法锥为 $\{0\}$; 边界法向量朝向所有可行方向的反侧.

\begin{proposition}\label{prop:poly-normal}
对多面体 $C=\{x:a_i\T x\le b_i\}$, 在可行点 $x$ 有 $N_C(x)=\cone\{a_i:a_i\T x=b_i\}$.
\end{proposition}
\begin{proof}
活跃法向量的非负组合显然是法向量. 反向, 任一法向量 $g$ 对局部可行方向 $d$ 满足 $g\T d\le0$. 这些方向由活跃不等式 $a_i\T d\le0$ 刻画, 因为不活跃约束有严格余量, 对任何固定方向足够小的正步长仍满足它们. Farkas 引理使 $g$ 是活跃法向量的非负组合.
\end{proof}
\begin{proposition}
若 $K=\cl\conv X$ 非空, 则 $\partial\sigma_X(y)=\argmax_{x\in K}y\T x$. 对有限最大仿射函数 $f(x)=\max_i(a_i\T x+b_i)$, 令 $I(x)=\{i:a_i\T x+b_i=f(x)\}$, 则 $\partial f(x)=\conv\{a_i:i\in I(x)\}$.
\end{proposition}
\begin{proof}
$\sigma_X=\ind K^*$, 且 $\ind K$ 闭真凸, 所以 Fenchel 等号说明 $x\in\partial\sigma_X(y)$ 当且仅当 $x\in K$ 且 $y\T x=\sigma_X(y)$. 最大值未取到时, 两边都是空集.

对最大仿射函数, 每个活跃斜率都是次梯度, 其凸组合也都是. 若 $g$ 不在这些斜率凸包中, 严格分离给方向 $d$ 使 $g\T d>\max_{i\in I(x)}a_i\T d$. 因为项数有限, 不活跃项的余量确保足够小的 $t>0$ 时, $f(x+td)=f(x)+t\max_{i\in I(x)}a_i\T d$. 这违反 $g$ 的次梯度不等式, 所以没有更多次梯度.
\end{proof}

\originalsection{链式法则与和法则}
\begin{theorem}\label{thm:subgrad-chain}
令 $F(x)=f(Ax)$ 为真凸函数. 总有 $A\T\partial f(Ax)\subseteq\partial F(x)$. 若 $\operatorname{Range}A\cap\ri\dom f\ne\varnothing$, 或 $f$ 为多面体凸函数, 则等号成立.
\end{theorem}
\begin{proof}
正向包含把 $f$ 的支撑不等式代入 $Az$ 即得. 反向设 $d\in\partial F(x)$, 则 $(y,z)=(Ax,x)$ 达到问题 $\min[f(y)-d\T z]$, 约束 $Az=y$, $y\in\dom f$ 的有限最小值. 相对内部资格条件使这个等式约束问题具有最优乘子 $v$. 若 $f$ 多面体, 可以在其多面体上图上写成线性规划, 也有最优乘子. 拉格朗日函数为 $f(y)-d\T z+v\T(Az-y)$. 对自由 $z$ 极小化要求 $A\T v=d$, 而对 $y$ 的最优性说明 $v\in\partial f(Ax)$. 因而 $d$ 属于右侧所述像.
\end{proof}
多面体函数指上图为多面体的真凸函数, 可写成多面体域上有限个仿射函数的最大值. 这类函数的资格条件可以比一般凸函数宽松, 因为其有限生成结构排除了参数向无穷逃逸而丢失边界的困难.

\begin{theorem}\label{thm:subgrad-sum}
设 $f_i$ 真凸且 $F=\sum_i f_i$ 真. 若 $\bigcap_i\ri\dom f_i\ne\varnothing$, 则 $\partial F(x)=\sum_i\partial f_i(x)$. 若前 $k$ 个函数多面体, 条件可放宽为
\[
\left(\bigcap_{i\le k}\dom f_i\right)\cap
\left(\bigcap_{i>k}\ri\dom f_i\right)\ne\varnothing.
\]
\end{theorem}
\begin{proof}
令 $A x=(x,\ldots,x)$, $h(x_1,\ldots,x_m)=\sum_i f_i(x_i)$. 对各独立变量, 次梯度不等式说明 $\partial h=\partial f_1\times\cdots\times\partial f_m$. 又 $F=h\circ A$, 用链式法则即得一般资格条件下的等式. 在含多面体项的情况下, 同一拆变量证明中的相等约束与多面体域可合并为一个多面体集合; 对其余项使用 MC/MC 多面体加强版本. 它只要求该多面体与剩余有效域相对内部相交, 正是所列放宽条件. 最优乘子的存在再给出各分量次梯度的分解.
\end{proof}
没有资格条件时只能保证次微分和包含于总函数的次微分, 不能任意作反向分解. 例如 $f(x)=-\sqrt x$ 在非负半轴上, $C=\{0\}$, 总函数 $f+\ind C$ 在0最小且有零次梯度, 但 $\partial f(0)=\varnothing$, 所以无法把零分成目标次梯度与法向量之和.

\begin{theorem}[约束次梯度最优性]
若 $f$ 与凸集 $X$ 满足下列任一条件: 相对内部相交; $f$ 多面体且 $\dom f\cap\ri X\ne\varnothing$; $X$ 多面体且 $\ri\dom f\cap X\ne\varnothing$; 二者都多面体且有共同可行点, 则 $x^*$ 最优当且仅当存在 $g\in\partial f(x^*)$ 使 $-g\in N_X(x^*)$, 即 $g\T(x-x^*)\ge0$ 对所有 $x\in X$.
\end{theorem}
\begin{proof}
在 $X$ 上最小化等价于全空间最小化 $f+\ind X$. 全空间最优性等价于 $0\in\partial(f+\ind X)(x^*)$. 上述条件分别允许使用和法则, 所以它等价于 $0\in\partial f(x^*)+N_X(x^*)$, 即所述条件.
\end{proof}
可微情况恢复第1章的变分不等式. 一般情况只须找到一个适当次梯度, 不要求整个次微分都朝向可行集外.

\originalsection{乘子的灵敏度意义}
对约束扰动函数 $p(u)$, 若它真凸且强对偶成立, 则最优乘子集为 $-\partial p(0)$. 因为 $q(\mu)=-p^*(-\mu)$, 最优性正是 $p(0)+p^*(-\mu)=0$, 由 Fenchel 等号即得. 因而每个最优乘子给出全局扰动下界
\[
p(u)\ge p(0)-\mu^{*\mathsf T}u.
\]
放宽约束的分量 $u_j>0$ 不会提高最优值, 非负乘子记录这一降低的局部边际幅度. 若 $p$ 在0可微, 最优乘子唯一, 且 $\mu_j^*=-\partial p(0)/\partial u_j$. 若不可微, 可能有多个乘子, 每个都给一个合法的一阶下界, 不能假定某个乘子在大幅扰动下仍精确预测最优值.

\begin{example}
最小化 $x^2/2$ 且 $x\ge b$, 在 $b>0$ 处求影子价格.
\end{example}
\begin{solution}
写约束 $b-x\le u$. 对足够小的 $u$, 最优点为 $x=b-u$, 所以 $p(u)=(b-u)^2/2$. 有 $-p'(0)=b$. 拉格朗日函数 $x^2/2+\mu(b-x)$ 在 $x=\mu$ 极小, 原最优点为 $b$, 因而最优乘子为 $b$, 与扰动导数一致. 当放宽量 $u\ge b$ 时, 最优点变为0, 旧的一阶公式只能给下界, 不能继续当作等式.
\end{solution}

\originalchapter{锥规划与问题结构}\label{chap:6}
本章承接原第11讲锥对偶及第13讲, 原PDF第153--157页、第173--184页. 后续算法的选择不仅取决于是否凸, 还取决于目标与约束如何分解、线性优化或投影是否容易计算, 以及约束锥的结构.

\originalsection{锥约束与对偶锥}
闭凸锥 $C$ 的非负对偶锥定义为 $C^+=\{y:y\T x\ge0\ \forall x\in C\}$. 它等于 $-C^\circ$, 且 $(C^+)^+=C$. 最小化闭真凸 $f$ 于 $x\in C$ 可写成 $\min(f+\ind C)$, Fenchel 对偶为
\[
\sup_{y\in C^+}\{-f^*(y)\}.
\]
原讲义有时将它改写为 $\min_{y\in C^+}f^*(y)$, 此时是将对偶目标整体换号, 两个最小化问题的值不能直接写成相等.

\begin{theorem}
若原值有限且 $\ri\dom f\cap\ri C\ne\varnothing$, 则上述锥对偶零间隙且有最优解. 若对偶值有限且 $\ri\dom f^*\cap\ri C^+\ne\varnothing$, 则反向保证零间隙及原解存在. 原始--对偶解对满足 $x\in C$, $y\in C^+$, $y\in\partial f(x)$ 与 $x\T y=0$.
\end{theorem}
\begin{proof}
前两条分别在 Fenchel 对偶的原侧与共轭侧使用资格条件. 最优解对的拉格朗日条件是 $y\in\partial f(x)$ 与 $-y\in\partial\ind C(x)=N_C(x)$. 对锥, 后式等价于 $y\in C^+$ 与 $x\T y=0$: 在法锥定义中取 $z=0,2x$ 得内积为0, 再对任意 $z\in C$ 得非负对偶条件. 反向把这些条件代回 Fenchel 等号, 得 $f(x)=-f^*(y)$, 弱对偶使二者最优.
\end{proof}

\begin{proposition}\label{prop:conic-dual}
线性锥规划有如下两种常用对偶形式:
\[
\inf_{Ax=b,\ x\in C}c\T x
\quad\longleftrightarrow\quad
\sup_{c-A\T\lambda\in C^+}b\T\lambda,
\]
\[
\inf_{Ax-b\in C}c\T x
\quad\longleftrightarrow\quad
\sup_{A\T y=c,\ y\in C^+}b\T y.
\]
\end{proposition}
\begin{proof}
第一式拉格朗日函数取 $c\T x+\lambda\T(b-Ax)$. 在锥上极小化 $(c-A\T\lambda)\T x$ 时, 系数属于 $C^+$ 则下确界为0, 否则沿某个锥方向趋于 $-\infty$. 第二式先令 $z=Ax-b\in C$, 用乘子 $y$ 对等式写 $c\T x-y\T(Ax-b-z)$; 对自由 $x$ 极小化要求 $A\T y=c$, 对锥变量 $z$ 极小化要求 $y\in C^+$, 剩余为 $b\T y$.
\end{proof}
若用仿射集合 $b+S$ 而非矩阵等式描述, $f(x)=c\T x+\ind{b+S}(x)$ 的共轭为 $(y-c)\T b$ 当 $y-c\in S^\perp$, 否则为 $+\infty$. 对偶最大值因此为 $\sup[b\T c-b\T y]$, 约束 $y-c\in S^\perp$, $y\in C^+$. 去掉常数并换号才得到原讲义的对偶最小化形式 $\min b\T y$. 这一计算保留了常数, 有助于检查对偶最优值的符号.

\originalsection{二阶锥与半正定锥}
\begin{proposition}
非负正交锥、二阶锥 $\mathcal Q_n=\{(u,t):\norm u\le t\}$ 以及对称半正定矩阵锥 $\mathbb S_+^n$ 都自对偶, 即 $C^+=C$. 矩阵空间的内积取 $\langle X,Y\rangle=\tr(XY)$.
\end{proposition}
\begin{proof}
非负正交域的结论逐坐标取单位向量即可. 若 $(u,t),(v,s)\in\mathcal Q_n$, 则 $u\T v+ts\ge-\norm u\norm v+ts\ge0$, 因而二阶锥包含于其对偶. 若 $(v,s)$ 不在锥内, 当 $s<0$ 取 $(0,1)$ 给负内积; 否则 $\norm v>s$, 取 $(-v/\norm v,1)$ 给负内积, 得反向包含.

若 $X,Y\succeq0$, 对 $X$ 作谱分解 $X=\sum_i\lambda_i v_iv_i\T$, 则 $\tr(XY)=\sum_i\lambda_i v_i\T Yv_i\ge0$. 若对称矩阵 $Y$ 非半正定, 取 $v$ 使 $v\T Yv<0$, 秩一半正定矩阵 $vv\T$ 与它的内积为负, 所以不属于对偶. 对称矩阵空间实际维数为 $n(n+1)/2$, 不必把重复的对称分量当作独立坐标.
\end{proof}
二阶锥规划的约束可写为 $A_ix-b_i\in\mathcal Q_{n_i}$, 目标线性. 把锥取笛卡尔积并用命题\ref{prop:conic-dual}, 得对偶 $\max\sum_i b_i\T y_i$, $\sum_iA_i\T y_i=c$, $y_i\in\mathcal Q_{n_i}$. 若有严格锥可行点且原值有限, 则零间隙与对偶解存在. 这些资格条件并不比一般线性锥规划更强, 也不能因锥自对偶就删掉.

\begin{example}
把 $\norm{Bx-d}\le a\T x+\beta$ 和 $\norm{x}^2\le t$ 写成锥约束.
\end{example}
\begin{solution}
第一条就是 $(Bx-d,a\T x+\beta)\in\mathcal Q$, 自动要求右侧非负. 第二条等价于 $(2x,t-1,t+1)\in\mathcal Q_{n+2}$, 因为 $(t+1)^2-[(t-1)^2+4\norm{x}^2]=4(t-\norm{x}^2)$, 且锥条件保证 $t+1\ge0$. 反向原不等式使 $t\ge0$, 故也满足锥约束. 不少凸二次模型可通过这种提升成为二阶锥规划.
\end{solution}
半正定规划在 $\mathbb S^n$ 中最小化 $\langle C,X\rangle$, 约束 $\mathcal A(X)=b$, $X\succeq0$. 其对偶为 $\max b\T y$, $C-\mathcal A^*(y)\succeq0$, 其中伴随算子由 $\langle\mathcal A^*y,X\rangle=y\T\mathcal A(X)$ 定义. 互补性是 $\langle X,C-\mathcal A^*y\rangle=0$. 若原始与对偶松弛矩阵都半正定, 该内积等于0还蕴含其矩阵乘积为零: 从 $\tr(X^{1/2}SX^{1/2})=0$ 得半正定矩阵 $X^{1/2}SX^{1/2}=0$, 继而 $S^{1/2}X^{1/2}=0$, 因而 $SX=0$. 后章会给最大特征值与离散二次规划的具体建模例子.

\originalsection{可分离结构与离散下界}
设变量按块 $x_i\in X_i$, 目标为 $\sum_i f_i(x_i)$, 耦合约束为 $\sum_i g_{ji}(x_i)\le0$. 对固定 $\mu\ge0$, 拉格朗日下确界可分解为
\[
q(\mu)=\sum_iq_i(\mu),\qquad
q_i(\mu)=\inf_{x_i\in X_i}\left[f_i(x_i)+\sum_j\mu_jg_{ji}(x_i)\right].
\]
每个块可独立计算; 原问题变量维数很大时, 耦合乘子维数可能较小. 若各块下确界取到, 令块解为 $x_i(\mu)$, 则向量 $s_j=\sum_i g_{ji}(x_i(\mu))$ 是凹函数 $q$ 的超梯度, 即 $q(\nu)\le q(\mu)+s\T(\nu-\mu)$. 证明只须在 $q(\nu)$ 的极小化中代入旧块解, 无须对块解关于乘子求导. 等价地 $-s$ 是凸函数 $-q$ 的次梯度.

当 $X_i$ 离散, 强对偶可能失败, 但 $q(\mu)$ 仍是原最优值的合法下界. 在分支定界中, 每个节点固定部分离散选择, 剩余子问题的对偶值给节点下界; 已找到的可行解给全局上界. 若节点下界不小于当前上界, 节点可剪枝. 对偶间隙影响下界紧度与搜索效率, 不影响下界的正确性. 本课程在这里介绍结构用途, 并未把离散问题变成无条件强对偶的凸问题.

\originalsection{大规模和与多约束}
统计学习常出现 $f(x)=\sum_{i=1}^m f_i(x)$, 每一项对应一个样本损失. 最小二乘为 $\sum_i(a_i\T x-b_i)^2$; 加 $\ell_1$ 正则得到
$\sum_i(a_i\T x-b_i)^2+\lambda\sum_j|x_j|$. 不可微项促使部分坐标精确为0, 后面的近端软阈值公式会把这一性质计算出来. 离散概率空间上的期望 $\mathbb E F(x,\omega)=\sum_i p_iF(x,\omega_i)$ 也具有同一大和结构. 两阶段随机规划还可写为 $F_1(x)+\mathbb E_\omega[\inf_yF_2(x,y,\omega)]$, 其中内层值函数的凸性与闭性需要使用部分极小化条件, 不能因为写成期望就自动成立.

当约束数很大时, 一种方式是罚函数, 如 $\sum_j[(a_j\T x-b_j)_+]^2$ 或 $\sum_j(a_j\T x-b_j)_+$. 二次惩罚通常需要参数趋大才逼近精确可行; 非光滑线性惩罚在乘子有界条件下可以精确, 其阈值将在下降方法章节证明. 另一种方式是先保留少量约束, 求解后检测违反项并加回, 形成外逼近; 若维护一组确定可行的点或内模型, 则是内逼近. 它们分别提供不同方向的界, 后面的切平面与单纯形分解将说明如何形成最优性间隔.

原第178页给出实践中的粗略问题层级: 线性与凸二次、二阶锥、半正定、一般凸、非凸、离散. 这是一种课程组织视角, 并非对每个实例的严格复杂度排名. 网络结构、可分离性、投影和线性优化子问题的可解性, 都可能显著改变实际难度. 第1讲所强调的算法与对偶结合, 在这些结构化问题中最有作用.

\originalsection{由紧性保证强对偶}
\begin{proposition}
设 $X$ 闭凸, $f,g_j$ 为闭凸函数, 原可行集 $C=\{x\in X:g(x)\le0\}$ 非空紧且存在有限目标值可行点, $f$ 真. 则原解集非空紧, 拉格朗日对偶零间隙. 此条件本身不保证对偶解存在.
\end{proposition}
\begin{proof}
在紧可行集上用下半连续性得到原解存在, 且最优值有限. 令 $F(x,u)=f(x)$ 当 $x\in X,g(x)\le u$, 其他点为 $+\infty$. 其上图由闭凸条件定义, 所以闭凸. 在 $u=0$ 处的一个非空有限目标水平集包含于紧可行集, 因而紧. 部分极小化定理使 $p(u)=\inf_xF(x,u)$ 闭真凸. 于是 $p^{**}(0)=p(0)$, 即强对偶. 对偶解需要 $p$ 在0有次梯度, 紧性只保证闭性而不保证这一点; 前章 $x^2\le0$ 的例子正说明区别.
\end{proof}

\originalchapter{锥模型、惩罚与次梯度迭代}\label{chap:7}
本章对应原第14--16讲, 原PDF第185--213页. 第14讲先承接锥对偶的建模应用, 随后转入非光滑算法. 本章补写的计算与收敛证明是对幻灯片省略步骤的教学重构.

\originalsection{二阶锥与半正定模型}
线性锥问题的两种常用格式是
\[
 \inf\{c\T x:Ax=b,\ x\in C\},\qquad
 \inf\{c\T x:Ax-b\in C\}.
\]
这里$C$是闭凸锥, 其对偶锥取$C^*=\{y:y\T z\ge0,\ z\in C\}$. 第一种格式的对偶是$\sup\{b\T\lambda:c-A\T\lambda\in C^*\}$; 第二种格式的对偶是$\sup\{b\T\lambda:A\T\lambda=c,\ \lambda\in C^*\}$. 例如对第二式取拉格朗日函数$c\T x-\lambda\T(Ax-b)$, 对$x$取下确界, 即可看到等式$A\T\lambda=c$的来源. 本节所说的对偶界总是弱对偶界; 把它升级为最优值相等还需前章的强对偶条件.

记$\mathcal Q_d=\{(u,t)\in\R^{d-1}\times\R:\norm{u}\le t\}$为二阶锥. 因$\mathcal Q_d^*=\mathcal Q_d$, 问题
\[
 \min c\T x\quad\text{满足}\quad A_i x-b_i\in\mathcal Q_{d_i},\quad i=1,\ldots,m
\]
的对偶是
\[
 \max\sum_{i=1}^m b_i\T\lambda_i\quad\text{满足}\quad
 \sum_{i=1}^m A_i\T\lambda_i=c,\quad \lambda_i\in\mathcal Q_{d_i}.
\]
多个锥约束只是在乘积锥上写一个约束, 并不需要另造一套对偶理论.

\begin{example}
在鲁棒线性规划中, 第$j$条约束要求$a\T x\le b$对所有$(a,b)\in T_j$成立. 设不确定集为椭球
\[
 T_j=\{(\bar a_j+P_j u,\bar b_j+q_j\T u):\norm{u}\le1\}.
\]
把这条无限多不等式写成一个二阶锥约束.
\end{example}
\begin{solution}
最坏约束违反量为
\begin{align*}
 g_j(x)&=\sup_{\norm{u}\le1}\bigl[(\bar a_j+P_j u)\T x-\bar b_j-q_j\T u\bigr]\\
 &=\bar a_j\T x-\bar b_j+\sup_{\norm{u}\le1}(P_j\T x-q_j)\T u\\
 &=\bar a_j\T x-\bar b_j+\norm{P_j\T x-q_j}.
\end{align*}
最后一步用柯西不等式; 向量非零时取同方向的单位向量达到上界. 因而$g_j(x)\le0$恰好等价于$(P_j\T x-q_j,\bar b_j-\bar a_j\T x)\in\mathcal Q_{d_j+1}$. 这说明椭球不确定性所引入的非线性仍可保留为标准锥模型.
\end{solution}

对称矩阵空间的内积是$\inner{U}{V}=\tr(UV)$. 半正定锥$\mathbb S_+^n$在此内积下自对偶, 内部为正定矩阵. 于是半正定规划
\[
 \min\inner{D}{X}\quad\text{满足}\quad
 \inner{A_i}{X}=b_i\ (i=1,\ldots,m),\quad X\succeq0
\]
的对偶是$\max\{b\T\lambda:D-\sum_i\lambda_i A_i\succeq0\}$. 若原问题有正定可行矩阵, 且原最优值有限, 则强对偶成立且对偶最优解存在. 若对偶有严格可行点, 且对偶最优值有限, 则反向的结论成立. 有限性不可从严格可行性自动推出.

\begin{example}
给定仿射对称矩阵$M(\lambda)=D+\sum_{i=1}^m\lambda_i M_i$, 最小化其最大特征值.
\end{example}
\begin{solution}
由谱定理, $\lambda_{\max}(M(\lambda))\le z$等价于$zI-M(\lambda)\succeq0$. 故问题可写为
\[
 \min_{z,\lambda}z\quad\text{满足}\quad zI-D-\sum_i\lambda_i M_i\succeq0.
\]
变量虽包含矩阵的特征值, 约束对$(z,\lambda)$却是仿射半正定约束.
\end{solution}

\begin{example}
考虑非凸二次等式问题
\[
 \min_x x\T Q_0x+a_0\T x+b_0\quad\text{满足}\quad
 x\T Q_i x+a_i\T x+b_i=0,\quad i=1,\ldots,m,
\]
其中$Q_i$对称. 推导一个半正定规划形式的拉格朗日下界.
\end{example}
\begin{solution}
令$Q(\lambda)=Q_0+\sum_i\lambda_iQ_i$, $a(\lambda)=a_0+\sum_i\lambda_i a_i$, $b(\lambda)=b_0+\sum_i\lambda_i b_i$. 对偶函数是
\[
 q(\lambda)=\inf_x\{x\T Q(\lambda)x+a(\lambda)\T x+b(\lambda)\}.
\]
条件$q(\lambda)\ge t$是一个对所有$x$成立的二次不等式. 它等价于
\[
 \begin{pmatrix}
 Q(\lambda)&a(\lambda)/2\\
 a(\lambda)\T/2&b(\lambda)-t
 \end{pmatrix}\succeq0.
\]
正定方向的等价性可由配方看出. 一般情形, 矩阵半正定显然推出二次不等式. 反向, 对$(v,s)$在$s\ne0$时把二次型写成$s^2$乘以$x=v/s$的表达式; 在$s=0$时由连续性取极限. 因而对偶就是在上述矩阵约束下最大化$t$. 它给原问题下界, 不保证非凸原问题的强对偶. 特别地, $x_i\in\{0,1\}$可写成$x_i^2-x_i=0$, 所以这种下界也适用于离散优化.
\end{solution}

\originalsection{精确惩罚与乘子的阈值}
考虑凸约束问题
\[
 f^*=\inf\{f(x):x\in X,\ g_j(x)\le0,\ j=1,\ldots,r\},
\]
其中$X$凸, $f,g_j$为实值凸函数. 惩罚函数$P:\R^r\to\R$在负正交锥上为零, 在出现正分量时为正; 为使$P\circ g$仍凸, 本节还取$P$凸且逐坐标不减. 常用选择是$P_c(u)=c\sum_j\max\{0,u_j\}$或$\frac c2\sum_j\max\{0,u_j\}^2$.

用扰动值函数$p(u)=\inf\{f(x):x\in X,\ g(x)\le u\}$, 有
\begin{equation}\label{eq:penalty-value}
 \inf_{x\in X}\{f(x)+P(g(x))\}=\inf_u\{p(u)+P(u)\}.
\end{equation}
一侧由取$u=g(x)$得到; 另一侧由$g(x)\le u$与$P$单调性得到, 再分别逼近下确界. 设$p$真凸, 并满足Fenchel强对偶的相对内部条件, 则右侧等于
\[
 \sup_{\mu\ge0}\{q(\mu)-P^*(\mu)\},\qquad
 q(\mu)=\inf_{x\in X}\{f(x)+\mu\T g(x)\}.
\]
$P$在负正交锥上为零, 故其共轭在含负分量的$\mu$处为$+\infty$. 惩罚的作用是在原对偶目标$q$上再扣掉一个凸函数$P^*$.

\begin{proposition}\label{prop:exact-penalty}
若$\mu^*\ge0$为一个满足$q(\mu^*)=f^*\in\R$的最优乘子, 且$\mu^*\in\partial P(0)$, 则惩罚问题的最优值等于$f^*$. 对$P_c(u)=c\sum_j\max\{0,u_j\}$, 条件为$0\le\mu_j^*\le c$. 若$c>\norm{\mu^*}_\infty$, 则惩罚问题的每个最优解都可行, 因而也是原问题的最优解.
\end{proposition}
\begin{proof}
次梯度不等式给$P(u)\ge\mu^{*\mathsf T}u$. 所以对每个$x\in X$,
\[
 f(x)+P(g(x))\ge f(x)+\mu^{*\mathsf T}g(x)\ge q(\mu^*)=f^*.
\]
另一方面, 用原问题的可行点逼近$f^*$, 惩罚项恒为零, 得到反向值界. 对分段线性惩罚, $\partial P_c(0)=[0,c]^r$. 若$x$有违反约束的分量, 则
\[
 c\sum_j(g_j(x))_+-\mu^{*\mathsf T}g(x)
 \ge\sum_{g_j(x)>0}(c-\mu_j^*)g_j(x)>0.
\]
此时惩罚目标严格大于$f^*$, 不可能最优. 这里严格阈值用于排除不可行最优解; 仅用非严格阈值足以保证最优值相等.
\end{proof}

一维共轭计算尤其直观. 对$c>0$,
\[
 (c\,u_+)^*(\mu)=
 \begin{cases}0,&0\le\mu\le c,\\+\infty,&\text{其他},\end{cases}
 \qquad
 \left(\frac c2(u_+)^2\right)^*(\mu)=
 \begin{cases}\mu^2/(2c),&\mu\ge0,\\+\infty,&\mu<0.\end{cases}
\]
第一式在一段乘子区间上完全平坦; 第二式只在零处取零. 原第193页还画了带线性项的二次惩罚. 为符合负半轴上惩罚为零的定义, 应取$P(u)=a u_++(u_+)^2$, $a\ge0$, 此时
\[
 P^*(\mu)=\begin{cases}
 +\infty,&\mu<0,\\0,&0\le\mu\le a,\\(\mu-a)^2/4,&\mu>a.
 \end{cases}
\]
这是在$u\ge0$上最大化$(\mu-a)u-u^2$所得. 原页印作$\max\{0,au+u^2\}$, 若按全实轴解释会在$u<-a$处产生正惩罚, 与该页图形和前页定义不相容; 这里明确采用与惩罚定义相容的分段形式. 更一般地, $P^*$在$\mu^*$取最小值等价于$0\in\partial P^*(\mu^*)$, 也等价于$\mu^*\in\partial P(0)$. 若$P$在零处可微且负半轴上为零, 则$\partial P(0)=\{0\}$, 无法容纳非零最优乘子. 这解释了精确惩罚常常必然非光滑.

\begin{example}
最小化$-x$, 约束$x\le0$. 比较$c x_+$和$\frac c2(x_+)^2$两种惩罚.
\end{example}
\begin{solution}
原最优解为$x=0$, 最优乘子为$\mu^*=1$. 线性惩罚在$c>1$时唯一最优于零; 在$c=1$时所有$x\ge0$都有目标值零, 因此值的精确性不代表解的精确性. 二次惩罚在$x>0$上目标为$-x+cx^2/2$, 导数为$-1+cx$, 最优解$x=1/c$始终违反约束, 目标值为$-1/(2c)$. 只有令$c\to\infty$才消除违反量.
\end{solution}

\originalsection{方向导数与最陡下降}
对凸函数$f$及$x\in\dom f$, 方向导数取单侧极限
\[
 f'(x;d)=\lim_{\alpha\downarrow0}\frac{f(x+\alpha d)-f(x)}{\alpha}.
\]
若$0<a<b$, 凸性给$f(x+ad)\le(1-a/b)f(x)+(a/b)f(x+bd)$, 故差商随$\alpha$增大不减, 上述极限存在于扩展实数中. 在$x\in\ri(\dom f)$时, 方向导数是次微分的支撑函数: $f'(x;d)=\sup_{g\in\partial f(x)}g\T d$. 在边界或域外方向上使用此公式必须另核对适用条件.

对本节的实值凸函数, $\partial f(x)$为非空紧凸集. $f'(x;d)<0$意味着沿$d$走足够小的正步长确实降低目标. 与之相反, 任意一个负次梯度方向不一定下降, 因为方向导数取所有次梯度的最大内积.

\begin{proposition}
设$g^*$为$\partial f(x)$中欧氏范数最小的向量. 若$g^*\ne0$, 单位球上的最陡下降方向是$d^*=-g^*/\norm{g^*}$, 并且$\min_{\norm{d}\le1}f'(x;d)=-\norm{g^*}$. 若$g^*=0$, 则$x$是全局最优点.
\end{proposition}
\begin{proof}
投影最优性给$(g-g^*)\T g^*\ge0$对所有$g\in\partial f(x)$成立. 因此$g\T d^*\le-\norm{g^*}$, 而$g=g^*$达到等号. 对任意单位球内$d$, 支撑函数至少为$g^{*\mathsf T}d\ge-\norm{g^*}$. 两个方向的界吻合. 若$0\in\partial f(x)$, 次梯度不等式给$f(z)\ge f(x)$对所有$z$成立.
\end{proof}

算法中常把未归一化的$-g^*$写作搜索方向, 再用步长吸收范数. 原第196页的表述应在这个意义下理解; 它不是单位球约束问题在$\norm{g^*}\ne1$时的原样最优解. 计算最小范数次梯度需要整个次微分, 代价可能很高. 更重要的是, 非光滑点附近活跃次梯度集合可突变; 仅依当前最陡方向和精确线搜索, 不能据此推出全局收敛. 原第197--198页的示意例用于强调这个失效机制, 并未给出可复算的函数表达式. 后面改用到最优解的距离作为进展量, 不要求每一步目标下降.

\originalsection{一个次梯度的计算}
若$f(x)=\sup_{z\in Z}\phi(x,z)$, 每个$\phi(\cdot,z)$凸, 且$z_x$在$x$达到上确界, 则任取$g\in\partial_x\phi(x,z_x)$都有$g\in\partial f(x)$. 证明只需把三件事连起来:
\[
 f(y)\ge\phi(y,z_x)\ge\phi(x,z_x)+g\T(y-x)=f(x)+g\T(y-x).
\]
因此不必算出整个次微分. 对有限个可微函数的最大值, 可选任何一个活跃分支的梯度.

拉格朗日对偶也是这种结构. 若$x_\mu$达到$q(\mu)=\inf_{x\in X}\{f(x)+\mu\T g(x)\}$, 则
\[
 q(\nu)\le f(x_\mu)+\nu\T g(x_\mu)
 =q(\mu)+g(x_\mu)\T(\nu-\mu).
\]
于是$g(x_\mu)$是凹函数$q$的超梯度, $-g(x_\mu)$是$-q$的次梯度. 在$\mu\ge0$上最大化$q$对应迭代$\mu_{k+1}=[\mu_k+\alpha_k g(x_{\mu_k})]_+$: 违反约束使相应乘子增大, 是计算公式与经济解释相吻合的地方.

\originalsection{投影次梯度法的距离估计}
设$X$非空闭凸, $f$凸且在迭代点有次梯度. 投影次梯度迭代为
\begin{equation}\label{eq:subgradient-iteration}
 x_{k+1}=P_X(x_k-\alpha_k g_k),\qquad g_k\in\partial f(x_k),\quad\alpha_k>0.
\end{equation}
投影保证可行性, 但不保证目标序列单调. 例如$f(x)=|x|$在零点允许次梯度$g=1$, 一步可把最优点移到$-\alpha$.

先回顾投影的非扩张性. 设$p=P_X(u)$, $q=P_X(v)$. 投影不等式给$(u-p)\T(q-p)\le0$及$(v-q)\T(p-q)\le0$. 相加得$\norm{p-q}^2\le(p-q)\T(u-v)\le\norm{p-q}\norm{u-v}$, 故$\norm{p-q}\le\norm{u-v}$.

\begin{lemma}\label{lem:subgradient-distance}
对迭代\eqref{eq:subgradient-iteration}, 任意$y\in X$都有
\[
 \norm{x_{k+1}-y}^2\le\norm{x_k-y}^2
 -2\alpha_k\bigl(f(x_k)-f(y)\bigr)+\alpha_k^2\norm{g_k}^2.
\]
若$f(x_k)>f(y)$且$0<\alpha_k<2(f(x_k)-f(y))/\norm{g_k}^2$, 则到$y$的距离严格减小.
\end{lemma}
\begin{proof}
用$P_X(y)=y$与非扩张性, 再展开平方, 得
\begin{align*}
 \norm{x_{k+1}-y}^2&\le\norm{x_k-\alpha_k g_k-y}^2\\
 &=\norm{x_k-y}^2-2\alpha_k g_k\T(x_k-y)+\alpha_k^2\norm{g_k}^2.
\end{align*}
次梯度不等式给$g_k\T(x_k-y)\ge f(x_k)-f(y)$, 代入即得. 严格步长条件使右边的修正量为负.
\end{proof}

\begin{theorem}\label{thm:subgradient-value}
设生成的次梯度满足$\norm{g_k}\le C$.
\begin{enumerate}
\item 若$\alpha_k=\alpha>0$, 则当$f^*=\inf_X f> -\infty$时, $\liminf_k f(x_k)\le f^*+\alpha C^2/2$. 若$f^*=-\infty$, 则$\inf_k f(x_k)=-\infty$.
\item 若$\alpha_k\to0$且$\sum_k\alpha_k=\infty$, 则$\liminf_k f(x_k)=f^*$, 因而历史最好值$\min_{0\le j\le k}f(x_j)$趋于$f^*$.
\end{enumerate}
\end{theorem}
\begin{proof}
固定$y\in X$. 若从某一步开始总有$f(x_k)-f(y)>\alpha C^2/2+\delta$, 距离引理使平方距离每步至少减少$2\alpha\delta$, 不可能永远非负. 故$\liminf f(x_k)\le f(y)+\alpha C^2/2$. 对$y$取下确界得到第一项; 当$f^*=-\infty$时可令$f(y)$任意小.

第二项若失败, 可选$y\in X$与$\delta>0$使从某步起$f(x_k)-f(y)\ge\delta$. 再取足够大的$k$使$\alpha_k C^2\le\delta$. 距离引理于是给$\norm{x_{k+1}-y}^2\le\norm{x_k-y}^2-\alpha_k\delta$. 求和与$\sum\alpha_k=\infty$矛盾. 下界$f(x_k)\ge f^*$完成证明. 这个论证不要求最优解存在.
\end{proof}

如果还要求迭代点本身收敛, 可增加$f$下半连续、最优解存在以及$\sum_k\alpha_k^2<\infty$. 距离引理此时给到每个最优点的平方距离仅有可求和的向上误差, 所以该距离有极限, 且序列有界. 由上一定理取得目标值趋于$f^*$的子列, 再取其收敛子列; 下半连续性保证极限最优. 到这个极限点的距离已有极限且子列距离趋零, 故整个序列趋于该最优点. 必须把这个点收敛结论与仅有最好值收敛的较弱结论区分开.

\originalsection{步长、目标值估计与误差次梯度}
常数步长带来固定误差邻域, 递减步长以渐慢的移动换取精度. 若知道$f^*$, 可选动态步长
\[
 \alpha_k=\frac{f(x_k)-f^*}{\norm{g_k}^2}
\]
并在$g_k=0$时停止. 距离引理给平方距离至少下降$(f(x_k)-f^*)^2/\norm{g_k}^2$. 用统一界$C^2$替代分母也可保证下降. 实际中常以目标值估计$f_k$替代$f^*$; 估计过低会造成围绕最优集合的振荡, 过高则可能只到达$f\le f_k$的水平集. 原讲义给的自适应办法是$f_k=\min_{j\le k}f(x_j)-\delta_k$, 以是否达到本次目标调整期望下降量$\delta_k$: 达到时放大, 未达到时缩小但保留一个正下限. 这是一类实现策略, 不能在未给参数条件时把它当作已证明的精确收敛定理.

\begin{definition}
对$\varepsilon\ge0$, 向量$g$称为$f$在$x\in\dom f$的$\varepsilon$次梯度, 若对所有$z$都有$f(z)\ge f(x)+g\T(z-x)-\varepsilon$. 全体这样的向量组成$\partial_\varepsilon f(x)$; 域外该集合为空.
\end{definition}
有$\partial_0 f(x)=\partial f(x)$, $\varepsilon_1\le\varepsilon_2$时集合包含关系为$\partial_{\varepsilon_1}f(x)\subseteq\partial_{\varepsilon_2}f(x)$, 且$\bigcap_{\varepsilon>0}\partial_\varepsilon f(x)=\partial f(x)$. 这些关系都直接来自定义的不等式.

若$\phi(x,z_x)\ge\sup_z\phi(x,z)-\varepsilon$, 任取$g\in\partial_x\phi(x,z_x)$, 则
\[
 f(y)\ge\phi(y,z_x)\ge\phi(x,z_x)+g\T(y-x)
 \ge f(x)+g\T(y-x)-\varepsilon.
\]
所以近似求解内部最大化就能得到误差可控的次梯度. 对偶算法里, 原子问题不必每次算到完全精确.

以$g_k\in\partial_{\varepsilon_k}f(x_k)$代替精确次梯度, 距离估计变成
\begin{equation}\label{eq:epsilon-distance}
 \norm{x_{k+1}-y}^2\le\norm{x_k-y}^2
 -2\alpha_k\bigl(f(x_k)-f(y)-\varepsilon_k\bigr)+\alpha_k^2 C^2.
\end{equation}
重复前面的矛盾证明, 常数$\alpha,\varepsilon$给$\liminf f(x_k)\le f^*+\varepsilon+\alpha C^2/2$; 若$\alpha_k\to0$, $\sum\alpha_k=\infty$, $\varepsilon_k\to\bar\varepsilon$, 则$\liminf f(x_k)\le f^*+\bar\varepsilon$. 这里所谓“到达误差最优集合”首先是目标最好值的保证, 并不自动意味着所有迭代点最后都留在该集合中.

最后说明误差次梯度与增量方法的联系. 若$g\in\partial f(\bar x)$, 则把次梯度不等式改写到$x$得
\[
 g\in\partial_\varepsilon f(x),\qquad
 \varepsilon=|f(x)-f(\bar x)|+\norm{g}\norm{x-\bar x}.
\]
若$f$在相关区域为$C$-Lipschitz且$\norm{g}\le C$, 则$\varepsilon\le2C\norm{x-\bar x}$. 又由定义逐项相加,
\[
 \partial_{\varepsilon_1}f_1(x)+\cdots+\partial_{\varepsilon_m}f_m(x)
 \subseteq\partial_{\sum_i\varepsilon_i}\left(\sum_i f_i\right)(x).
\]
在一轮增量迭代中, 各分量在彼此接近却不同的点计算次梯度, 因而可在轮起点理解为带小误差的次梯度. 约束投影逐次进行时, 整轮并不严格等于对次梯度和只做一次投影; 严格的整轮距离估计将在增量方法一章给出.

\originalchapter{多面体近似与内外线性化}\label{chap:8}
本章对应原第16讲后半、第17--18讲, 原PDF第214--259页. 切平面、单纯形分解及广义单调规划在本章统一处理. 原第18讲所依据的研究论文是Bertsekas与Yu的\emph{A Unifying Polyhedral Approximation Framework for Convex Optimization}, SIAM Journal on Optimization 21 (2011), 333--360. 本章给出幻灯片实际包含的算法与收敛结果, 不把混合情形的研究级分析压缩成无条件保证.

\originalsection{切平面与上下界}
次梯度$g_i\in\partial f(x_i)$给全局仿射下界$\ell_i(x)=f(x_i)+g_i\T(x-x_i)$. 对$f$取有限个下界的最大值,
\[
 F_k(x)=\max_{0\le i\le k}\ell_i(x),
\]
则$F_k\le f$, 且$F_{k+1}\ge F_k$. 其上图为若干半空间的交集, 包含真正的$\epi f$, 所以称为外线性化. “外”描述上图的位置, 不是函数值较大.

切平面法从$x_0\in X$开始, 反复求解
\[
 x_{k+1}\in\argmin_{x\in X}F_k(x),
\]
然后在新点取次梯度加入模型. 若$X$为多面体, 每轮可写成线性规划: 最小化$t$, 满足$x\in X$和$t\ge\ell_i(x)$对所有$i\le k$成立. 原变量维数基本不变, 但约束数逐轮增多.

设$L_k=F_k(x_{k+1})$及$U_k=\min_{i\le k}f(x_i)$. 则
\begin{equation}\label{eq:cut-bounds}
 L_k\le f^*\le U_k,\qquad L_{k+1}\ge L_k,\quad U_{k+1}\le U_k.
\end{equation}
因此$U_k-L_k$是可检验的绝对误差证书. 它来自下界模型与实际可行点, 不需知道$f^*$.

\begin{theorem}\label{thm:cutting-convergence}
设$X$非空闭凸, $f:\R^n\to\R$凸, 每轮模型最小值达到. 切平面迭代的每个聚点都是$f$在$X$上的最优解. 更一般地, 对闭真凸$f$, 同一结论在所取次梯度沿收敛子列有界时成立, 但须保证所有迭代点及极限点在相关定义域中可用.
\end{theorem}
\begin{proof}
取收敛子列$x_{k_j}\to\bar x$, 令$j>h$且$k_j>k_h$. 第$k_h$步加入的平面已在生成$x_{k_j}$的模型中, 所以对每个$z\in X$,
\[
 f(x_{k_h})+g_{k_h}\T(x_{k_j}-x_{k_h})
 \le F_{k_j-1}(x_{k_j})\le f(z).
\]
实值凸函数连续, 且次梯度在紧邻域上有界. 令$h,j\to\infty$, 混合内积趋零, 得$f(\bar x)\le f(z)$. 闭性给$\bar x\in X$. 扩展实值版本用下半连续性取下极限; 由有界次梯度和同一不等式保证$f(\bar x)\le f(z)<\infty$, 从而极限落在有效域内.
\end{proof}

如果$X$有界, 上述条件使序列有聚点, 因而最好值趋于$f^*$. 但从一个迭代点的目标值突然变差并不违背定理: 模型在远离采样点处可能非常乐观. 即使$x_0$恰为最优点, 第一张切平面也可能把下一步推向远处. 这是下一章引入近端稳定化的原因.

若$f=\max_{1\le j\le N}(a_j\T x+b_j)$为有限多面体函数, 并且次梯度计算规则每次选一个活跃的原始斜率$a_j$, 则有限终止: 非最优时新点的真实值严格高于当前模型值, 因而必须加入一张尚未包含的活跃平面; 至多加入$N$张. 若任意选择活跃斜率的凸组合, 可产生无限多不同平面, 不能用同一个有限计数证明.

原讲义还给出两个变体. 一是同时用分离超平面对$X$作外近似, 在违反约束时加入约束切面. 二是中心切平面法, 在下界模型与当前上界形成的定位区域
\[
 \{(x,t):x\in X,\ F_k(x)\le t\le U_k\}
\]
内选较居中的点, 再作查询和切割. 它控制查询点的几何位置, 不再总取下界模型的极小点. 具体的中心定义、删面规则和复杂度必须随具体算法给出.

\originalsection{单纯形分解}
当线性优化比非线性优化便宜, 还可从内部逐渐扩大可行集合. 设$X$是有界多面体, $f$可微凸. 初始$X_0=\{x_0\}\subseteq X$, 第$k$轮已有$x_k$在$\conv(X_k)$上最优. 先解大型但线性的子问题
\[
 \widetilde x_{k+1}\in\argmin_{x\in X}\nabla f(x_k)\T(x-x_k),
\]
选择一个达到最小值的极点. 加入$X_{k+1}=X_k\cup\{\widetilde x_{k+1}\}$, 再求
\[
 x_{k+1}\in\argmin_{x\in\conv(X_{k+1})}f(x).
\]
最后一个问题可用系数变量$\alpha$写成$\min f(\sum_{v\in X_{k+1}}\alpha_vv)$, 约束$\alpha_v\ge0$, $\sum_v\alpha_v=1$. 原空间即使很大, 有效优化变量数也仅为当前保存的点数.

\begin{theorem}
在上述条件下, 不删点的单纯形分解法有限次找到最优解.
\end{theorem}
\begin{proof}
记$d_k=\min_{x\in X}\nabla f(x_k)\T(x-x_k)\le0$, 因$x_k\in X$. 若$d_k=0$, 则可微凸函数的最优性条件说明$x_k$在$X$上最优. 若$d_k<0$, 因$x_k$已在$\conv(X_k)$上最优, 对该凸包内每个$x$有$\nabla f(x_k)\T(x-x_k)\ge0$. 所以$\widetilde x_{k+1}$不属于该凸包, 尤其未存入$X_k$. 每个非终止步骤加入一个新极点, 而有界多面体的极点有限, 故最终必须终止. 子问题最小值存在由紧性和连续性保证.
\end{proof}

可用$-d_k$作误差上界, 因对最优解$x^*$, 凸性给$f(x_k)-f^*\le\nabla f(x_k)\T(x_k-x^*)\le-d_k$. 它在不求$f^*$的情况下给停止准则.

原讲义建议丢弃满足$\nabla f(x_{k+1})\T(v-x_{k+1})>0$的旧极点, 因这些点不能以正权重出现在当前最优凸组合中. 这可减小子问题, 但删点后有限终止的上述计数论证不再直接适用, 不能沿用不删点版的结论. 无界多面体需要处理极射线或人工界; 非多面体集合仍可作线性优化, 但不再有有限极点论证. 非光滑$f$则需下面的原对偶匹配次梯度, 随意选一个次梯度一般不够.

\originalsection{函数的内线性化及共轭关系}
对闭真凸$h$与有限点集$V=\{v_1,\ldots,v_N\}\subseteq\dom h$, 定义内线性化
\begin{equation}\label{eq:inner-linearization}
 \overline h_V(z)=\begin{cases}
 \displaystyle\min\left\{\sum_j\alpha_jh(v_j):\sum_j\alpha_jv_j=z,\ \sum_j\alpha_j=1,\ \alpha_j\ge0\right\},&z\in\conv V,\\
 +\infty,&z\notin\conv V.
 \end{cases}
\end{equation}
其上图是有限条竖直向上射线$\{(v_j,t):t\ge h(v_j)\}$的凸包. 系数在紧单纯形中, 所以可行时最小值达到. 凸性给$h\le\overline h_V$, 而在每个采样点$v_j$有等号. 增加点使内模型下降, 同时使其有效域扩大.

对$f$的外模型, 若$y_j\in\partial f(x_j)$, Fenchel等式给$f(x_j)+f^*(y_j)=x_j\T y_j$, 因而
\[
 F(x)=\max_j\{x\T y_j-f^*(y_j)\}.
\]
这段表述既说明外模型由斜率决定, 也为下面的对偶关系消除采样点$x_j$.

\begin{proposition}\label{prop:inner-outer-conjugacy}
设$Y=\{y_1,\ldots,y_N\}$为上述斜率集合. 则$F^*=\overline{f^*}_Y$. 反向地, 对任意有限$V\subseteq\dom h$, 有
\[
 (\overline h_V)^*(y)=\max_{v\in V}\{v\T y-h(v)\}.
\]
\end{proposition}
\begin{proof}
先算第二式. 任一$z$可行表示成$\sum\alpha_vv$, 所以
\[
 y\T z-\sum_v\alpha_vh(v)=\sum_v\alpha_v\bigl(y\T v-h(v)\bigr)
 \le\max_{v\in V}\{y\T v-h(v)\}.
\]
取$z=v$达到对应项. 对$h=f^*$、$V=Y$, 第二式右边恰为$F$. 内模型为闭真凸多面体函数, 取双共轭即得第一式. 这也证明了“外模型的共轭是内模型”并非仅是图形比喻.
\end{proof}

\begin{center}
\begin{tikzpicture}[x=2.2cm,y=1.5cm]
 \draw[->] (-1.65,0)--(1.65,0) node[right]{$x$};
 \draw[->] (0,-.15)--(0,2.35) node[above]{$t$};
 \draw[thick] plot[domain=-1.48:1.48,samples=70] (\x,{\x*\x});
 \draw[structurecolor,thick] (-1.5,2)--(-.5,0)--(.5,0)--(1.5,2);
 \draw[dashed,thick] (-1,1)--(0,0)--(1,1);
 \fill (-1,1) circle (1.2pt); \fill (0,0) circle (1.2pt); \fill (1,1) circle (1.2pt);
 \node at (-1.05,2.2) {$f(x)=x^2$};
 \node[structurecolor] at (1.05,.45) {$F(x)$};
 \node at (-.75,.3) {$\overline f_V(x)$};
 \node[below] at (-1,0) {$-1$}; \node[below] at (1,0) {$1$};
\end{tikzpicture}
\end{center}
图中取$V=\{-1,0,1\}$, 外模型$F(x)=\max\{0,2x-1,-2x-1\}$位于抛物线下方, 内模型在$[-1,1]$上为$|x|$, 区间外为$+\infty$. 实线与虚线分别反映上图外包和内包.

\originalsection{广义单纯形分解与对偶切平面}
现在最小化$f(x)+c(x)$, $f,c$闭真凸. 用内线性化$C_k$替代$c$, 求$x_k\in\argmin(f+C_k)$. 在满足次梯度和规则的正则性条件下, 可取
\[
 g_k\in\partial f(x_k),\qquad -g_k\in\partial C_k(x_k).
\]
随后求$\widetilde x_{k+1}$满足$-g_k\in\partial c(\widetilde x_{k+1})$, 即
\[
 \widetilde x_{k+1}\in\argmin_x\{c(x)+g_k\T x\},
\]
并将此点加入内模型. 若$f$可微, $g_k=\nabla f(x_k)$; 若$f$不可微, 必须选能与内模型最优性相匹配的$g_k$. 这就是非光滑版本比“把梯度换成任意次梯度”更细的地方.

对$c=\ind{X}$, 内模型为$\ind{\conv V_k}$, 上式变成在$X$上最小化$g_k\T x$. 于是经典单纯形分解是其特殊情形. 同时Fenchel对偶是
\[
 \min_y\{f^*(y)+c^*(-y)\}.
\]
内模型$C_k$的共轭是$c^*$的外模型; 所以原问题的广义单纯形分解正对应对偶问题的切平面精化. 原、对偶计算只要使用同一匹配最优对, 就产生相同的数学迭代信息; 应选较方便求解的一侧实现.

\originalsection{广义单调规划的对称对偶}
令$x=(x_1,\ldots,x_m)$, 各块$x_i\in\R^{n_i}$, $S$为乘积空间中的线性子空间. 广义单调规划, 简称EMP, 写成
\begin{equation}\label{eq:emp-primal}
 \inf_{x\in S}\sum_{i=1}^m f_i(x_i),\qquad f_i\text{闭、真、凸}.
\end{equation}
若各块是一维, 它包含经典单调规划; 网络环流守恒给一个子空间, 每条弧的凸成本给一个分量. 取$S=\{(x,x):x\in\R^n\}$得到Fenchel问题. 取$S=\{(x,\ldots,x):x\in\R^n\}$则得到函数和最小化. 仿射约束$Ax=b$可增加变量$z$, 把$Ax-z=0$纳入子空间, 并加入$\ind{\{b\}}(z)$.

引入$z_i=x_i$和乘子$y_i$, 得
\[
 L(x,z,y)=\sum_i\bigl(f_i(z_i)+y_i\T(x_i-z_i)\bigr).
\]
对$z$取下确界给$-\sum_i f_i^*(y_i)$. 对$x\in S$取下确界, 当$y\in S^\perp$时为零, 否则为$-\infty$, 因子空间可沿任意正负方向缩放. 因而对偶的最小化写法是
\begin{equation}\label{eq:emp-dual}
 \inf_{y\in S^\perp}\sum_i f_i^*(y_i).
\end{equation}
注意\eqref{eq:emp-dual}的最小值是原拉格朗日对偶最大值的相反数. 强对偶时, 两个最小化问题的最优值相加为零, 不是相等.

\begin{proposition}\label{prop:emp-optimality}
一对有限目标值的可行点$x\in S$, $y\in S^\perp$同时最优且无对偶间隙, 当且仅当$y_i\in\partial f_i(x_i)$对所有$i$成立. 等价条件是$x_i\in\partial f_i^*(y_i)$.
\end{proposition}
\begin{proof}
Fenchel不等式给每个$r_i=f_i(x_i)+f_i^*(y_i)-x_i\T y_i\ge0$. 正交性给$\sum_i x_i\T y_i=0$, 所以原对偶目标之差是$\sum_i r_i$. 差为零当且仅当每项为零, 而每项为零等价于共轭次梯度关系. 反向同样由零间隙和弱对偶得到最优性.
\end{proof}

通常可用$S\cap\ri(\dom f_1\times\cdots\times\dom f_m)\ne\varnothing$作充分条件. 原第246页另给更宽的向量和条件: 若原问题可行, 对所有有效可行$x$和$\varepsilon>0$, 集合
\[
 S^\perp+\bigl(\partial_\varepsilon f_1(x_1)\times\cdots\times\partial_\varepsilon f_m(x_m)\bigr)
\]
闭, 则零对偶间隙成立. 这里的笛卡尔积等于把各分量嵌入对应块后作向量和; 不应误读为在不同维空间里直接相加. 这项高级对偶定理在本节作为已声明的工具使用, 不用它替代后面的算法证明.

该条件覆盖各$f_i$分别属于以下类别的情形: 在整个空间实值; 多面体; 本质一维; 定义域一维. 为说明原页的术语, 本质一维指$f_i(x)=h(a\T x)$, $h$为闭真凸一维函数; 定义域一维指$\aff(\dom f_i)$是点或直线. 两种一维情形的误差次微分都由一个闭区间及线性约束描述, 因而为多面体; 实值情形的误差次微分紧. 紧凸集与多面体的和闭, 解释了这些类别为何适用. 定义核对来源为Bertsekas的\href{https://web.mit.edu/dimitrib/www/Extended_Mono.pdf}{\emph{Extended Monotropic Programming and Duality}, 2010年修订稿}, 定义4.1--4.2. 交换原、对偶后, “实值”可改成“共轭实值”, 即co-finite, 但需要相应的对偶可行性.

\originalsection{混合多面体近似算法}
把指标集分成互不相交的三类$E,O,I$: 保持精确、外线性化、内线性化. 对$i\in O$保存有限斜率集$Y_i$, 对$i\in I$保存有限点集$V_i$. 每轮求近似EMP的原对偶最优对$(\widehat x,\widehat y)$:
\[
 \min_{x\in S}\left\{
 \sum_{i\in E}f_i(x_i)+\sum_{i\in O}f_{i,Y_i}(x_i)
 +\sum_{i\in I}\overline f_{i,V_i}(x_i)\right\}.
\]
然后按分量精化:
\begin{enumerate}
\item 对$i\in O$, 取$\widetilde y_i\in\partial f_i(\widehat x_i)$, 加入$Y_i$.
\item 对$i\in I$, 取$\widetilde x_i\in\partial f_i^*(\widehat y_i)$, 加入$V_i$.
\end{enumerate}
第二步等价于求$f_i(x_i)-\widehat y_i\T x_i$的极小点. 算法定义隐含了近似原对偶最优对及这些查询的存在, 应在具体应用中核对; 闭真凸本身不保证任意共轭次梯度都存在.

对偶近似问题把$f_{i,Y_i}^*$变成内模型, 把$(\overline f_{i,V_i})^*$变成外模型. 原对偶角色交换不改变点和斜率的精化过程. 与经典每轮只加一张整体切平面不同, 这里一轮可同时精化多个分量; 若分量是一维, 查询很便宜, 并能保留原成本或网络的分离结构. 混合模型既有下界分量又有上界分量, 因此其目标值一般既不是原最优值的下界也不是上界; 不能直接套用\eqref{eq:cut-bounds}.

\begin{proposition}
假设每个外线性化分量是实值多面体函数, 每个内线性化分量的共轭是实值多面体函数, 且每次加入的斜率或点从相应有限多面体表示中选取. 若近似最优对及查询一直存在, 算法有限次得到原问题的最优原对偶对.
\end{proposition}
\begin{proof}
如果某一轮所有查询都未精化模型, 则对外分量, 新查询平面已在模型中, 所以模型在$\widehat x_i$与$f_i$接触. 近似问题的匹配次梯度$\widehat y_i$于是也属于$\partial f_i(\widehat x_i)$. 对内分量, 在共轭侧作同一论证, 得$\widehat x_i\in\partial f_i^*(\widehat y_i)$. 精确分量本已满足同一关系. 命题\ref{prop:emp-optimality}证明该对原问题最优. 否则至少有一个分量加入新的表示元素, 而总元素数有限, 故不能无限精化.
\end{proof}

\begin{proposition}\label{prop:emp-outer-limit}
对纯外线性化情形$I=\varnothing$, 若每个外分量生成的查询次梯度序列有界, 则近似原最优点序列的每个聚点原最优. 对纯内线性化情形$O=\varnothing$, 若生成的查询点序列有界, 则近似对偶最优点序列的每个聚点对偶最优.
\end{proposition}
\begin{proof}
证明第一项. 对任一有效可行$x\in S$, 两个先后轮次$k>\ell$的模型最优性给
\begin{multline*}
 \sum_{i\in E}f_i(\widehat x_i^k)+
 \sum_{i\in O}\left[f_i(\widehat x_i^\ell)
 +\widetilde y_i^{\ell\mathsf T}(\widehat x_i^k-\widehat x_i^\ell)\right]\\
 \le\sum_i f_i(x_i).
\end{multline*}
沿同一收敛子列取$k,\ell\to\infty$, 有界次梯度使所有内积趋零. 每个闭真凸函数有仿射下界, 因此在有界子列上其值有统一下界, 可以对有限和使用下半连续性, 得$\sum_i f_i(\bar x_i)\le\sum_i f_i(x_i)$. 子空间闭, 故$\bar x\in S$. 第二项交换共轭和子空间, 由内外共轭关系得到.
\end{proof}

这两个结果分别控制原聚点与对偶聚点, 不是混合情形的无条件收敛证明. 原第258页明确将一般混合情形的分析交给Bertsekas--Yu论文. 在需要混合算法的具体实现时, 应另检查查询有界性、模型精化误差和原对偶可行性, 不能仅以“每轮更多切面”替代这些条件. 近端项也可加入近似问题以稳定迭代, 并把已得的内外模型带入下一轮; 下一章解释这种稳定化为何有用.

\originalchapter{近端稳定化与增广拉格朗日}\label{chap:9}
本章对应原第19--20讲及第21讲的广义近端部分, 原PDF第260--284页. 主线是同一个二次正则化子问题的原问题实现、共轭实现和约束优化实现; bundle方法则把精确目标换成切平面模型.

\originalsection{近端子问题的存在与基本恒等式}
设$f:\R^n\to(-\infty,+\infty]$闭、真、凸. 对$c>0$及任意$z\in\R^n$, 定义
\[
 \operatorname{prox}_{cf}(z)=\argmin_x\left\{f(x)+\frac1{2c}\norm{x-z}^2\right\}.
\]
闭真凸函数有仿射下界$f(x)\ge a\T x+b$. 加上正二次项后, 目标下水平集有界, 且目标下半连续, 故最小值达到. 二次项的严格凸性给唯一性. 因此起点即使不在$\dom f$内, 第一步也定义良好并落入有效域. 约束$x\in X$可用$f+\ind{X}$吸收, 但需此和函数为真.

近端点法是
\begin{equation}\label{eq:proximal-point}
 x_{k+1}=\operatorname{prox}_{c_k f}(x_k),\qquad c_k>0.
\end{equation}
最优性条件为
\begin{equation}\label{eq:prox-subgradient}
 g_{k+1}:=\frac{x_k-x_{k+1}}{c_k}\in\partial f(x_{k+1}).
\end{equation}
这是一种隐式次梯度步: 使用新点的次梯度, 不是旧点的次梯度. 隐式性使每步更贵, 也带来稳定性. 若$x_k$已经最优, 则$0\in\partial f(x_k)$, 因而$x_{k+1}=x_k$; 近端点法不会像无稳定化切平面法那样主动离开最优点.

\begin{lemma}\label{lem:prox-distance}
对任意$y\in\dom f$, 有
\begin{equation}\label{eq:prox-distance}
 \norm{x_{k+1}-y}^2\le\norm{x_k-y}^2
 -2c_k\bigl(f(x_{k+1})-f(y)\bigr)-\norm{x_k-x_{k+1}}^2.
\end{equation}
若$x_k\in\dom f$, 还满足
\[
 f(x_{k+1})+\frac1{2c_k}\norm{x_{k+1}-x_k}^2\le f(x_k).
\]
\end{lemma}
\begin{proof}
由\eqref{eq:prox-subgradient}, $g_{k+1}\T(x_{k+1}-y)\ge f(x_{k+1})-f(y)$. 展开$x_k-y=(x_{k+1}-y)+c_k g_{k+1}$的平方, 移项得到第一式. 第二式是用$x_k$作为子问题的比较点. 若$x_k$不是最优点, 子问题唯一极小点不等于$x_k$, 此时目标确实严格下降.
\end{proof}

与显式次梯度距离估计相比, 这里没有$+c_k^2\norm{g}^2$误差, 反而多了一个负的移动距离项. 不需假设所有次梯度有统一范数界.

\originalsection{收敛、增长阶与有限终止}
\begin{theorem}\label{thm:prox-convergence}
若$\sum_{k=0}^\infty c_k=\infty$, 则$f(x_k)\to f^*=\inf f$. 若最优集合$X^*$非空, 则$x_k$趋于其中一个点.
\end{theorem}
\begin{proof}
从第一步起目标值单调不增. 将\eqref{eq:prox-distance}求和并删去非正项, 对任意$y\in\dom f$及$N\ge1$有
\[
 2\sum_{k=0}^{N-1}c_k\bigl(f(x_{k+1})-f(y)\bigr)\le\norm{x_0-y}^2.
\]
因$f(x_N)\le f(x_{k+1})$, 得
\[
 f(x_N)\le f(y)+\frac{\norm{x_0-y}^2}{2\sum_{k=0}^{N-1}c_k}.
\]
令$N\to\infty$, 再对$y$取下确界, 得值收敛, 也包括$f^*=-\infty$的情形.

若$x^*\in X^*$, 距离不等式使$\norm{x_k-x^*}$单调不增, 故序列有界. 取一个收敛子列, 下半连续性与值收敛保证其极限$\bar x\in X^*$. 到$\bar x$的距离同样单调并沿该子列趋零, 所以整个序列趋于$\bar x$.
\end{proof}

参数$c_k$越大, 对离开当前点的二次惩罚越弱, 子问题越接近原目标. 但大$c_k$也可能使子问题数值求解困难. 若目标在最优集合附近有较陡增长, 一个近端步可跨越更多距离.

\begin{proposition}\label{prop:prox-growth}
设$X^*\ne\varnothing$, $d(x)=\dist(x,X^*)$. 假设在相关近邻内
\[
 f(x)-f^*\ge\theta d(x)^p,\qquad\theta>0,\quad p\ge1.
\]
若$x_{k+1}$在该近邻且非最优, 则
\begin{equation}\label{eq:prox-growth-distance}
 d(x_k)\ge d(x_{k+1})+c_k\theta d(x_{k+1})^{p-1}.
\end{equation}
\end{proposition}
\begin{proof}
记$u=x_{k+1}$, $v=x_k$, $g=(v-u)/c_k$. 取$z_u=P_{X^*}(u)$, 次梯度不等式与柯西不等式给
\[
 \theta d(u)^p\le f(u)-f^*\le g\T(u-z_u)\le\norm{g}d(u),
\]
故$\norm{g}\ge\theta d(u)^{p-1}$. 再取$z_v=P_{X^*}(v)$, 展开并用同一次梯度不等式,
\begin{align*}
 d(v)^2&=\norm{u-z_v+c_k g}^2\\
 &\ge d(u)^2+2c_k\theta d(u)^p+c_k^2\norm{g}^2\\
 &\ge\bigl(d(u)+c_k\theta d(u)^{p-1}\bigr)^2.
\end{align*}
取非负平方根即得结论.
\end{proof}

当$p=2$且$c_k\to\bar c>0$, 有$\limsup d(x_{k+1})/d(x_k)\le1/(1+\theta\bar c)<1$, 即线性距离收敛. 当$1<p<2$且$c_k\ge\underline c>0$, 有
\[
 d(x_{k+1})\le\left(\frac{d(x_k)}{\underline c\theta}\right)^{1/(p-1)},
\]
幂指数大于一, 所以距离收敛超线性. 下界$\underline c$在此结论中不能省略; 若$c_k$不断趋零, 正则化可能越来越强.

当$p=1$, 非终止步满足$d(x_k)-d(x_{k+1})\ge\theta c_k$. 若$\sum c_k=\infty$, 永远不终止将使非负距离下降到负数, 所以有限终止. 若增长界全局成立且$c_0\ge d(x_0)/\theta$, 第一步即最优. 有最优解的多面体凸函数具有相应的线性误差界, 这是原第265页有限终止结论的基础.

\begin{example}
对$f(x)=|x|$, 显式求一个近端步并解释有限终止.
\end{example}
\begin{solution}
子问题$\min_u|u|+(u-x)^2/(2c)$的最优性条件是$0\in\partial|u|+(u-x)/c$. 在$u>0$时解为$u=x-c$, 在$u<0$时为$u=x+c$, 在$u=0$时条件是$|x|\le c$. 因而
\[
 \operatorname{prox}_{c|\cdot|}(x)=\operatorname{sign}(x)\max\{|x|-c,0\}.
\]
非零步使绝对值减少$c$, 累计参数之和超过$|x_0|$时便到达零. 此例同时表明参数变化不能在有限终止讨论里被忽略.
\end{solution}

\originalsection{共轭实现与Moreau包络}
对二次函数$r_{k}(x)=\norm{x-x_k}^2/(2c_k)$, 共轭由配方给$r_k^*(v)=x_k\T v+(c_k/2)\norm{v}^2$. 所以近端子问题的Fenchel对偶可写为
\begin{equation}\label{eq:dual-proximal}
 y_{k+1}=\argmin_y\left\{f^*(y)-x_k\T y+\frac{c_k}{2}\norm{y}^2\right\}.
\end{equation}
二次函数处处有限, 保证适用强对偶. 原对偶问题都有唯一解, 并满足
\[
 y_{k+1}\in\partial f(x_{k+1}),\qquad
 x_{k+1}=x_k-c_k y_{k+1}.
\]
因此可以先求共轭侧的子问题, 再重建原点. 选择哪侧只由结构便利决定, 不是另一种不同的近端迭代. 若$c_k$有正下界且$x_k$收敛, 则$y_{k+1}\to0$. 没有参数条件时, 不能仅由$x_k-x_{k+1}\to0$就除以可能趋零的$c_k$.

定义Moreau包络
\[
 e_c f(z)=\inf_x\left\{f(x)+\frac1{2c}\norm{x-z}^2\right\},\qquad u(z)=\operatorname{prox}_{cf}(z).
\]
它把非光滑凸函数变成可微凸函数, 并且$\nabla e_c f(z)=(z-u(z))/c$. 为说明可微性, 比较在$z$和$z+h$处的最优解, 得
\begin{align*}
 e_cf(z+h)-e_cf(z)&\le\frac1c(z-u(z))\T h+\frac1{2c}\norm h^2,\\
 e_cf(z+h)-e_cf(z)&\ge\frac1c(z-u(z+h))\T h+\frac1{2c}\norm h^2.
\end{align*}
近端映射非扩张: 对两点的最优性条件作差, 次梯度单调性给
\[
 \norm{u(z)-u(w)}^2\le(u(z)-u(w))\T(z-w).
\]
故$u(z+h)-u(z)=O(\norm h)$, 上下界的线性主部相同, 余量为$O(\norm h^2)$, 得到梯度公式. 同一单调性还给\mbox{$\norm{\nabla e_cf(z)-\nabla e_cf(w)}\le\norm{z-w}/c$}.

包络与$f$有相同最优值和最优集合: 显然$e_cf(z)\ge f^*$, 最优点达到等号; 反向, 包络最优点的梯度为零, 故$z=u(z)$且$0\in\partial f(z)$. 近端步于是也是步长$c_k$的包络梯度步. 原第278页由此指出可以在平滑包络上考虑Newton或拟Newton加速, 但该建议不等于给出它们的无条件收敛结论.

\originalsection{近端切平面与bundle方法}
切平面下界$F_k\le f$与二次稳定项组合成
\[
 u_k\in\argmin_{x\in X}\left\{F_k(x)+\frac1{2c_k}\norm{x-y_k}^2\right\}.
\]
$u_k$是查询点, $y_k$是近端中心, 两者不必相同. 如果每轮都把中心移到查询点, 就是近端切平面法. 参数很大时稳定作用弱, 又接近普通切平面法; 参数很小时可能过于保守. 存储的切面很多时, 二次子问题也会变贵.

bundle方法保留一个稳定中心, 用真实下降检验决定是否移动. 假设中心已被查询且$F_k(y_k)=f(y_k)$, 定义预测下降
\begin{equation}\label{eq:bundle-prediction}
 \delta_k=f(y_k)-\left(F_k(u_k)+\frac1{2c_k}\norm{u_k-y_k}^2\right)\ge0.
\end{equation}
取固定$\beta\in(0,1)$. 若$f(y_k)-f(u_k)\ge\beta\delta_k$, 令$y_{k+1}=u_k$, 称实步. 否则中心不动, 称空步. 两种情形都在$u_k$求次梯度并加入新切面. 空步虽然没有接受查询点, 却修正了过度乐观的模型.

当$\delta_k=0$, 强凸性使$u_k=y_k$. 此时$y_k$最小化模型加稳定项, 稳定项在中心梯度为零; 因模型与$f$相切, 其最优性条件也是原问题的最优性条件. 所以零预测下降可以作精确停止证书.

\begin{theorem}\label{thm:basic-bundle-convergence}
为给出可完整证明的基本版本, 设$X=\R^n$, $f$实值凸, 初始水平集$\{x:f(x)\le f(y_0)\}$紧, $c_k=c>0$, 始终保留全部切面. 则上述bundle法的中心目标值趋于$f^*$, 中心的每个聚点最优.
\end{theorem}
\begin{proof}
中心目标值单调不增, 因而中心留在紧水平集内. 记$s_k=(y_k-u_k)/c\in\partial F_k(u_k)$. 强凸性给
\[\delta_k\ge\frac{\norm{u_k-y_k}^2}{2c}.\]
对任意$z$, 模型次梯度不等式给
\[
 f(z)\ge F_k(u_k)+s_k\T(z-u_k)
 =f(y_k)+s_k\T(z-y_k)-\varepsilon_k,
\]
其中$\varepsilon_k=\delta_k-\norm{u_k-y_k}^2/(2c)\in[0,\delta_k]$. 故$s_k$是中心处的$\varepsilon_k$次梯度, 并有$\norm{s_k}\le\sqrt{2\delta_k/c}$.

若有无限多个实步, 中心目标的总下降有限, 而每个实步下降至少$\beta\delta_k$, 所以实步上的$\delta_k\to0$. 用任一最优点$z=x^*$代入上式, 中心有界使内积趋零, 从而实步前的中心目标趋于$f^*$. 单调性使所有中心目标也趋于$f^*$.

若只有有限多个实步, 从某一步起中心固定为$y$. 写$Q_k=F_k+\norm{\cdot-y}^2/(2c)$, $q_k=Q_k(u_k)$. 模型单调增加, 故$q_k$不减且不超过$f(y)$. 一张固定旧切面与二次项给统一的强制有界下水平集, 所以查询点有界, 在这些点的次梯度也有界. 强凸性给
\[
 q_{k+1}-q_k\ge\frac1{2c}\norm{u_{k+1}-u_k}^2,
\]
故相邻查询点之差趋零. 新切面在$u_k$接触$f$, 所以
\[
 q_{k+1}\ge f(u_k)+g_k\T(u_{k+1}-u_k)+\frac1{2c}\norm{u_{k+1}-y}^2.
\]
与$q_k=F_k(u_k)+\norm{u_k-y}^2/(2c)$相减, 得模型误差$f(u_k)-F_k(u_k)\to0$. 但空步检验说明
\[
 f(u_k)-F_k(u_k)>(1-\beta)\delta_k+\frac1{2c}\norm{u_k-y}^2.
\]
所以$\delta_k\to0$, 再用上述误差次梯度不等式得到$f(z)\ge f(y)$对所有$z$, 即固定中心已经最优. 两种情形均完成值收敛, 聚点最优由连续性得到.
\end{proof}

实用bundle算法常聚合旧切面、调整参数或设置允许误差. 它们应保留足够的模型信息, 需要各自的分析; 原讲义的通式不意味着任意删面、任意参数更新仍满足本定理.

\originalsection{增广拉格朗日作为对偶近端法}
考虑凸等式问题$\min\{f(x):x\in X,\ Ex=d\}$. 定义扰动值函数
\[
 p(u)=\inf\{f(x):x\in X,\ Ex-d=u\},\qquad
 q(\mu)=\inf_{x\in X}\{f(x)+\mu\T(Ex-d)\}=-p^*(-\mu).
\]
假设$p$闭、真、凸, 对偶最优乘子存在, 所需子问题下确界达到. 在凹对偶目标上作近端最大化,
\[
 \mu_{k+1}=\arg\max_\mu\left\{q(\mu)-\frac1{2c_k}\norm{\mu-\mu_k}^2\right\},
\]
其共轭实现是在扰动变量上最小化$p(u)+\mu_k\T u+(c_k/2)\norm u^2$, 然后令$\mu_{k+1}=\mu_k+c_k u_{k+1}$. 代入$p$的定义可实现为
\begin{equation}\label{eq:augmented-lagrangian}
 \begin{split}
 x_{k+1}&\in\argmin_{x\in X}\left\{f(x)+\mu_k\T(Ex-d)+\frac{c_k}{2}\norm{Ex-d}^2\right\},\\
 \mu_{k+1}&=\mu_k+c_k(Ex_{k+1}-d).
 \end{split}
\end{equation}
花括号中的函数是增广拉格朗日函数. 普通二次惩罚只加最后一项; 增广拉格朗日还不断校正线性乘子, 因而不必单靠让惩罚参数趋于无穷来消除违反量.

验证其对偶近端身份也可直接用最优性. 设$u_{k+1}=Ex_{k+1}-d$. 扰动子问题的最优性给$-\mu_k-c_k u_{k+1}\in\partial p(u_{k+1})$. 令$\mu_{k+1}=\mu_k+c_k u_{k+1}$, 共轭次梯度关系表明$x_{k+1}$最小化普通拉格朗日函数$L(x,\mu_{k+1})$, 并且$u_{k+1}$为$q$在新乘子处的超梯度. 因而$u_{k+1}=(\mu_{k+1}-\mu_k)/c_k$正好满足对偶近端最优性.

若$\sum c_k=\infty$, 对$-q$应用定理\ref{thm:prox-convergence}, 乘子趋于一个对偶最优解. 若还满足$c_k\ge\underline c>0$, 则残差$Ex_{k+1}-d\to0$. 由
\[
 q(\mu_{k+1})=f(x_{k+1})+\mu_{k+1}\T(Ex_{k+1}-d)
\]
得$f(x_{k+1})\to f^*$. 若原点序列有聚点且$X$闭、$f$下半连续, 该聚点原最优. 乘子收敛与原变量有界性是不同事项, 不能把后者省略.

对切平面模型$F_k+\ind{X}$同样可作原或共轭侧的近端子问题. 不受约束时, $F_k^*$正是共轭目标的内线性化; 有约束时还涉及与$\ind X$的共轭卷积, 不能机械把它当作单纯的有限点凸包. 在完整Fenchel模型中两种实现仍等价. 若只在通过进展检验后移动中心, 对应外模型的bundle与共轭内模型的bundle两种表达.

\originalsection{广义正则化与上界最小化}
广义近端迭代写成$x_{k+1}\in\argmin_x\{f(x)+D_k(x,x_k)\}$. 这里不再有二次项自动保证存在性, 必须另假设子问题达到最小值. 若$D_k(x,x_k)\ge D_k(x_k,x_k)$, 则比较点$x_k$给
\[
 f(x_{k+1})\le f(x_{k+1})+D_k(x_{k+1},x_k)-D_k(x_k,x_k)\le f(x_k).
\]
若$f$凸, $D_k(\cdot,x_k)$凸且在$x_k$可微, 两者满足相对内部相交条件, 则算法在$x_k$停住意味着$x_k$最优: 正则项在其极小点$x_k$梯度为零, 和函数的最优性于是给$0\in\partial f(x_k)$. 这些条件也保证非最优点的严格目标改善. 对非凸$f$, 这种全局最优结论一般失效, 不能只凭下降性声称收敛到全局极小点.

Bregman正则化取$D_k(x,y)=[\phi(x)-\phi(y)-\nabla\phi(y)\T(x-y)]/c_k$, 在$\phi$可微的区域使用. 凸性使其非负并在$x=y$为零; 它不一定对称, 也不一定满足三角不等式. 熵与镜像下降将在后一章计算.

上界最小化, 即MM方法, 则构造$M_k(x,y)\ge f(x)$及$M_k(y,y)=f(y)$, 令$x_{k+1}\in\argmin_x M_k(x,x_k)$. 比较得到$f(x_{k+1})\le M_k(x_{k+1},x_k)\le f(x_k)$. 若要写成$f+D$形式, 应令$D_k(x,y)=M_k(x,y)-f(x)$, 而不是把上界值减去一个只依赖$y$的常数后再次加到$f$上. 原第284页有关MM的两种通式没有保持这一对应, 本稿采用与上界最小化定义一致的写法.

\begin{example}
对$f(x)=R(x)+\norm{Ax-b}^2$, 其中$R$闭真凸, 构造一个容易求解的MM上界.
\end{example}
\begin{solution}
取$\tau\ge\norm A^2$并且$\tau>0$, 令
\[
 M(x,y)=R(x)+\norm{Ax-b}^2-\norm{A(x-y)}^2+\tau\norm{x-y}^2.
\]
最后两项之和非负, 故$M(x,y)\ge f(x)$且$M(y,y)=f(y)$. 展开平方后,
\[
 M(x,y)=R(x)+\norm{Ay-b}^2+2(Ay-b)\T A(x-y)+\tau\norm{x-y}^2.
\]
因此最小化$M$等价于计算$R$的近端映射, 中心为$y-A\T(Ay-b)/\tau$, 参数为$1/(2\tau)$. 原页写的系数一对应$\norm A\le1$的情形; 不核对谱范数便不能保证它是上界. EM是统计模型中另一类上界最小化方法, 原页只提名称, 并未给特定似然模型, 本章不据此捏造一个通用的全局收敛定理.
\end{solution}

\originalchapter{内点障碍与增量方法}\label{chap:10}
本章对应原第21讲内点部分与第22讲, 原PDF第285--300页, 并接续第16讲第212--213页的增量方法引入. 障碍法从严格可行区域逼近边界最优点, 增量法则利用目标的巨大分量和结构; 两者处理的是不同的计算困难.

\originalsection{障碍法及其极限}
设$f,g_1,\ldots,g_r$为连续凸函数, 可行集$C=\{x:g_j(x)\le0,\ j=1,\ldots,r\}$, 严格可行集$S=\{x:g_j(x)<0,\ j=1,\ldots,r\}$非空. 在$S$上取连续障碍函数$B$, 要求当点从内部趋于边界时$B\to+\infty$. 常用的是
\[
 B(x)=-\sum_{j=1}^r\ln(-g_j(x)),\qquad
 B(x)=-\sum_{j=1}^r\frac1{g_j(x)}.
\]
这些函数都使靠近约束边界变得昂贵. 与违反约束后再处罚的惩罚法不同, 障碍函数只在严格可行区使用.

取$\varepsilon_k>0$, $\varepsilon_k\downarrow0$, 每轮求
\[
 x_k\in\argmin_{x\in S}\{f(x)+\varepsilon_k B(x)\}.
\]
子问题最小值达到仍需假设, 例如需要相应有界下水平集; 障碍阻止接近有限边界, 并不自动阻止逃向无穷远.

\begin{theorem}\label{thm:barrier-limit}
在上述条件下, 障碍法所生成序列的每个聚点都是原约束问题的最优解.
\end{theorem}
\begin{proof}
取$x_{k_j}\to\bar x$. 因$C$闭, $\bar x\in C$. 若它不是最优, 则有$z\in C$满足$f(z)<f(\bar x)$. 取一个严格可行点$s$, 对充分小的$t>0$, $\widetilde z=(1-t)z+ts$严格可行: 凸性给$g_i(\widetilde z)\le(1-t)g_i(z)+tg_i(s)<0$. 连续性还给$f(\widetilde z)<f(\bar x)$. 子问题最优性于是给
\[
 f(x_{k_j})+\varepsilon_{k_j}B(x_{k_j})
 \le f(\widetilde z)+\varepsilon_{k_j}B(\widetilde z).
\]
若$\bar x\in S$, 障碍值连续, 故左边附加项趋零. 若$\bar x\in\partial S$, 障碍值趋于$+\infty$, 故该附加项下极限非负. 两种情形都推出$f(\bar x)\le f(\widetilde z)$, 矛盾.
\end{proof}

严格可行性在证明中有明确作用: 保证每个较优可行点都能由严格可行点逼近. 缺少这种可逼近性时, 不能沿用这段证明.

\begin{example}
最小化$\frac12(x_1^2+x_2^2)$, 约束$x_1\ge2$, 用对数障碍求近似解并考察病态性.
\end{example}
\begin{solution}
原最优解是$(2,0)$. 子问题为
\[
 \min_{x_1>2,x_2}\left\{\frac12(x_1^2+x_2^2)-\varepsilon\ln(x_1-2)\right\}.
\]
驻点条件给$x_2=0$及$x_1-\varepsilon/(x_1-2)=0$, 所以$x_1=1+\sqrt{1+\varepsilon}$, 另一根不在严格可行区. 随$\varepsilon\downarrow0$趋于2. 子问题在该点的Hessian为
\[
 \operatorname{diag}\left(1+\frac\varepsilon{(x_1-2)^2},1\right).
\]
因为$x_1-2=\varepsilon/(\sqrt{1+\varepsilon}+1)\sim\varepsilon/2$, 第一特征值约为$4/\varepsilon$, 而第二个恒为1. 越接近边界, 条件数越差, 说明简单固定步长梯度法会困难, 应用Newton方向和保持严格可行的线搜索.
\end{solution}

若$f,g_j$可微, 对数障碍子问题的驻点可写成
\[
 \nabla f(x_\varepsilon)+\sum_j\mu_j\nabla g_j(x_\varepsilon)=0,
 \qquad \mu_j=\frac\varepsilon{-g_j(x_\varepsilon)}>0.
\]
于是$\mu_j g_j(x_\varepsilon)=-\varepsilon$, 是互补松弛的平滑近似. 在凸问题的其他可行条件成立时, 这些乘子给对偶可行解, 原对偶间隙不超过$r\varepsilon$. 这解释了沿中心路径降低障碍参数为何具有可检验的精度含义.

\originalsection{锥的对数障碍}
二阶锥$\mathcal Q_d$的内部写成$t>\norm u$. 标准障碍为
\[
 B(u,t)=-\ln(t^2-\norm u^2),\qquad t>\norm u.
\]
定义域必须指定正分支; 仅写$t^2>\norm u^2$会错误地包含$t<0$的另一支. 对SOCP$\min\{c\T x:A_i x-b_i\in\mathcal Q_{d_i}\}$, 近似问题是
\[
 \min_x\left\{c\T x+\varepsilon\sum_i B_i(A_i x-b_i)\right\},\quad
 A_i x-b_i\in\operatorname{int}\mathcal Q_{d_i}.
\]
Newton方向利用障碍的Hessian, 线搜索必须让所有锥残差保持在内部.

半正定锥的障碍是$-\ln\det Z$, 定义在$Z\succ0$. 对对偶SDP
\[
 \max_\mu b\T\mu\quad\text{满足}\quad Z(\mu)=D-\sum_i\mu_i A_i\succeq0,
\]
障碍问题为最大化$b\T\mu+\varepsilon\ln\det Z(\mu)$, 限制$Z(\mu)\succ0$. 其导数是$b_i-\varepsilon\tr(Z^{-1}A_i)$. 在驻点令$X_\varepsilon=\varepsilon Z^{-1}\succ0$, 则$\inner{A_i}{X_\varepsilon}=b_i$, 因而得到一个原可行矩阵. 同时
\[
 X_\varepsilon Z=\varepsilon I,\qquad
 \inner{D}{X_\varepsilon}-b\T\mu
 =\inner{Z}{X_\varepsilon}=n\varepsilon.
\]
这既是矩阵互补条件的近似, 也是直接的原对偶误差证书. 原讲义在第288页提到内点复杂度分析, 未陈述自协调障碍或完整复杂度定理; 本章不把仅有这些障碍公式当成已经证明多项式复杂度.

\originalsection{大规模分量和与极限环}
设$f(x)=\sum_{i=1}^m f_i(x)$, $m$很大. 分量可代表每个数据样本的损失、分离对偶成本或情景成本. 例如最小二乘和$\ell_1$正则化,
\[
 \lambda\norm{x}_1+\frac12\sum_{i=1}^m(c_i\T x-d_i)^2,
\]
其中$\norm{x}_1=\sum_j|x_j|$. 随机规划中的期望$\mathbb E[F(x,\omega)]$也可用样本分量近似, 或逐次查询随机分量. 两阶段模型$F_1(x)+\mathbb E_\omega[\inf_y F_2(x,y,\omega)]$先选择$x$, 待情景$\omega$出现后再优化补救变量$y$; 若相应联合模型凸且部分最小化有效, 情景成本同样提供凸分量. 多约束问题和分布式优化也可把不同约束或节点的信息组织成逐次更新.

全梯度或全次梯度需要遍历所有分量. 增量法每次只处理$i_k$:
\[
 x_{k+1}=P_X(x_k-\alpha_k g_{i_k}(x_k)),\qquad
 g_i(x)\in\partial f_i(x).
\]
一个由$m$次分量更新组成的轮次, 工作量与一次全次梯度近似相当, 却在中间不断更新点的位置. 这可能显著提升实际效果, 但固定步长一般产生极限环.

\begin{example}
依次对$f_1(x)=\frac12(1-x)^2$和$f_2(x)=\frac12(1+x)^2$作增量梯度步, 使用固定$\alpha\in(0,2)$.
\end{example}
\begin{solution}
从轮起点$a$出发, 第一步为$u=(1-\alpha)a+\alpha$, 第二步为$v=(1-\alpha)u-\alpha=(1-\alpha)^2a-\alpha^2$. 因$(1-\alpha)^2<1$, 轮起点趋于$a_*=-\alpha/(2-\alpha)$, 中间点趋于$u_*=\alpha/(2-\alpha)$. 原总目标是$1+x^2$, 唯一最优点为零; 两点极限环的振幅随步长趋零而趋零, 固定步长却不会消失.
\end{solution}

\begin{example}
对$|1-x|+|1+x|+|x|$按三个分量的顺序更新. 取$0<\alpha<1/2$, 轮起点$a\in(0,\alpha)$. 展示固定步长极限环对起点的依赖.
\end{example}
\begin{solution}
这一轮依次得到$a+\alpha$, $a$, $a-\alpha$. 下一轮从$a-\alpha<0$出发依次得到$a$, $a-\alpha$, $a$. 因而轮末点在$a$和$a-\alpha$间交替, 振荡位置取决于$a$. 总目标在$[-1,1]$上为$2+|x|$, 最优点仍为零. 此处无需虚构非唯一最优集合来解释振荡.
\end{solution}

\originalsection{循环次序的整轮估计}
为把原讲义的“有界次梯度”写成可核对的充分条件, 假设各$f_i$在所有相关点之间为$C$-Lipschitz, 所用次梯度范数不超过$C$. 若$X$紧且函数在其开邻域实值凸, 这类统一界通常成立. 一轮从$x_k$开始, 以同一个步长$\alpha_k$依次作
\[
 \psi_0=x_k,\quad \psi_i=P_X(\psi_{i-1}-\alpha_k g_i(\psi_{i-1})),\quad
 x_{k+1}=\psi_m.
\]
本节下标$k$计整轮, 与前节计单次分量更新的下标不同.

\begin{lemma}\label{lem:incremental-cycle}
对任意$y\in X$, 有
\[
 \norm{x_{k+1}-y}^2\le\norm{x_k-y}^2
 -2\alpha_k\bigl(f(x_k)-f(y)\bigr)+\alpha_k^2m^2C^2.
\]
\end{lemma}
\begin{proof}
非扩张性给$\norm{\psi_i-\psi_{i-1}}\le\alpha_k C$, 所以$\norm{\psi_{i-1}-x_k}\le(i-1)\alpha_k C$. 每步次梯度距离估计求和, 再用
\[
 f_i(\psi_{i-1})\ge f_i(x_k)-C\norm{\psi_{i-1}-x_k}
 \ge f_i(x_k)-(i-1)\alpha_kC^2,
\]
得到误差$\alpha_k^2 C^2[m+2\sum_{i=1}^m(i-1)]=\alpha_k^2m^2C^2$.
\end{proof}

由与次梯度法同样的距离矛盾论证, 固定步长给$\liminf_k f(x_k)\le f^*+\alpha m^2C^2/2$. 若$\alpha_k\to0$, $\sum\alpha_k=\infty$, 则轮末历史最好值趋于$f^*$. 若还有最优解存在、下半连续性和$\sum\alpha_k^2<\infty$, 则轮末点趋于一个最优点; 轮内距离至多$mC\alpha_k\to0$, 所以所有分量迭代点同样收敛. 逐轮投影与“先加全部次梯度再投影”并不相同, 但整轮距离估计把两者的差异控制在二阶步长误差内.

\originalsection{独立随机次序与概率保证}
若每步$i_k$独立于过去, 在$\{1,\ldots,m\}$中均匀抽取, 则在当前信息$\mathcal F_k$条件下,
\begin{equation}\label{eq:random-incremental-distance}
 \mathbb E[\norm{x_{k+1}-y}^2\mid\mathcal F_k]
 \le\norm{x_k-y}^2-\frac{2\alpha_k}{m}(f(x_k)-f(y))+\alpha_k^2C^2.
\end{equation}
这由对每个分量的距离估计取条件平均得到. 与循环整轮误差$m^2C^2$相比, 随机单步的误差只有$C^2$, 但目标下降项也除以$m$.

\begin{proposition}
在上述统一有界和Lipschitz条件下, 固定步长$\alpha$满足
\[
 \liminf_k f(x_k)\le f^*+\frac{\alpha mC^2}{2}\quad\text{几乎必然}.
\]
若$\alpha_k\to0$且$\sum\alpha_k=\infty$, 则$\liminf_k f(x_k)=f^*$几乎必然.
\end{proposition}
\begin{proof}
固定$y\in X$与$\delta>0$. 固定步长时, 在$f(x_k)>f(y)+\alpha mC^2/2+\delta$的区域, 条件期望中的平方距离每步至少减少$2\alpha\delta/m$. 停在第一次进入该水平集的时刻, 对停止后的非负距离取期望并逐步求和. 若永远不进入有正概率, 该求和中的期望累计减少将无穷大, 与初始有限距离矛盾. 从任意确定时刻重新应用同一论证, 得每个尾部都进入该水平集, 即下极限不超过这个阈值. 对可数个逼近$f^*$的$y$和趋零的$\delta$取交集, 得第一项.

递减步长时, 取足够晚的起点使$\alpha_k C^2\le\delta/m$. 在$f(x_k)>f(y)+\delta$区域, 条件漂移不超过$-\alpha_k\delta/m$. 由于$\sum\alpha_k=\infty$, 同样的停止求和论证保证每个尾部都进入这个水平集. 令$\delta\downarrow0$并用可数逼近点, 得第二项.
\end{proof}

如果要求点序列几乎必然收敛, 可再加$\sum\alpha_k^2<\infty$及最优解存在, 用带可求和误差的非负超鞅收敛定理获得到最优集合的距离收敛. 本命题没有把下极限保证误称为每次目标都收敛. 每轮开始随机打乱全部分量的顺序, 即不放回随机重排, 在实践中也常用; 它与这里的独立有放回抽样不同, 不能直接以\eqref{eq:random-incremental-distance}逐步条件平均.

\originalsection{增量近端与次梯度的组合}
对$f_i$的约束近端步$z=\argmin_{x\in X}\{f_i(x)+\norm{x-v}^2/(2\alpha)\}$, 在次梯度和规则适用时, 有$g\in\partial f_i(z)$使
\[
 z=P_X(v-\alpha g).
\]
因为最优性条件是$v-z-\alpha g\in N_X(z)$, 恰为投影条件. 这是新点次梯度的隐式更新; 对只有扩展实值定义的函数, 必须核对相对内部等条件以保证可拆出这样的$g$.

设$F_i=f_i+h_i$, 总目标$F=\sum_iF_i$. 可对较容易的$f_i$用近端, 对其他$h_i$用显式次梯度:
\[
 z_k\in\argmin_{x\in X}\left\{f_{i_k}(x)+\frac1{2\alpha_k}\norm{x-x_k}^2\right\},\qquad
 x_{k+1}=P_X(z_k-\alpha_k h'_{i_k}(z_k)).
\]
也可以调换次序或把近端步先放在整个空间上, 但相应的可行性和误差分析需重新检查.

为把原第298页的“小常数”明确化, 假设$f_i,h_i$都在相关区域为$C$-Lipschitz, 所用隐式和显式次梯度均有界于$C$, 并满足上述近端拆分条件. 近端距离估计加上次梯度距离估计, 再把$F_i(z_k)$比较到$F_i(x_k)$, 得单步界
\[
 \norm{x_{k+1}-y}^2\le\norm{x_k-y}^2
 -2\alpha_k(F_{i_k}(x_k)-F_{i_k}(y))+5\alpha_k^2C^2.
\]
其中$\norm{z_k-x_k}\le\alpha_k C$, 所以替换两个函数值至多引入$4\alpha_k^2C^2$, 显式步另有$\alpha_k^2C^2$. 每个完整分量步移动至多$2\alpha_k C$. 循环一轮, 再比较到轮起点, 误差上界是
\[
 [5m+8\sum_{i=1}^m(i-1)]\alpha_k^2C^2
 =(4m^2+m)\alpha_k^2C^2\le5m^2\alpha_k^2C^2.
\]
因此原讲义的统一写法$\beta\alpha_k^2m^2C^2$可在这些明确充分条件下取$\beta=5$. 固定步长的循环值界为$F^*+5\alpha m^2C^2/2$, 独立均匀抽样的几乎必然下极限界为$F^*+5\alpha mC^2/2$. 递减步长时二阶误差消失, 得到与前面相同的最好值收敛结论. 这些数值常数是本稿为补齐论证给出的保守界, 不是原作者宣称的最佳常数.

\begin{example}
对$\lambda\norm{x}_1+\frac12\sum_i(c_i\T x-d_i)^2$, 给出两个可直接实现的分量近端公式.
\end{example}
\begin{solution}
一维绝对值近端公式逐坐标应用, 给
\[
 \operatorname{prox}_{\alpha\lambda\norm{\cdot}_1}(v)_j
 =\operatorname{sign}(v_j)\max\{|v_j|-\alpha\lambda,0\}.
\]
因此小坐标直接被置零. 对$q_i(x)=\frac12(c_i\T x-d_i)^2$, 近端最优性是$z-v+\alpha c_i(c_i\T z-d_i)=0$. 左乘$c_i\T$先解标量残差, 得
\[
 c_i\T z-d_i=\frac{c_i\T v-d_i}{1+\alpha\norm{c_i}^2},\qquad
 z=v-\frac{\alpha(c_i\T v-d_i)}{1+\alpha\norm{c_i}^2}c_i.
\]
若将正则项均摊到$m$个分量, 每个分量应使用$\lambda/m$而不是$\lambda$; 或将它作为独立的第$m+1$个分量每轮处理一次. 两种组织都保持原目标, 反复对每个样本施加完整$\lambda$却会把正则化强度乘上$m$.
\end{solution}

\originalchapter{一阶复杂度、复合近端与镜像下降}\label{chap:11}
本章对应原第23--24讲, 原PDF第301--316页, 并接续第21讲的Bregman正则化. 用统一的距离或势函数估计比较非光滑、光滑与加速算法, 最后计算熵近端及乘法更新. 复杂度所计的一步必须指明所需查询; 一个难解近端子问题不能视作与一次廉价次梯度计算成本相同.

\originalsection{从渐近收敛到迭代次数}
设最优点$x^*$存在, 记$R=\norm{x_0-x^*}$. 对次梯度法使用固定步长$\alpha$及$\norm{g_k}\le G$, 距离估计求和给
\[
 2\alpha\sum_{k=0}^{N-1}(f(x_k)-f^*)\le R^2+N\alpha^2G^2.
\]
因最小值不大于平均值,
\begin{equation}\label{eq:subgradient-complexity}
 \min_{0\le k<N}f(x_k)-f^*\le\frac{R^2}{2\alpha N}+\frac{\alpha G^2}{2}.
\end{equation}
若预先给定次数$N$, 取$\alpha=R/(G\sqrt N)$得到误差至多$RG/\sqrt N$. 若要求误差$\varepsilon$, 可取$\alpha=\varepsilon/G^2$, 再令$N\ge R^2G^2/\varepsilon^2$. 因而有$O(\varepsilon^{-2})$的迭代次数上界. 这里输出的是历史最好点, 或用凸性输出平均点, 不是自动把最后点当作最好点.

参数$R,G$若只知上界, 可用相应上界选步长; 若未知, 需要自适应或阶段式策略, 不应把依赖未知最优点的公式直接当作可无条件实现的程序.

\originalsection{光滑性、投影与二次模型}
设$f$凸可微, 梯度在包含所用线段的凸区域内满足$\norm{\nabla f(x)-\nabla f(y)}\le L\norm{x-y}$. 令$\omega(z;x)=f(x)+\nabla f(x)\T(z-x)$. 沿线段积分,
\begin{align*}
 f(z)-\omega(z;x)
 &=\int_0^1[\nabla f(x+t(z-x))-\nabla f(x)]\T(z-x)\,\mathrm dt\\
 &\le\int_0^1 Lt\norm{z-x}^2\,\mathrm dt
 =\frac L2\norm{z-x}^2.
\end{align*}
凸性给另一侧$f(z)\ge\omega(z;x)$. 所以$f$被切平面与一个二次上界夹住.

梯度投影步
\[
 u=P_X(x-\alpha\nabla f(x))
 =\argmin_{z\in X}\left\{\omega(z;x)+\frac1{2\alpha}\norm{z-x}^2\right\}
\]
是对目标作线性化后的近端步. 投影最优性给$\nabla f(x)\T(u-x)\le-\norm{u-x}^2/\alpha$, 因而
\[
 f(u)\le f(x)-\left(\frac1\alpha-\frac L2\right)\norm{u-x}^2.
\]
所以$0<\alpha<2/L$保证非零步的严格下降. 不过下一项简洁复杂度估计使用更保守的$\alpha\le1/L$.

二次模型的强凸性对所有$z\in X$给
\[
 \omega(u;x)+\frac1{2\alpha}\norm{u-x}^2
 \le\omega(z;x)+\frac1{2\alpha}\norm{z-x}^2-\frac1{2\alpha}\norm{z-u}^2.
\]
结合光滑上界和凸性下界, 得到核心比较式
\begin{equation}\label{eq:gradient-threepoint}
 f(u)-f(z)\le\frac{\norm{x-z}^2-\norm{u-z}^2}{2\alpha}
 -\left(\frac1{2\alpha}-\frac L2\right)\norm{u-x}^2.
\end{equation}
它比只证明$f(u)\le f(x)$更强, 因为直接把目标误差与到任意比较点的距离变化相连.

\begin{theorem}\label{thm:gradient-complexity}
设$X$非空闭凸, 最优解存在, 取$\alpha=1/L$. 梯度投影法满足对所有$N\ge1$,
\[
 f(x_N)-f^*\le\frac{L\dist(x_0,X^*)^2}{2N}.
\]
\end{theorem}
\begin{proof}
在\eqref{eq:gradient-threepoint}取$z=x^*$, 对$k=0,\ldots,N-1$求和, 距离项望远镜消去, 得$\sum_{k=0}^{N-1}(f(x_{k+1})-f^*)\le L\norm{x_0-x^*}^2/2$. 目标值单调, 每项至少为$f(x_N)-f^*$, 再对$x^*$取最近最优点即得.
\end{proof}

于是光滑性把$O(\varepsilon^{-2})$改善到$O(\varepsilon^{-1})$. 若$L$未知, 可从试探步长出发, 按固定因子缩小直到$f(u)\le\omega(u;x)+\norm{u-x}^2/(2\alpha)$. 该检验直接验证本次二次上界, 足以使用同一比较式. 在有统一Lipschitz常数且试探规则给已接受步长正下界时, 同阶复杂度仍成立.

\originalsection{复合目标的梯度近端法}
设$H=f+g+\ind X$, 其中$f$光滑凸, $g$闭真凸, $X$闭凸, 和函数为真. 保留易处理但可能不可微的$g$, 只线性化$f$:
\begin{equation}\label{eq:forward-backward}
 u=\argmin_{z\in X}\left\{\omega(z;x)+g(z)+\frac1{2\alpha}\norm{z-x}^2\right\}.
\end{equation}
配方给等价实现: 先$v=x-\alpha\nabla f(x)$, 再$u=\operatorname{prox}_{\alpha(g+\ind X)}(v)$. 这把不方便求近端的光滑部分交给梯度计算, 把有闭式或廉价近端的非光滑部分保留下来.

令$q(z)=\omega(z;x)+g(z)+\norm{z-x}^2/(2\alpha)+\ind X(z)$. 子问题强凸性给$q(u)\le q(z)-\norm{z-u}^2/(2\alpha)$. 与前节完全相同的夹逼得到
\begin{equation}\label{eq:composite-threepoint}
 H(u)-H(z)\le\frac{\norm{x-z}^2-\norm{u-z}^2}{2\alpha}
 -\left(\frac1{2\alpha}-\frac L2\right)\norm{u-x}^2,
 \quad z\in\dom H.
\end{equation}
因此$\alpha\le1/L$保证复合目标下降, 并给$H(x_N)-H^*\le\norm{x_0-x^*}^2/(2\alpha N)$. 取$g=0$就是梯度投影, 取$f=0$就是近端点法.

多个光滑分量可以先合成一个$f$, 但每步必须计算其全梯度. 这不是每次只选一个分量的增量算法; 循环增量的常步长极限环与本节完整梯度近端的定理并不矛盾. 原第312页关于“不能推广到多个分量”的警示应在保持这种逐分量更新结构的意义下理解.

\originalsection{梯度法上界的量级与动量}
\begin{example}
固定$L>0$, 取参数$h>0$, 定义
\[
 f_h(x)=\begin{cases}
 Lx^2/2,&|x|\le h,\\
 Lh|x|-Lh^2/2,&|x|>h.
 \end{cases}
\]
从$x_0=R>h$以步长$1/L$作梯度下降, 说明$O(\varepsilon^{-1})$上界的量级无法只靠固定步长梯度法改善.
\end{example}
\begin{solution}
两段在$\pm h$处的值和一阶导数一致, 导数为$L\max\{-h,\min\{x,h\}\}$, 故梯度的Lipschitz常数为$L$. 只要$x_k>h$, 就有$x_{k+1}=x_k-h$. 为把函数族参数与目标精度分开, 给定小$\varepsilon$, 取$h=2\varepsilon/(LR)$. 在线性段$f_h(x)\le\varepsilon$要求$x\le R/2+h/2$. 因此至少要约$R/(2h)=LR^2/(4\varepsilon)$步, 与前节上界同阶. 这是随精度变化的一族问题, 不应把函数参数$h$与目标误差$\varepsilon$混成一个固定问题的两种距离.
\end{solution}

经典重球法加入固定动量,
\[
 x_{k+1}=x_k-\alpha\nabla f(x_k)+\beta(x_k-x_{k-1}),\qquad0<\beta<1.
\]
另一种形式先外推$y_k=x_k+\beta(x_k-x_{k-1})$, 再在$y_k$处作梯度投影. 对上例的线性段, 两种形式都使位移满足$v_{k+1}=\beta v_k-h$. 从零初速度开始, $|v_k|\le h/(1-\beta)$, 所以固定$\beta$至多改善常数, 迭代次数仍为$O(\varepsilon^{-1})$量级. 动量的稳定性也取决于参数与目标, 不能仅凭$\beta\in(0,1)$就宣称所有凸问题收敛.

\originalsection{递增外推与加速势函数}
对复合目标$H=f+g+\ind X$, 假设$f$在所有外推点及相应线段上凸、$L$-光滑. 外推点未必在$X$内, 因此只假设梯度在$X$上Lipschitz不足以定义或分析算法. 本节以$f$在整个空间$L$-光滑作为方便的充分条件.

令$x_{-1}=x_0\in\dom H$, $\theta_{-1}=1$, $\theta_k=2/(k+2)$对$k\ge0$. 取
\[
 \beta_0=0,\qquad\beta_k=\frac{k-1}{k+2}\ (k\ge1),\qquad
 y_k=x_k+\beta_k(x_k-x_{k-1}),
\]
然后在$y_k$处作步长$1/L$的复合梯度近端步. 一般系数关系为$\beta_k=\theta_k(1-\theta_{k-1})/\theta_{k-1}$. 这里$\beta_k\to1$, 但不是随意加大动量, 而是与下面的势函数恒等式匹配. 原第307页行文中的初始索引与其显式例子不完全一致, 本节按上述明确索引定义算法.

\begin{theorem}\label{thm:accelerated-complexity}
上述算法满足对所有$N\ge1$,
\[
 H(x_N)-H^*\le\frac{2L\norm{x_0-x^*}^2}{(N+1)^2}.
\]
所以达到目标误差$\varepsilon$所需次数为$O(\varepsilon^{-1/2})$.
\end{theorem}
\begin{proof}
记$\Delta_k=H(x_k)-H^*$, 并定义
\[
 w_k=x_k+\frac{1-\theta_{k-1}}{\theta_{k-1}}(x_k-x_{k-1}),\qquad w_0=x_0.
\]
于是$y_k=(1-\theta_k)x_k+\theta_k w_k$, 而
$w_{k+1}=x_k+(x_{k+1}-x_k)/\theta_k$. 在\eqref{eq:composite-threepoint}取线性化中心$y_k$, 比较点$z=(1-\theta_k)x_k+\theta_k x^*$. 凸性给$H(z)\le(1-\theta_k)H(x_k)+\theta_k H^*$. 又
\[
 y_k-z=\theta_k(w_k-x^*),\qquad x_{k+1}-z=\theta_k(w_{k+1}-x^*).
\]
所以
\begin{equation}\label{eq:accelerated-energy}
 \frac{\Delta_{k+1}}{\theta_k^2}+\frac L2\norm{w_{k+1}-x^*}^2
 \le\frac{1-\theta_k}{\theta_k^2}\Delta_k+\frac L2\norm{w_k-x^*}^2.
\end{equation}
对$k\ge1$, 所选参数满足
\[
 \frac{1-\theta_k}{\theta_k^2}=\frac{k(k+2)}4
 \le\frac{(k+1)^2}4=\frac1{\theta_{k-1}^2}.
\]
故势函数逐步不增. 第零步$\theta_0=1$使$\Delta_0$项消失, 给初始势至多$L\norm{x_0-x^*}^2/2$. 迭代得到$\Delta_N/\theta_{N-1}^2\le L\norm{x_0-x^*}^2/2$, 代入$\theta_{N-1}=2/(N+1)$即得. 证明也解释了为什么要使用这组看似不直观的外推系数.
\end{proof}

加速算法的目标值未必每步下降, 势函数下降代替了单步下降. 对纯光滑约束问题取$g=0$, 就得到原第307页的加速投影法; 对有廉价近端的$g$同一证明仍有效.

在无约束且最优解存在的情形, 对实值、全局$G$-Lipschitz的非光滑凸$f$, 还可先用Moreau包络$e_cf$平滑. 下界与误差界为
\[
 0\le f(x)-e_cf(x)\le\frac{cG^2}{2}.
\]
因为$f(u)\ge f(x)-G\norm{u-x}$, 对$u$最小化右边加二次项即可得到下界$e_cf(x)\ge f(x)-cG^2/2$. 包络梯度常数为$1/c$. 若取$c=\varepsilon/G^2$, 对包络用加速法达到$\varepsilon/2$误差, 则原目标总误差不超过$\varepsilon$, 外层次数为$O(GR/\varepsilon)$. 但每次包络梯度本身需要解近端子问题. 原第308页的$O(\varepsilon^{-1})$应按这种更强的查询能力理解, 不能据此说任意廉价次梯度黑箱都已有同样复杂度. 平滑也可能破坏逐分量的增量结构.

\originalsection{Bregman正则化与三点恒等式}
设$\phi$为闭真凸函数, 在其有效域内部的开凸集上可微严格凸, 定义
\[
 D_\phi(x,y)=\phi(x)-\phi(y)-\nabla\phi(y)\T(x-y).
\]
它非负, 对严格凸情形$x\ne y$时为正; 通常不是对称距离. 二次函数$\phi(x)=\norm{x}^2/2$给普通平方距离. 广义近端步最小化$H(x)+D_\phi(x,x_k)/c_k$, 在定义域、子问题存在性和次梯度和规则成立时, 有
\[
 \frac{\nabla\phi(x_k)-\nabla\phi(x_{k+1})}{c_k}\in\partial H(x_{k+1}).
\]
也就是在梯度坐标$\nabla\phi(x)$中作隐式次梯度更新.

Bregman三点恒等式为
\begin{equation}\label{eq:bregman-threepoint}
 D_\phi(z,x)-D_\phi(z,u)-D_\phi(u,x)
 =[\nabla\phi(u)-\nabla\phi(x)]\T(z-u).
\end{equation}
逐项展开定义即得, 所有$\phi$值相消后只剩梯度项. 对广义近端步, 次梯度不等式与此式结合给
\[
 c_k(H(x_{k+1})-H(z))\le
 D_\phi(z,x_k)-D_\phi(z,x_{k+1})-D_\phi(x_{k+1},x_k).
\]
这是欧氏近端距离估计的真正推广. 若要进一步推出点收敛, 还需用$D_\phi$控制点间距离、核对边界行为及参数条件, 不能把形式相似当作全部证明.

把$D_\phi(\cdot,x_k)/c_k$的共轭直接配算, 得
\[
 r_k^*(s)=\frac1{c_k}\left[\phi^*(\nabla\phi(x_k)+c_k s)-\phi^*(\nabla\phi(x_k))\right].
\]
故Fenchel对偶最小化$H^*(y)+\phi^*(\nabla\phi(x_k)-c_k y)/c_k$, 可忽略与$y$无关的常数. 若$\nabla\phi^*$在相关点存在, 原点由$\nabla\phi^*(\nabla\phi(x_k)-c_k y)$重建. 内外线性化、bundle及增量组合也可接入这种正则化, 但每个组合都要保持相应的存在性与误差条件.

\originalsection{熵近端及指数增广拉格朗日}
取非负正交锥上的熵函数
\[
 \phi(x)=\sum_i x_i(\ln x_i-1),\qquad0\ln0:=0.
\]
负分量处取$+\infty$. 对$y_i>0$,
\[
 D_\phi(x,y)=\sum_i\left[x_i\ln\frac{x_i}{y_i}-x_i+y_i\right].
\]
原第314页的目标省略了仅依赖旧点的$\sum_i y_i/c_k$, 不影响极小点. 熵近端最小化$H(x)+D_\phi(x,x_k)/c_k$. 若$H$的非负有效域包含一个各分量严格正的点, 则一个已存在的有限极小点不能有零分量: 沿通往严格正点的线段, 零分量产生$t\ln t$项, 其右导数为$-\infty$, 而凸性控制$H$的增长至多为线性, 因此边界点不是极小点. 若没有严格正可行点, 则不能照搬“所有后续点都正”的假设.

一维$p(t)=t(\ln t-1)$在$t\ge0$上定义, 共轭计算为
\[
 p^*(s)=\sup_{t\ge0}\{st-t\ln t+t\}=e^s,
\]
因为导数$ s-\ln t=0$给$t=e^s$, 且目标严格凹. 因而熵近端的对偶, 使用标准共轭$H^*(y)=\sup_x(y\T x-H(x))$, 是
\begin{equation}\label{eq:entropy-prox-dual}
 \min_y\left\{H^*(y)+\frac1{c_k}\sum_i x_{k,i}e^{-c_k y_i}\right\},\qquad
 x_{k+1,i}=x_{k,i}e^{-c_k y_i}.
\end{equation}
若换变量$v=-y$, 则写成$H^*(-v)+\sum_i x_{k,i}e^{c_kv_i}/c_k$, 重建点使用正指数. 原第314--315页把$H^*(y)$与正指数同时写在同一约定下, 少了这一负号; 本节明确保留相容的共轭和指数符号.

对凸不等式问题$\min\{f(x):x\in X,\ g_j(x)\le0\}$, 在非负乘子上对凹对偶函数$q$作熵近端最大化, 则得到指数增广拉格朗日法:
\begin{equation}\label{eq:exponential-alm}
 \begin{split}
 x_{k+1}&\in\argmin_{x\in X}\left\{f(x)+\frac1{c_k}\sum_j\mu_{k,j}e^{c_k g_j(x)}\right\},\\
 \mu_{k+1,j}&=\mu_{k,j}e^{c_k g_j(x_{k+1})},\qquad\mu_{0,j}>0.
 \end{split}
\end{equation}
这里正指数是正确的: 若$x_{k+1}$使子问题最优性与原凸模型的和规则成立, 新乘子恰为各指数项的导数系数, 所以$x_{k+1}$最小化$L(x,\mu_{k+1})$. 因而$g(x_{k+1})$是$q$在新乘子处的超梯度, 而熵近端最优性为
\[
 -g_j(x_{k+1})+\frac1{c_k}\ln\frac{\mu_{k+1,j}}{\mu_{k,j}}=0,
\]
与更新式吻合. 违反量为正使乘子按比例增加, 严格满足约束使之减少. 所有乘子保持正, 但可在极限中趋零. 这一推导建立算法等价性; 具体收敛还需非二次近端的相关条件, 不由指数公式单独保证.

\originalsection{镜像下降与单纯形上的乘法更新}
镜像下降进一步把$H$的主要目标线性化. 设$g_k\in\partial f(x_k)$, 在闭凸集$X$上取
\[
 x_{k+1}\in\argmin_{x\in X}\{\alpha_k g_k\T(x-x_k)+D_\phi(x,x_k)\}.
\]
最优性是$[\alpha_k g_k+\nabla\phi(x_{k+1})-\nabla\phi(x_k)]\T(z-x_{k+1})\ge0$, $z\in X$. 对二次$\phi$, 它正是投影次梯度法.

设$\phi$相对某范数$\norm{\cdot}$为$\sigma$-强凸, 即$D_\phi(u,x)\ge\sigma\norm{u-x}^2/2$, 次梯度的对偶范数满足$\norm{g_k}_*\le G$. 用最优性及\eqref{eq:bregman-threepoint}, 得
\begin{align*}
 \alpha_k(f(x_k)-f(z))
 &\le\alpha_k g_k\T(x_k-x_{k+1})
   +D_\phi(z,x_k)-D_\phi(z,x_{k+1})-D_\phi(x_{k+1},x_k)\\
 &\le D_\phi(z,x_k)-D_\phi(z,x_{k+1})+\frac{\alpha_k^2G^2}{2\sigma}.
\end{align*}
第二步是对$r=\norm{x_k-x_{k+1}}$最大化$\alpha_kGr-\sigma r^2/2$. 求和后用凸性, 对加权平均点$\bar x_N=(\sum_{k<N}\alpha_kx_k)/(\sum_{k<N}\alpha_k)$得到
\begin{equation}\label{eq:mirror-rate}
 f(\bar x_N)-f(z)\le
 \frac{D_\phi(z,x_0)+(G^2/(2\sigma))\sum_{k<N}\alpha_k^2}{\sum_{k<N}\alpha_k}.
\end{equation}
欧氏平方距离被与约束几何匹配的Bregman量替换, 算法的基本进展机制仍是望远镜求和.

在概率单纯形$X=\{x\ge0:\sum_i x_i=1\}$上, 取熵$\phi$及$x_{0,i}>0$. 子问题去掉常数后为
\[
 \min_{x\in X}\sum_i\left[x_i g_{k,i}+\frac1{\alpha_k}x_i\ln\frac{x_i}{x_{k,i}}\right].
\]
对等式约束加乘子$\lambda$, 驻点给$g_{k,i}+(\ln(x_i/x_{k,i})+1)/\alpha_k+\lambda=0$. 用总和为一解出公因子,
\begin{equation}\label{eq:entropy-mirror-update}
 x_{k+1,i}=\frac{x_{k,i}e^{-\alpha_k g_{k,i}}}{\sum_jx_{k,j}e^{-\alpha_k g_{k,j}}}.
\end{equation}
分母归一化保证概率总和, 指数保证正性. 目标次梯度较大的坐标减小权重, 较小的坐标增大相对权重.

熵在单纯形上相对$\ell_1$范数是1-强凸. 对内部线段$x(t)=x+th$, $\sum_i x_i(t)=1$, 柯西不等式给
\[
 h\T\nabla^2\phi(x(t))h=\sum_i\frac{h_i^2}{x_i(t)}\ge\left(\sum_i|h_i|\right)^2.
\]
沿线段二次积分得$D_\phi(u,x)\ge\norm{u-x}_1^2/2$, 边界点可取极限. 对偶范数是$\ell_\infty$, 因而若$\norm{g_k}_\infty\le G$, 可在\eqref{eq:mirror-rate}取$\sigma=1$. 从均匀分布$x_{0,i}=1/n$开始, 对任意最优点有
\[
 D_\phi(x^*,x_0)=\ln n+\sum_i x_i^*\ln x_i^*\le\ln n.
\]
取常数步长$\alpha=\sqrt{2\ln n}/(G\sqrt N)$, 平均点误差至多$G\sqrt{2\ln n/N}$, 当$n=1$时问题仅有一个可行点, 无需该步长公式. 这段补充估计说明选择熵并非只为获得一个漂亮的指数更新式: 它也让初始距离与维数的依赖从欧氏几何改成$\ln n$.

\originalchapter{线性规划建模与单纯形法}\label{chap:12}
本章补充 MIT 15.053 \emph{Optimization Methods in Management Science}, Spring 2013, James Orlin 与 Ebrahim Nasrabadi 课程中的第2--6、9讲及教程1、2、4--7. 它们补足6.253已讲的线性对偶理论所未展开的建模、标准形、单纯形计算、Phase I与灵敏度. 原课件中的休息插图、时事、项目安排不属于数学正文; 明示排除开放许可的出版社题面只保留来源指引, 相关数学方法在此独立推导.

\originalsection{变量、参数与线性模型}
模型要分别说明能控制的变量、预先给定的参数、比较方案的目标及限制方案的约束. 变量名称不是小事: 若约束跨多天、多个阶段或多个产品, 只记录单日总量可能丢失原问题的逻辑. 单位也应匹配; 时间必须统一为分钟或小时, 比例与总量不能随意相加. 线性模型采用可分性、比例性与确定系数等假设, 输出最优解后还要检查这些假设是否合理.

\begin{example}
教程1第5--9页的生产例: 产品P、Q分别需要机器A的50、24分钟和机器B的30、33分钟, 两台机器可用40、35小时; 单位利润为25、30美元. 写出连续生产模型.
\end{example}
\begin{solution}
令 $x,y$ 为两产品数量, 时间统一成分钟, 得 $\max(25x+30y)$, 约束 $50x+24y\le2400$, $30x+33y\le2100$, $x,y\ge0$. 原教程第7页的一行目标写成 $40x+35y$, 与同页说明及完整模型不符; 此处按单位利润25、30及后续完整模型计算. 若产品不可拆分, 应另要求整数; 连续模型只是其线性松弛, 不能把小数产量解释为可直接实施的整数方案.
\end{solution}

\begin{example}
教程1第10--17页的三产品模型, 机器时间矩阵与销售上限为
\[
A=\begin{pmatrix}20&10&10\\12&28&16\\15&6&16\\10&15&0\end{pmatrix},\quad
b=2400\mathbf1,\quad u=(100,40,60)\T.
\]
单位销售收入为 $(90,100,70)$, 材料成本为 $(45,40,20)$, 固定周费用6000美元. 求并验证最优连续生产方案.
\end{example}
\begin{solution}
模型为 $\max(45p+60q+50r-6000)$, $A(p,q,r)\T\le b$, $0\le(p,q,r)\T\le u$. 固定费用不影响最优变量, 但报告实际利润时不可遗漏. 取 $r=60$ 并使前两台机器满负荷, 解 $20p+10q=1800$, $12p+28q=1440$, 得 $p=900/11$, $q=180/11$. 后两台机器使用 $25140/11$ 与 $11700/11$ 分钟, 都不超过2400; 变量上限也满足.

为证明最优, 对前两台机器约束分别用乘子 $27/22,75/44$, 对 $r\le60$ 用乘子 $115/11$, 其余置零. 加权后各变量目标系数恰为45、60、50, 给上界 $84300/11$. 候选方案达到该界, 所以最优. 扣除6000后利润为 $18300/11$. 原教程的数值约7664及1664是四舍五入; 产品Q表中16.63与公式不一致, 正确近似为16.36.
\end{solution}

广告、营养、混合配料模型的共同结构是 $\min\sum_jc_jx_j$, $\ell_i\le\sum_ja_{ij}x_j\le u_i$, $0\le x_j\le d_j$. 索引集、系数与单位要在公式之前定义. 广告模型中 $a_j$ 表示一次展示带来的计数, 不自动等于不重复触达人数; 若同一受众会重复看到广告, 需说明计数假设. 混合配料中若总质量固定为 $P$, 成分约束可写 $L_jP\le\sum_i a_{ji}x_i+M_j\le U_jP$, 重量守恒为 $\sum_i x_i+\sum_jM_j=P$. 这里只将纯成分 $M_j$ 加在第 $j$ 个成分中, 不把所有纯材料同时加进每个成分约束. 教程2的外部熔炉书摘以这一抽象模型替代逐字重译, 原书出处保留在来源台账中.

\begin{example}
第2讲的历史饮食模型允许分数份量, 成本向量 $(1,3,2.5,3,1)$, 热量 $(250,770,360,190,230)$, 脂肪热量 $(81,360,144,45,99)$, 蛋白质 $(31,44,14,27,3)$, 钠 $(480,1170,800,580,160)$. 要求总热量600--900, 脂肪热量不超过总热量的40\%, 蛋白质至少30, 钠不超过1150. 写出模型并核对课件给的近似方案.
\end{example}
\begin{solution}
令份量向量 $x\ge0$, 总热量 $F=a\T x$. 目标 $\min c\T x$, 约束 $600\le F\le900$, $d\T x\le0.4F$, $p\T x\ge30$, $s\T x\le1150$. 仅用前两项, 令热量下界与脂肪比例取等号, 解得 $x_1=1040/921$, $x_2=380/921$, 其他为0, 成本 $2180/921\approx2.367$. 蛋白质及钠限制均满足. 可用热量与脂肪两条活跃约束构造非负对偶乘子证明其最优: 设热量下界乘子 $\alpha=109/27630$, 脂肪限制乘子 $\beta=2/2763$, 则 $250\alpha+19\beta=1$, $770\alpha-52\beta=3$; 其他食品的同一加权系数不超过其成本, 给下界 $600\alpha=2180/921$.

原第14页钠数据是1170, 第15页公式却写1770; 后者也与给定近似解不相容, 本稿按第14页数据建模. 此处数字是2013课程例题, 不作为现行营养、价格或饮食建议. 整数模型中一份第1项与两份第5项成本为3, 满足所有条件; 连续最优值2.367本身只给整数下界. 但成本低于3时, 第2与第4项不能选, 第3项最多一份且不能与任何其他正份量组合, 热量仅360; 或只选第1与第5项且总份数至多2, 热量最多500. 这些都违反600热量下界, 所以成本3确为整数最优.
\end{solution}

\begin{example}
第2讲的投资组合例给历史期望收益率 $(12.7,9.9,11.8)$ 与协方差矩阵
$Q=\left(\begin{smallmatrix}350&50&100\\50&150&30\\100&30&600\end{smallmatrix}\right)$. 写出收益--风险模型并计算权重 $(0.5,0.2,0.3)$ 的两个指标.
\end{example}
\begin{solution}
权重满足 $\mathbf1\T x=1,x\ge0$, 收益为 $r\T x$, 方差为 $x\T Qx$. 给定最大方差 $v$, 可以最大化 $r\T x$ 并要求 $x\T Qx\le v$. 协方差矩阵半正定使这是凸问题, 虽非线性规划. 样本权重收益为11.87, 方差为191.1. 每一对不同资产的协方差在二次型中出现两次, 不能只计算对角项. 改变风险上界形成有效前沿; 这是历史数学模型, 不涉及对实际证券作投资推荐.
\end{solution}

\originalsection{排班与辅助变量}
\begin{example}
第2讲与教程1的邮局例要求七天人数至少为 $(17,13,15,19,14,16,11)$, 每人连续工作五天再休息两天. 建立最少雇员模型并比较连续与整数方案.
\end{example}
\begin{solution}
令 $x_j$ 表示在第 $j$ 天开始五天班次的人数, 下标按模7循环. 第 $i$ 天上班人数为 $\sum_{k=0}^4x_{i-k}$, 因而模型是 $\min\sum_jx_j$, 约束 $\sum_{k=0}^4x_{i-k}\ge b_i$, $x\ge0$. 单纯定义每天上班总人数不能保证同一个人连续工作五天, 这就是选择班次变量的作用.

连续方案 $(4/3,10/3,2,22/3,0,10/3,5)$ 满足所有约束, 总数 $67/3$. 取非负乘子 $y=(1/3,0,1/3,1/3,0,1/3,0)$. 每个班次的覆盖乘子和都不超过1, 而需求加权和为 $67/3$. 因而对所有可行 $x$, $\sum_jx_j\ge y\T Ax\ge y\T b=67/3$, 候选方案达到下界, 证明连续最优.

整数可行方案 $(4,4,2,6,0,4,3)$ 总数23. 由连续最优值 $67/3$, 整数最优值至少为 $\lceil67/3\rceil=23$, 该方案遂最优. 原教程第23页写的向量把第六项写为6, 总和实际为25; 其后23周轮换表却给第六类4周. 本稿采用可行的4, 并给出上述连续下界证书以证明整数方案最优.

公平轮换可用长度23的班次序列, 各开始日重复次数依次为4、4、2、6、0、4、3. 让23名员工取该序列的23个循环平移, 每周各类人数保持不变, 每人在23周周期内经历相同班次多重集. 这证明了轮换的公平计数, 不自动处理跨周劳动法规或其他现实排班约束.
\end{solution}

若允许不足与超员, 可令每天 $w_i+d_i-e_i=b_i$, $d_i,e_i\ge0$, 并在目标加入正系数的不足与超员惩罚, 如 $5d_i^2+e_i^2$. 可行集中会出现 $d_i,e_i$ 同时正的“多余表示”, 但任何这样的解都可把两者同时减去 $\min(d_i,e_i)$, 保持等式而严格降低目标. 所以每个最优解至少一者为0, 正好具有原意. 若惩罚系数不正, 或目标不随这些变量严格增加, 这个等价性证明可能失效.

\originalsection{线性变换与标准形}
本章标准形统一为 $\max c\T x$, $Ax=b$, $x\ge0$, 并预处理使 $b\ge0$ 及 $A$ 行满秩. 极大与极小可通过目标换号相互转换. $a\T x\le b$ 加松弛变量 $s\ge0$ 变成 $a\T x+s=b$; $a\T x\ge b$ 减剩余变量 $s\ge0$ 变成 $a\T x-s=b$. 右端为负时先将整行乘以 $-1$, 同时反转不等号. $x_j\le0$ 可替换为 $-y_j$, $y_j\ge0$; 自由变量可拆成 $x_j=u_j-v_j$, $u_j,v_j\ge0$. 这一表示不唯一, 但投影回原变量时可行性与目标值完全相同.

\begin{example}
教程6的练习为 $\min(x_1-x_2+x_3)$, 约束 $x_1+2x_2-x_3\le3$, $-x_1+x_2+x_3\ge2$, $x_1-x_2=10$, $x_1\ge0,x_2\le0$, $x_3$ 自由. 转为极大标准形.
\end{example}
\begin{solution}
令 $y_2=-x_2$, $x_3=y_3-y_4$, 再加松弛及剩余变量, 得
\begin{align*}
\max\quad&-x_1-y_2-y_3+y_4,\\
 x_1-2y_2-y_3+y_4+s_1&=3,\\
 -x_1-y_2+y_3-y_4-s_2&=2,\\
 x_1+y_2&=10,
\end{align*}
所有新变量非负. 由等式消去自由变量也是可行方法, 但消去有符号限制的变量时必须保留对应不等式; 否则会把不可行方案引入. 目标常数即使不影响最优变量, 也要在恢复原目标值时计入.
\end{solution}

教程4还介绍三类常用变换. $|a\T x+b|\le d$ 等价于两条线性不等式, 要求 $d\ge0$; $|a\T x+b|\ge d>0$ 通常是两个半空间的并, 非凸, 不能用仅含连续变量的线性规划提升来一般表示, 因为线性像与投影保持凸性. 最小化有限个仿射函数的最大值可引入 $t\ge a_i\T x+b_i$ 并最小化 $t$; 最大化有限个仿射函数的最小值则令 $t\le a_i\T x+b_i$ 并最大化 $t$. 若反转优化方向, 新变量可能不被压到正确边界, 不能机械套用.

比例约束 $(a\T x+b)/(d\T x+e)\ge\rho$ 只有在分母已知严格正时才等价于 $a\T x+b\ge\rho(d\T x+e)$. 零分母处原比例未定义, 交叉相乘后的线性式却可能允许该点; 若采用这种扩展必须明确约定, 不能称严格等价. 教程中的“至少20\%广告属于某类”因此还要有正总广告量假设或明示零投放时的比例约定.

\originalsection{基本可行解与顶点}
设 $A\in\R^{m\times n}$ 行满秩. 选择 $m$ 个线性无关列组成可逆矩阵 $B$, 对应变量为基本变量, 其他为非基本变量. 设置 $x_N=0$ 后, $x_B=B^{-1}b$ 称基本解; 若它非负, 称基本可行解. 此时字典为
\[
x_B=B^{-1}b-B^{-1}A_Nx_N,\qquad
z=c_B\T B^{-1}b+\bar c_N\T x_N,
\]
其中 $\bar c_j=c_j-c_B\T B^{-1}a_j$ 是检验数或约化成本. 原课程的 $-z$ 行与这里的目标字典是同一种表达, 只需按等式解回 $z$, 不可直接凭行首符号猜方向.

\begin{theorem}
标准形可行点为顶点当且仅当其正坐标对应的列线性无关. 每个非空标准形可行集有基本可行解; 若最优值有限, 存在一个最优基本可行解.
\end{theorem}
\begin{proof}
若正坐标的列相关, 存在非零 $d$ 支持于这些坐标且 $Ad=0$. 足够小的正负步长保持非负, 点是两个不同可行点的中点, 不是顶点. 若列无关, 将 $x$ 写成可行点 $y,z$ 的中点, 非负性强制所有零坐标在 $y,z$ 也为0, 列无关又使 $A(y-z)=0$ 只能 $y=z$, 所以是顶点. 正坐标的独立列可补足为一个基, 因而顶点就是某个基本可行解.

任一可行点若支持列相关, 选择上述方向, 沿一个方向移动直到一个正坐标变为0, 得支持更小的可行点. 有限次后列独立. 若从最优点出发, 足够小的正负移动都可行, 最优性强制 $c\T d=0$, 因而沿方向降支持时保持最优. 有限最优值的取到已由多面体上的线性函数存在性证明, 所以得到最优基本可行解.
\end{proof}
一般多面体若含直线, 可以没有顶点. 若不含直线, 通过有限面下降或标准形的支持消去可得顶点; 在有有限最优值时, 最优面也不含直线, 因而有最优顶点. “顶点处必有最优解”需要非空、有限最优值及顶点存在, 不能用于不可行或无界问题. 开约束 $0<x<1$ 的 $\sup x=1$ 未取到, 也不属于闭多面体理论.

\originalsection{换基、检验数与终止证书}
当前基本可行解满足 $\bar c_j\le0$ 对全部非基本变量时, 任何可行解都有 $z=z_0+\sum_j\bar c_jx_j\le z_0$, 因而当前解最优. 这也给对偶证书 $\pi=B^{-\mathsf T}c_B$, 因为 $A\T\pi\ge c$, $b\T\pi=z_0$. 有正检验数时, 选择进入变量 $x_s$ 并增加至 $t\ge0$. 基本变量变为 $x_B(t)=\beta-t d$, $\beta=B^{-1}b$, $d=B^{-1}a_s$.

若所有 $d_i\le0$, 任意 $t\ge0$ 都可行而目标 $z_0+t\bar c_s\to\infty$, 给无界证书. 否则最大可行步长为 $t^*=\min_{d_i>0}\beta_i/d_i$, 达到该最小值的基本变量离基. 换基等价于以其行、进入列元素为主元作 Gauss--Jordan 消元: 主元行除以主元, 再消去其他行及目标行的同列系数. 只在正的 $d_i$ 上作比值检验, 负比值不能当作更小可行步长.

\begin{example}
第4讲用方程组 $2x_1+2x_2+x_3=9$, $2x_1-x_2+2x_3=6$, $x_1-x_2+2x_3=5$ 说明消元. 第5讲的字典为
\begin{align*}
z&=11-2x_2+6x_5,\\
x_3&=4-2x_2-2x_5,\quad x_4=1+x_2+2x_5,\\
x_1&=9-6x_2-3x_5.
\end{align*}
分别完成线性消元与单纯形的一次换基.
\end{example}
\begin{solution}
方程组先以首行的一半消去第一列, 得 $x_1+x_2+x_3/2=9/2$, $-3x_2+x_3=-3$, $-2x_2+3x_3/2=1/2$. 第二行规范化后消去第二列, 得 $x_1+5x_3/6=7/2$, $x_2-x_3/3=1$, $5x_3/6=5/2$. 最后 $x_3=3,x_2=2,x_1=1$.

字典中令 $x_5$ 进入. 正限制来自 $x_3$ 与 $x_1$, 比值为 $4/2=2$ 与 $9/3=3$, 所以 $x_3$ 离基. 解出 $x_5=2-x_2-x_3/2$, 代回得 $x_4=5-x_2-x_3$, $x_1=3-3x_2+3x_3/2$, $z=23-8x_2-3x_3$. 检验数都非正, 所以 $x_1=3,x_4=5,x_5=2$、其余0为最优, 值23. 若原字典改为 $z=6-2x_2+x_5$, $x_3=4-2x_2+2x_5$, $x_4=2+x_2+2x_5$, $x_1=3-6x_2$, 则令 $x_2=0,x_5\to\infty$ 给可行无界射线, 正好对应第5讲的无界例.
\end{solution}

\originalsection{Phase I与人工变量}
松弛变量表示原约束的剩余资源, 在最终解中可以为正. 人工变量只是为了构造初始基, 若它在辅助最优解中仍为正, 就不能当作原问题变量. 对 $Ax=b$, $x\ge0$, 先使 $b\ge0$, 加人工变量 $y\ge0$, 解
$\min\mathbf1\T y$, $Ax+y=b$, $x,y\ge0$. 初始基为 $y=b$. 目标非负, 且原问题可行当且仅当辅助最优值为0. 若原问题可行, 把其解与 $y=0$ 代入得0; 若辅助取到0, 非负性强制所有 $y=0$, 恢复原可行解.

若辅助值0但人工变量仍在基中且值为0, 可在该行选择一个非人工的非零系数换出它; 若没有这种系数, 该行是冗余等式, 可以在保留等价可行集的前提下移除. 然后删除人工列, 恢复原目标, 用基本行消去目标中的基本变量, 开始 Phase II. 原讲义把人工变量比作暂时的松弛, 但两者的语义与最终允许状态完全不同.

\begin{example}
第5讲的初始等式为 $x_1+x_2+x_3=4$, $-2x_1+x_2-x_3=1$, $x\ge0$. 构造并求解 Phase I.
\end{example}
\begin{solution}
加 $y_1,y_2$ 后初始值为 $(4,1)$, 辅助目标5. 首先让 $x_2$ 进入, 增至1使 $y_2=0$, 此时 $y_1=3$. 解出 $x_2=1+2x_1+x_3-y_2$, 代入另一行, 得 $y_1=3-3x_1-2x_3+y_2$, 目标为 $3-3x_1-2x_3+2y_2$. 令 $x_1$ 增至1, $y_1$ 变为0, $x_2=3,x_3=0$, 辅助值0. 所以原问题可行, 可使用 $x_1,x_2$ 作为起始基, 再按原目标继续.
\end{solution}

\originalsection{退化、循环与 Bland 规则}
基本可行解的某个基本变量为0时称退化. 最小比值可以为0, 此时换基但变量点和目标值不变. 同一个点可以对应多个基, 因而“每次换基就到了新点”不成立. 没有退化时每次正检验数换基都严格提高目标, 基只有有限个, 因而算法有限. 退化时需防循环.

\begin{example}
教程7第10--12页的循环例, 初始基为 $x_5,x_6,x_7$, 目标字典及三行约束为
\[
z=3+\tfrac34x_1-20x_2+\tfrac12x_3-6x_4,
\quad
\begin{pmatrix}
1/4&-8&-1&9&1&0&0\\
1/2&-12&-1/2&3&0&1&0\\
0&0&1&0&0&0&1
\end{pmatrix}x=\begin{pmatrix}0\\0\\1\end{pmatrix}.
\]
取进入序列 $x_1,x_2,x_3,x_4,x_5,x_6$, 离基序列 $x_5,x_6,x_1,x_2,x_3,x_4$, 验证循环.
\end{example}
\begin{solution}
每次按对应行的正主元消元, 所有比值都取0. 基序列依次为
\[
\begin{aligned}
&(5,6,7)\to(1,6,7)\to(1,2,7)\to(3,2,7),\\
&(3,2,7)\to(3,4,7)\to(5,4,7)\to(5,6,7).
\end{aligned}
\]
 第六次后矩阵与目标行恢复初始值, 变量点始终为 $x_7=1$、其余0, 目标始终3. 这说明合法的进入选择加普通比值规则可以无限循环, 并未说明所有选择规则都会循环. 来源课件对该例的历史作者归属未在本稿另行确认, 这里只使用已核的数学数据.
\end{solution}

\begin{theorem}[Bland 防循环规则]
极大单纯形中, 选择具有正检验数的最小编号非基本变量进入; 最小比值并列时选择编号最小的基本变量离开. 此规则不会循环, 因而在有限次换基后给最优或无界证书.
\end{theorem}
\begin{proof}
若有循环, 目标沿整个循环非降又回原值, 所以每次都退化, 可行变量点不变. 称循环中曾改变基本身份的变量为变化变量, 取编号最大的变化变量 $x_t$. 选一次 $x_t$ 离基、$x_s$ 进入的换基; 有 $s<t$. 该次进入方向 $d$ 满足 $d_s=1$, $d_i=-a_{is}$ 对当时基本变量, 其他非基本分量0, 且 $d_t<0$. 所有变化变量在固定可行点的值都为0. 若某个当时基本的变化变量 $i<t$ 有 $a_{is}>0$, 它也具有零最小比值, 按离基规则应先于 $t$ 离开, 矛盾. 因而对全部变化变量 $i<t$, 都有 $d_i\ge0$.

另选循环中 $x_t$ 进入的一个字典. 按进入规则, 该字典中编号小于 $t$ 的非基本变量检验数都非正, 而 $\bar c_t>0$. 编号大于 $t$ 的非基本变量不是变化变量, 始终非基本, 所以它们在前一个方向 $d$ 中为0. 同一满足 $Ad=0$ 的代数方向的目标变化可在任意字典中计算. 退化时它不一定允许正的可行步长, 此处只使用零空间恒等式. 在前一个字典它等于进入变量的正检验数 $\bar c_s>0$; 在后一个字典却为 $\sum_{j\in N}\bar c_jd_j<0$, 因为 $t$ 项严格负, 较小编号项非正, 较大编号项为0. 矛盾. 基有限, 无循环就保证有限性.
\end{proof}
教程7第15页把并列离基写成“编号最小行, 不考虑变量编号”; 这不是上述标准 Bland 规则. 若行随换基不按当前基本变量编号重新排序, 不能用行位置替代变量编号. 本稿采用可证明防循环的变量编号版本.

退化与多最优解也不同. 非基本变量有零检验数且允许正步长时, 可以沿该方向得到不同的同值可行解; 若只有零步长, 只能得到同一点的另一基. 因而“存在零检验数”本身不能无条件证明有不同的最优变量点.

\originalsection{基不变范围与影子价格}
固定最优基 $B$, 右端扰动 $b+\Delta b$ 时, 基本解为 $B^{-1}b+B^{-1}\Delta b$. 检验数不变, 因而只要这一向量仍非负, 原基仍最优, 最优值精确变化为 $\pi\T\Delta b$, 其中 $\pi=B^{-\mathsf T}c_B$. 对单个右端分量, 每个基本变量的非负条件给一个线性区间, 相交即有效范围. 超出范围后应换基, 旧影子价格可能仅为一个对偶界.

成本扰动 $c+\Delta c$ 不改变基本可行解, 但改变检验数. 新基仍最优当且仅当 $c_N+\Delta c_N-(c_B+\Delta c_B)\T B^{-1}A_N\le0$. 非基本变量单独改变利润时, 它的允许增加正好是负检验数的绝对值; 基本变量的改变会同时影响多个检验数, 需逐一求区间. 软件报告中的极大数如 $10^{30}$ 通常表示无上限, 不能把它当作可测的有限容许量.

\begin{example}
教程5的二维模型为 $\max(30K+50S)$, $2K+3S\le10$, $K+2S\le6$, $K+S\le5$, $K\le4$, $S\le3$, $K,S\ge0$. 求基点、第一条约束的影子价格与 $S$ 利润的允许变动范围.
\end{example}
\begin{solution}
前两条取等号得 $(K,S)=(2,2)$, 值160. 乘子 $(10,10)$ 对前两约束给目标法向量 $(30,50)$, 其他置零, 所以由对偶证书知最优. 将第一右端改为 $10+\Delta$, 同一交点为 $K=2+2\Delta$, $S=2-\Delta$, 值 $160+10\Delta$. 所有其余不等式及非负性给 $-1\le\Delta\le1$, 故影子价格10在该范围精确有效.

若目标变为 $30K+(50+\eta)S$, 保持该点最优要求 $(30,50+\eta)$ 是活跃法向量 $(2,3),(1,2)$ 的非负组合. 解乘子得第一为 $10-\eta$, 第二为 $10+2\eta$, 因而 $-5\le\eta\le10$. 两端某个乘子为0, 目标与边平行, 可出现多最优解. 区间内部基点不变而目标值按 $2\eta$ 改变.
\end{solution}

若多个右端同时改变, 单项允许范围不能逐项独立套用. 设方向对应的单项最大幅度为 $r_i>0$, 若 $\sum_i|\Delta b_i|/r_i\le1$, 则扰动向量是零扰动与各轴允许端点的凸组合. 同一基可行右端区域由线性不等式定义而凸, 所以新右端仍可行, 检验数不变, 影子价格仍有效. 这是100\%规则的证明; 它是充分条件, 不是必要条件. 无穷允许幅度可通过有限近似端点取极限处理.

\begin{example}
第6讲虚构的 mc2 生产例最优变量为 $(18,16.4,3,0,32)$, 利润46.12百万美元. 其新磁盘约束影子价格0.8、右端20、允许减少16.4及增加0.6; 产品A需求上限的影子价格1.04、右端18、允许减少18及增加0.75; 产品C上限影子价格0.6、右端3、允许减少3及增加0.6. 应用这些报告数字.
\end{example}
\begin{solution}
单位产量为千件, 利润单位为百万美元. 新磁盘从20降至15仍在 $[3.6,20.6]$ 内, 利润变化为 $-5\times0.8=-4$ 百万美元. 将A需求上限从18增至18.5, 变化 $0.5\times1.04=0.52$ 百万美元; 若广告成本0.4百万美元, 在其余假设不变时净增加0.12百万美元. 这只是模型内比较, 不预设现实需求一定按假设改变.

同时把A上限降至15、C上限降至2, 100\%比率为 $3/18+1/3=1/2$, 所以原影子价格有效, 利润降低 $3\times1.04+1\times0.6=3.72$ 百万美元. 若A降1而C增0.5, 比率 $1/18+0.5/0.6<1$, 利润变化 $-1.04+0.3=-0.74$ 百万美元. 非活跃资源约束的影子价格为0, 小范围增加该资源并不提高利润; 活跃约束也可能乘子为0, 不能把“活跃”与“有价值”当作等价.
\end{solution}

约化成本还有经济解释 $\bar c_j=c_j-\pi\T a_j$: 收入减去按影子价格计的资源成本. mc2 的产品D收益0.6, 使用0.5单位新磁盘及1单位电子阅读器容量, 对应资源价0.8及0.3, 所以 $\bar c_D=0.6-0.5\times0.8-0.3=-0.1$. 收益升至0.7才消除负检验数; 等号处未必立即要求正生产, 还要看允许进入方向与退化. 新产品F收益0.65、使用0.4单位新磁盘与1单位阅读器容量, 约化成本为 $0.65-0.4\times0.8-0.3=0.03>0$, 表示当前对偶价格不再能证明原方案最优, 可以通过换基寻求改进. 若模型同时改变多个成本或矩阵中的基列, 必须重新检查全部基可行性与检验数, 单一报告范围不足以说明结论.

\originalchapter{分析与计算的联系}\label{chap:13}
本章对应原第25讲复习与收束, 第317--339页. 该讲把全课程重新归为凸分析、对偶、算法三部分. 以下以可核的解证书串联它们, 不重复抄写前文已经证明的结果.

\originalsection{从模型到证书}
一个凸问题能否求解, 至少有四个不同问题. 首先, 是否存在可行点并且最优值有限; 其次, 最优值是否取到; 然后, 是否有零对偶间隙与最优乘子; 最后, 计算过程是否以可证明的速度产生足够好的近似. 凸性只使局部最优等于全局最优, 并不自动回答后三个问题.

第\ref{chap:1}章把凸性变成全局切平面下界与投影条件. 第\ref{chap:2}章用相对内部控制低维域, 用衰退锥判定紧水平集与解的存在. 第\ref{chap:3}章说明线性像和部分极小化可能丢失闭性, 并以分离与共轭把点集描述转换成半空间描述. 第\ref{chap:4}章的扰动函数将零对偶间隙解释为原点的高度不在闭凸化中下降, 将乘子存在解释为原点有非竖直支持平面. 第\ref{chap:5}章把这个支持斜率写成次梯度, 将其用于最优性与灵敏度.

对可行候选 $x$ 与对偶可行乘子, 原目标是上界, 对偶值是下界. 二者之差才是直接的停止证书. 若有不可行迭代, 目标值可能低于原最优值, 不能独立当成上界. 次梯度距离估计、近端势函数和障碍的原始--对偶间隙, 分别提供不同的可证明进展量. “目标看起来下降”与“已达到指定精度”不是同一个结论.

\originalsection{算法利用的结构}
当目标不可微而次梯度廉价, 第\ref{chap:7}章以距离递推分析投影次梯度. 它允许目标上升, 步长选择影响值误差与点收敛, 所以逐步下降并非必要的进展指标. 当支持平面容易保存, 第\ref{chap:8}章建立外部多面体下模型; 当线性优化子问题容易, 单纯形分解维护内部可行凸包. 两者在共轭下相互对应, 但它们给出的界方向和子问题条件不同.

第\ref{chap:9}章通过二次正则化产生唯一的近端点, 并用到最优集合的距离解释其收敛. 在对偶侧应用同一近端步骤, 就出现增广拉格朗日乘子更新; 加上多面体近似则产生束方法. 第\ref{chap:10}章的障碍法从严格可行域内部逼近边界, 其参数趋零可能带来病态, 需要适合的内层求解. 大和目标则可利用循环或随机增量, 此时单次工作量减少, 但固定步长可能只到误差邻域.

对光滑目标, 第\ref{chap:11}章的二次上界使梯度投影由次梯度的 $O(\varepsilon^{-2})$ 次数改进到 $O(\varepsilon^{-1})$; 外推与势函数进一步给 $O(\varepsilon^{-1/2})$ 次数阶. 这些次数界以相应梯度、投影或近端查询为单位, 并不忽略每次查询的成本. 非光滑平滑化需要同时计算近似误差与平滑函数的Lipschitz常数. Bregman正则化和熵镜像下降适应单纯形等几何结构, 但不能不核对定义域、正分量与三点不等式条件就直接搬用欧氏分析.

线性规划还具有有限基结构, 第\ref{chap:12}章的单纯形法沿基换元, 用约化成本与比值检验给最优或无界证书. 退化时算法可能换基而不移动点, 防循环规则在这里有实际数学作用. 灵敏度则是在基保持可行与最优的范围内, 把对偶乘子变为精确的参数斜率.

\originalsection{一个复合模型}
\begin{example}
对 $\lambda>0$, 考虑 $F(x)=\tfrac12\norm{Ax-b}^2+\lambda\norm{x}_1$. 给出最优性证书、一个可实现迭代及可证明的停止判断.
\end{example}
\begin{solution}
光滑部分梯度为 $A\T(Ax-b)$, Lipschitz常数可取 $L=\norm A^2$. 全空间有限连续与闭真 $\ell_1$ 项满足和法则, 所以最优性等价于存在 $s\in\partial\norm{x}_1$ 使 $A\T(Ax-b)+\lambda s=0$. 各坐标上, $x_j\ne0$ 时 $s_j=\operatorname{sgn}(x_j)$, $x_j=0$ 时 $s_j\in[-1,1]$. 因而非零坐标满足精确梯度等式, 零坐标满足 $|[A\T(Ax-b)]_j|\le\lambda$.

取 $0<\alpha\le1/L$ (若 $A=0$, 直接 $x=0$ 最优), 迭代 $z=x-\alpha A\T(Ax-b)$, 再逐坐标取 $x_j^+=\operatorname{sgn}(z_j)(|z_j|-\alpha\lambda)_+$. 这是梯度近端法, 软阈值由一维次梯度条件直接求得. 若使用外推, 应按第\ref{chap:11}章的参数与势函数证明分析, 不能只在该式前随意加一个动量项.

还可构造一个直接的对偶界. 将残差变量 $r=Ax-b$ 拆出并用乘子 $y$, 极小化 $\tfrac12\norm r^2+\lambda\norm x_1+y\T(Ax-b-r)$, 得对偶 $q(y)=-\tfrac12\norm y^2-b\T y$, 约束 $\norm{A\T y}_\infty\le\lambda$. 对任意残差候选 $r$, 若 $\norm{A\T r}_\infty>\lambda$, 将其按比例缩成对偶可行 $y$; 若已满足, 直接取 $y=r$. 因而 $F(x)-q(y)\ge0$ 是一个可计算误差上界. 仅检查软阈值后坐标变动小, 并不等于检查了这一目标误差.
\end{solution}
这个例子同时使用凸性、次微分、对偶、近端与次数分析. 它说明全书各部分的联系, 也说明模型与算法的每项假设都应在实际应用处重新核对.

\begin{thebibliography}{9}
\bibitem{mit} Dimitri P. Bertsekas. \emph{6.253 Convex Analysis and Optimization: Complete Lecture Notes}. MIT OpenCourseWare, Spring 2012, 340页. \url{https://ocw.mit.edu/courses/6-253-convex-analysis-and-optimization-spring-2012/pages/lecture-notes/}.
\bibitem{license} Massachusetts Institute of Technology. MIT OpenCourseWare Privacy and Terms of Use. \url{https://ocw.mit.edu/pages/privacy-and-terms-of-use/}. 访问日期: 2026-10-08.
\bibitem{orlin} James Orlin, Ebrahim Nasrabadi. \emph{15.053 Optimization Methods in Management Science}. MIT OpenCourseWare, Spring 2013. Lecture notes L2--6, L9; tutorials 1, 2, 4--7. \url{https://ocw.mit.edu/courses/15-053-optimization-methods-in-management-science-spring-2013/pages/lecture-notes/}; \url{https://ocw.mit.edu/courses/15-053-optimization-methods-in-management-science-spring-2013/pages/tutorials/}.
\bibitem{mono} Dimitri P. Bertsekas. \emph{Extended Monotropic Programming and Duality}. 2010 revised manuscript, definitions 4.1--4.2. \url{https://web.mit.edu/dimitrib/www/Extended_Mono.pdf}. 仅用于核对术语与相关高级工具, 不转载全文.
\end{thebibliography}

\end{document}
